% Editable text transcription generated from the supplied PDF.
% Compile with XeLaTeX. Layout is intentionally simplified for editing.
\documentclass[12pt,a4paper]{report}
\usepackage[a4paper,top=2.5cm,bottom=2.5cm,left=3cm,right=2.5cm]{geometry}
\usepackage{fontspec}
\usepackage{polyglossia}
\setdefaultlanguage{thai}
\setotherlanguage{english}
\IfFontExistsTF{Garuda}{\setmainfont{Garuda}}{\setmainfont{Thonburi}}
\IfFontExistsTF{Garuda}{\setsansfont{Garuda}}{\setsansfont{Thonburi}}
\IfFontExistsTF{Garuda}{\setmonofont{Garuda}}{\setmonofont{Thonburi}}
\IfFontExistsTF{Garuda}{\newfontfamily\thaifont[Script=Thai]{Garuda}}{\newfontfamily\thaifont[Script=Thai]{Thonburi}}
\IfFontExistsTF{Garuda}{\newfontfamily\thaifontsf[Script=Thai]{Garuda}}{\newfontfamily\thaifontsf[Script=Thai]{Thonburi}}
\IfFontExistsTF{Garuda}{\newfontfamily\thaifonttt[Script=Thai]{Garuda}}{\newfontfamily\thaifonttt[Script=Thai]{Thonburi}}
\IfFontExistsTF{Garuda}{\newfontfamily\englishfont{Garuda}}{\newfontfamily\englishfont{Thonburi}}
\usepackage{setspace}
\setstretch{1.25}
\usepackage{enumitem}
\usepackage{hyperref}
\usepackage{graphicx}
\usepackage{booktabs,longtable,tabularx,array}
\usepackage{amsmath,amssymb}
\setlength{\parindent}{1.5em}
\setlength{\parskip}{0.25em}
\begin{document}
\sloppy
\chapter*{คำชี้แจงสำหรับไฟล์ Editable}
ไฟล์นี้เป็นการถอดข้อความจาก PDF อัตโนมัติเพื่อให้แก้ไขเนื้อหาใน Overleaf ได้ง่ายขึ้น โดยรูปแบบหน้า ตาราง สมการ และตำแหน่งภาพอาจไม่ตรงกับต้นฉบับ 100\%. ฉบับนี้เพิ่มหัวข้อ 3.6--3.8 เกี่ยวกับนโยบายการตัดสินใจด้วย GA แล้ว ส่วนภาพและตารางเดิมที่ถอดข้อความไม่สมบูรณ์ให้ตรวจเทียบกับ PDF ต้นฉบับ ไฟล์ PDF ฉบับรวมในโฟลเดอร์นี้ใช้หน้าต้นฉบับร่วมกับหน้าส่วนเพิ่มเติมเพื่อคงรูปแบบเดิมไว้.
\clearpage
% ---- source PDF page 1 ----
\begin{center}\bfseries การพัฒนาแบบจำลองสำหรับการเชื่อมโยงข้อมูลจากแหล่งต่าง ๆ ของ\end{center}
\begin{center}\bfseries ลูกค้าเข้าด้วยกันเพื่อสร้างเป็นคุณลักษณะของลูกค้าคนเดียว และกำหนด\end{center}
\begin{center}\bfseries คะแนนแก่ลูกค้าเป้าหมายเพื่อใช้สำหรับระบบบริหารลูกค้าสัมพันธ์แบบ\end{center}
\begin{center}\bfseries อัจฉริยะ\end{center}
\begin{center}\bfseries DEVELOPMENT OF IDENTITY RESOLUTION AND LEAD SCORING\end{center}
\begin{center}\bfseries MODELS FOR AGENTIC CRM PLATFORM\end{center}
\begin{center}\bfseries โดย\end{center}
\begin{center}\bfseries วิมลสิริ ปานมา\end{center}
\begin{center}\bfseries ถาวรีย์ ขันสาคร\end{center}
\begin{center}\bfseries VIMOLSIRI PANMA\end{center}
\begin{center}\bfseries THAWAREE KHANSAKORN\end{center}
\begin{center}\bfseries ปริญญานิพนธ์นี้เป็นส่วนหนึ่งของการศึกษาตามหลักสูตรวิทยาศาสตรบัณฑิต\end{center}
\begin{center}\bfseries สาขาวิชาวิทยาการข้อมูลและการวิเคราะห์เชิงธุรกิจ\end{center}
\begin{center}\bfseries คณะเทคโนโลยีสารสนเทศ\end{center}
\begin{center}\bfseries สถาบันเทคโนโลยีพระจอมเกล้าเจ้าคุณทหารลาดกระบัง\end{center}
\begin{center}\bfseries ภาคเรียนที่ 2 ปีการศึกษา 2568\end{center}
% ---- source PDF page 2 ----
\clearpage
\begin{center}\bfseries การพัฒนาแบบจำลองสำหรับการเชื่อมโยงข้อมูลจากแหล่งต่าง ๆ ของ\end{center}
\begin{center}\bfseries ลูกค้าเข้าด้วยกันเพื่อสร้างเป็นคุณลักษณะของลูกค้าคนเดียว และกำหนด\end{center}
\begin{center}\bfseries คะแนนแก่ลูกค้าเป้าหมายเพื่อใช้สำหรับระบบบริหารลูกค้าสัมพันธ์แบบ\end{center}
\begin{center}\bfseries อัจฉริยะ\end{center}
\begin{center}\bfseries DEVELOPMENT OF IDENTITY RESOLUTION AND LEAD SCORING\end{center}
\begin{center}\bfseries MODELS FOR AGENTIC CRM PLATFORM\end{center}
\begin{center}\bfseries โดย\end{center}
\begin{center}\bfseries วิมลสิริ ปานมา\end{center}
\begin{center}\bfseries ถาวรีย์ ขันสาคร\end{center}
\begin{center}\bfseries อาจารย์ที่ปรึกษา\end{center}
\begin{center}\bfseries ดร.ปาณิตา ธูสรานนท์\end{center}
\begin{center}\bfseries อ.เฉลิมพล ศิริกายน\end{center}
\begin{center}\bfseries ปริญญานิพนธ์นี้เป็นส่วนหนึ่งของการศึกษาตามหลักสูตรวิทยาศาสตรบัณฑิต\end{center}
\begin{center}\bfseries สาขาวิชาวิทยาการข้อมูลและการวิเคราะห์เชิงธุรกิจ\end{center}
\begin{center}\bfseries คณะเทคโนโลยีสารสนเทศ\end{center}
\begin{center}\bfseries สถาบันเทคโนโลยีพระจอมเกล้าเจ้าคุณทหารลาดกระบัง\end{center}
\begin{center}\bfseries ภาคเรียนที่ 2 ปีการศึกษา 2568\end{center}
% ---- source PDF page 3 ----
\clearpage
\begin{center}\bfseries DEVELOPMENT OF IDENTITY RESOLUTION AND LEAD SCORING\end{center}
\begin{center}\bfseries MODELS FOR AGENTIC CRM PLATFORM\end{center}
\begin{center}\bfseries VIMOLSIRI PANMA\end{center}
\begin{center}\bfseries THAWAREE KHANSAKORN\end{center}
\begin{center}\bfseries A PROJECT SUBMITTED IN PARTIAL FULFILLMENT\end{center}
\begin{center}\bfseries OF THE REQUIREMENT FOR THE DEGREE OF\end{center}
\begin{center}\bfseries BACHELOR OF SCIENCE PROGRAM IN DATA SCIENCE AND BUSINESS\end{center}
\begin{center}\bfseries ANALYTICS\end{center}
\begin{center}\bfseries SCHOOL OF INFORMATION TECHNOLOGY\end{center}
\begin{center}\bfseries KING MONGKUT’S INSTITUTE OF TECHNOLOGY LADKRABANG\end{center}
\begin{center}\bfseries 2/2025\end{center}
% ---- source PDF page 4 ----
\clearpage
\begin{center}\bfseries COPYRIGHT 2025\end{center}
\begin{center}\bfseries SCHOOL OF INFORMATION TECHNOLOGY\end{center}
\begin{center}\bfseries KING MONGKUT’S INSTITUTE OF TECHNOLOGY LADKRABANG\end{center}
% ---- source PDF page 5 ----
\clearpage
\begin{center}\bfseries ใบรับรองปริญญานิพนธ์ ประจําปีการศึกษา 2568\end{center}
\begin{center}\bfseries คณะเทคโนโลยีสารสนเทศ\end{center}
\begin{center}\bfseries สถาบันเทคโนโลยีพระจอมเกล้าเจ้าคุณทหารลาดกระบัง\end{center}
\begin{center}\bfseries เรื่อง การพัฒนาแบบจำลองสำหรับการเชื่อมโยงข้อมูลจากแหล่งต่าง ๆ\end{center}
\begin{center}\bfseries ของลูกค้าเข้าด้วยกันเพื่อสร้างเป็นคุณลักษณะของลูกค้าคนเดียว\end{center}
\begin{center}\bfseries และกำหนดคะแนนแก่ลูกค้าเป้าหมายเพื่อใช้สำหรับระบบบริหาร\end{center}
\begin{center}\bfseries ลูกค้าสัมพันธ์แบบอัจฉริยะ\end{center}
\begin{center}\bfseries DEVELOPMENT OF IDENTITY RESOLUTION AND LEAD\end{center}
\begin{center}\bfseries SCORING MODELS FOR AGENTIC CRM PLATFORM\end{center}
\begin{center}\bfseries ผู้จัดทํา\end{center}
\noindent 1. นางสาววิมลสิริ ปานมา รหัสนักศึกษา 66070187\par
\noindent 2. นางสาวถาวรีย์ ขันสาคร รหัสนักศึกษา 66070260\par
\begin{center}\bfseries …….........................................อาจารย์ที่ปรึกษา\end{center}
\begin{center}\bfseries (.......................................)\end{center}
\begin{center}\bfseries ........................................อาจารย์ที่ปรึกษาร่วม\end{center}
\begin{center}\bfseries (.......................................)\end{center}
% ---- source PDF page 6 ----
\clearpage
\begin{center}\bfseries ใบรับรองโครงงาน (PROJECT)\end{center}
\begin{center}\bfseries เรื่อง\end{center}
\begin{center}\bfseries การพัฒนาแบบจำลองสำหรับการเชื่อมโยงข้อมูลจากแหล่งต่าง ๆ ของ\end{center}
\begin{center}\bfseries ลูกค้าเข้าด้วยกันเพื่อสร้างเป็นคุณลักษณะของลูกค้าคนเดียว และกำหนด\end{center}
\begin{center}\bfseries คะแนนแก่ลูกค้าเป้าหมายเพื่อใช้สำหรับระบบบริหารลูกค้าสัมพันธ์แบบ\end{center}
\begin{center}\bfseries อัจฉริยะ\end{center}
\begin{center}\bfseries DEVELOPMENT OF IDENTITY RESOLUTION AND LEAD SCORING\end{center}
\begin{center}\bfseries MODELS FOR AGENTIC CRM PLATFORM\end{center}
นางสาววิมลสิริ ปานมา รหัสนักศึกษา 66070187\par
นางสาวถาวรีย์ ขันสาคร รหัสนักศึกษา 66070260\par
ขอรับรองว่ารายงานฉบับนี้ ข้าพเจ้าไม่ได้คัดลอกมาจากที่ใด\par
รายงานฉบับนี้ได้รับการตรวจสอบและอนุมัติให้เป็นส่วนหนึ่งของ\par
การศึกษาวิชาโครงงาน หลักสูตรวิทยาศาสตรบัณฑิต (สาขาวิชาวิทยาการข้อมูลและการวิเคราะห์\par
เชิงธุรกิจ)\par
ภาคเรียนที่ 2 ปีการศึกษา 2568\par
.........................................\par
(นางสาววิมลสิริ ปานมา)\par
.........................................\par
(นางสาวถาวรีย์ ขันสาคร)\par
% ---- source PDF page 7 ----
\clearpage
\begin{center}\bfseries หัวข้อโครงงาน การพัฒนาแบบจำลองสำหรับการเชื่อมโยงข้อมูลจากแหล่งต่าง\end{center}
ๆ ของลูกค้าเข้าด้วยกันเพื่อสร้างเป็นคุณลักษณะของลูกค้าคน\par
เดียว และกำหนดคะแนนแก่ลูกค้าเป้าหมายเพื่อใช้สำหรับ\par
ระบบบริหารลูกค้าสัมพันธ์แบบอัจฉริยะ\par
\begin{center}\bfseries นักศึกษา นางสาววิมลสิริ ปานมา\end{center}
นางสาวถาวรีย์ ขันสาคร\par
\begin{center}\bfseries ปริญญา วิทยาศาสตรบัณฑิต\end{center}
\begin{center}\bfseries สาขาวิชา วิทยาการข้อมูลและการวิเคราะห์เชิงธุรกิจ\end{center}
\begin{center}\bfseries ปีการศึกษา 2568\end{center}
\begin{center}\bfseries อาจารย์ที่ปรึกษา ดร.ปาณิตา ธูสรานนท์ , อ.เฉลิมพล ศิริกายน\end{center}
\section*{บทคัดย่อ}
ในปัจจุบันองค์กรมักประสบปัญหาข้อมูลลูกค้ากระจัดกระจายอยู่ตาม\par
แพลตฟอร์มต่าง ๆ ทำให้ไม่สามารถมองเห็นภาพรวมของลูกค้าได้ชัดเจน ส่งผลให้การ\par
วิเคราะห์ Customer Journey ขาดความต่อเนื่องและทีมงานต้องเสียเวลาในการ\par
รวบรวมข้อมูลด้วยตนเอง โครงการนี้จึงมุ่งเน้นพัฒนาระบบส่วนหนึ่งของระบบ Agentic\par
CRM ที่ทำหน้าที่เป็นเสมือนศูนย์กลางอัจฉริยะ โดยการรวม Lead จากทุกแหล่งข้อมูล\par
ให้เป็นโปรไฟล์ลูกค้าเดียวกันโดยอัตโนมัติ และใช้ AI/ML ตรวจจับความสัมพันธ์ของ\par
ข้อมูล ผ่านกระบวนการ Identity Resolution เพื่อให้ได้ภาพรวมลูกค้าที่ชัดเจน ระบบ\par
นี้จะช่วยให้ทีมขายและการตลาดสามารถเข้าถึงข้อมูลเชิงลึกได้ทันที ซึ่งรวมถึงการ\par
วิเคราะห์ Customer Journey และแนะนำ Script เพื่อเพิ่มขีดความสามารถในการ\par
แข่งขันและปิดการขายได้อย่างมีประสิทธิภาพ\par
% ---- source PDF page 8 ----
\clearpage
\begin{center}\bfseries Project DEVELOPMENT OF IDENTITY RESOLUTION AND LEAD\end{center}
SCORING MODELS FOR AGENTIC CRM PLATFORM\par
\begin{center}\bfseries Student Miss Vimolsiri Panma Student ID 66070187\end{center}
Miss Thawaree Khansakorn Student ID 66070260\par
\begin{center}\bfseries Degree Bachelor of Science\end{center}
\begin{center}\bfseries Program Data Science and Business Analytics\end{center}
\begin{center}\bfseries Academic Year 2025\end{center}
\begin{center}\bfseries Advisor Dr. Panita Thusaranon, Mr. Chaloemphon Sirikayon\end{center}
\section*{ABSTRACT}
Nowadays, organizations often face the problem of customer data\par
being scattered across platforms, making it impossible to clearly see the\par
customer's overview. As a result, the Customer Journey analysis lacks\par
continuity and the team has to waste time collecting data manually. This\par
project focuses on developing a system part of the system. Agentic CRM\par
acts as a virtual hub of intelligence by automatically integrating leads from\par
all data sources into the same customer profile and using AI/ML to detect\par
the relationship of the data through the Identity Resolution process to\par
provide a clear customer overview. This system will allow sales and\par
marketing teams to access. Immediate insights, including Customer\par
Journey analysis and Script introduction to increase competitiveness and\par
effectively close sales.\par
% ---- source PDF page 9 ----
\clearpage
\section*{กิตติกรรมประกาศ}
การทําโครงงานเรื่อง การพัฒนาแบบจำลองสำหรับการเชื่อมโยงข้อมูลจาก\par
แหล่งต่าง ๆ ของลูกค้าเข้าด้วยกันเพื่อสร้างเป็นคุณลักษณะของลูกค้าคนเดียว และ\par
กำหนดคะแนนแก่ลูกค้าเป้าหมายเพื่อใช้สำหรับระบบบริหารลูกค้าสัมพันธ์แบบ\par
อัจฉริยะ และรายงานวิจัยฉบับนี้สามารถสําเร็จลุล่วงด้วยดี เนื่องจากได้รับความ\par
อนุเคราะห์และสนับสนุนเป็นอย่างดียิ่งจาก ดร.ปาณิตา ธูสรานนท์ และ อ.เฉลิมพล ศิริ กายน ที่ให้ความกรุณาอย่างสูงในการให้คําแนะนํา คําปรึกษา และให้ความรู้ในการทํา\par
วิจัย การตรวจทาน แก้ไขข้อบกพร่องต่าง ๆ จนการทําโครงงาน และการเขียนรายงาน\par
วิจัยจนสําเร็จลุล่วงด้วยดี\par
ขอขอบคุณเพื่อนร่วมงานที่อํานวยความสะดวก และช่วยเหลือในด้านต่าง ๆ ทั้ง\par
การศึกษา หาข้อมูล การลงมือทําโครงงาน และการเขียนรายงานวิจัยฉบับนี้จนสําเร็จ\par
ลุล่วงด้วยดี\par
สุดท้ายนี้คณะผู้จัดทําหวังว่าวิจัยฉบับนี้จะสามารถเป็นประโยชน์สําหรับ\par
หน่วยงานที่เกี่ยวข้องและผู้ที่สนใจศึกษาต่อไป\par
ชื่อ วิมลสิริ นามสกุล ปานมา\par
ชื่อ ถาวรีย์ นามสกุล ขันสาคร\par
% ---- source PDF page 10 ----
\clearpage
\section*{สารบัญ}
หน้า\par
บทคัดย่อ ........................................................................................................................................ I\par
ABSTRACT ................................................................................................................................... II\par
กิตติกรรมประกาศ ....................................................................................................................... III\par
สารบัญ ........................................................................................................................................ IV\par
สารบัญตาราง .............................................................................................................................. VI\par
สารบัญรูป ................................................................................................................................. VIII\par
บทที่ 1 ......................................................................................................................................... 1\par
\noindent 1.1 ที่มาและความสําคัญ ................................................................................................ 1\par
\noindent 1.2 วัตถุประสงค์ ............................................................................................................ 1\par
\noindent 1.3 วิธีการดําเนินงาน ..................................................................................................... 2\par
\noindent 1.4 ขอบเขตของงาน ...................................................................................................... 3\par
\noindent 1.5 ประโยชน์ที่คาดว่าจะได้รับ ....................................................................................... 3\par
\noindent 1.6 แผนการดำเนินงาน .................................................................................................. 4\par
บทที่ 2 เอกสารและงานที่เกี่ยวข้อง .............................................................................................. 5\par
\noindent 2.1 แนวคิดที่เกี่ยวข้อง .................................................................................................... 5\par
\noindent 2.2 ทฤษฎีที่เกี่ยวข้อง ..................................................................................................... 7\par
\noindent 2.3 สถาปัตยกรรมระบบที่เกี่ยวข้อง.............................................................................. 10\par
\noindent 2.4 งานวิจัยที่เกี่ยวข้อง ................................................................................................ 11\par
บทที่ 3 วิธีการดำเนินวิจัย .......................................................................................................... 15\par
\noindent 3.1 การทำความเข้าใจธุรกิจ (Business Understanding) .......................................... 16\par
\noindent 3.2 การทำความเข้าใจข้อมูล (Data Understanding) ................................................. 19\par
\noindent 3.3 การเตรียมข้อมูล (Data Preparation) .................................................................. 24\par
% ---- source PDF page 11 ----
\clearpage
\begin{center}\bfseries สารบัญ(ต่อ)\end{center}
หน้าที่\par
\noindent 3.4 การสร้างแบบจำลอง (Modeling) .......................................................................... 39\par
\noindent 3.5 การประเมินผล (Evaluation) ................................................................................ 49\par
บทที่ 4 ผลการทดลอง ................................................................................................................ 55\par
\noindent 4.1 การเตรียมข้อมูล .................................................................................................... 55\par
\noindent 4.2 วิธีวัดประสิทธิภาพ ................................................................................................. 58\par
\noindent 4.3 Candidate-Scoring Model ที่ใช้ก่อน Decision Thresholding ........................ 61\par
\noindent 4.4 Exact-First ร่วมกับ Calibrated Candidate Scoring สำหรับ Decision Thresholding ............................................................................................................. 70\par
\noindent 4.5 ผลลัพธ์ปลายทางของ MATCH, REVIEW และ NO\_MATCH สำหรับ CRM .......... 71\par
\noindent 4.6 ความสำคัญของคุณลักษณะทั้ง 41 ตัวที่สนับสนุนการตัดสินใจ .............................. 76\par
\noindent 4.7 สรุปผลของ Project 1 ในขอบเขต Decision Thresholding ............................... 78\par
บทที่ 5 วิเคราะห์และสรุปผล ..................................................................................................... 79\par
\noindent 5.1 สรุปผลการดำเนินงาน ............................................................................................ 79\par
\noindent 5.2 ปัญหาและอุปสรรคในการดำเนินงาน..................................................................... 84\par
\noindent 5.3 ข้อจำกัดของงานวิจัย ............................................................................................. 85\par
\noindent 5.4 ข้อเสนอแนะสำหรับการทำวิจัยในอนาคต .............................................................. 86\par
บรรณานุกรม.............................................................................................................................. 88\par
% ---- source PDF page 12 ----
\clearpage
\section*{สารบัญตาราง}
ตารางที่ หน้า ตารางที่ 1.1 แผนการดำเนินงาน ................................................................................................ 4 ตารางที่ 2.1 เปรียบเทียบ Similarity Measures ที่ใช้ในงาน Identity Resolution ................. 8 ตารางที่ 2.2 สรุปงานวิจัยที่เกี่ยวข้องกับโครงงาน ..................................................................... 13 ตารางที่ 3.1 จำนวนข้อมูลแยกตามแพลตฟอร์มในชุดข้อมูลมาตรฐานหลัก ................................ 20 ตารางที่ 3.2 ฟิลด์อ้างอิงหลักและบทบาทในการใช้งานจริงของระบบ ........................................ 20 ตารางที่ 3.3 ตัวอย่างการแปลงข้อมูลหลักในขั้นทำความสะอาดและปรับมาตรฐาน ................... 25 ตารางที่ 3.4 ตัวอย่างการสร้างฟิลด์สนับสนุนจากข้อมูลบริบท ................................................... 27 ตารางที่ 3.5 โครงสร้างชุดคู่ข้อมูลสำหรับการฝึกแบบจำลอง ...................................................... 29 ตารางที่ 3.6 ขนาดของชุดข้อมูลที่ได้หลังแบ่ง train, validation และ test ............................... 31 ตารางที่ 3.7 ขนาดของ merged feature matrices ใน stage 9 ............................................. 38 ตารางที่ 3.8 สายการทดลองหลักในขั้น Modeling .................................................................... 40 ตารางที่ 3.9 ชุดการทดลองใน multimodal suite ................................................................... 41 ตารางที่ 3.10 สรุปเกณฑ์ที่ใช้คัดเลือกแบบจำลองในขั้น Modeling ........................................... 41 ตารางที่ 3.11 ระดับการทดลองในสาย Multimodal Suite ...................................................... 45 ตารางที่ 3.12 สรุปค่าตั้งต้นไฮเปอร์พารามิเตอร์ของแบบจำลองหลักในขั้นตอน Modeling ....... 46 ตารางที่ 3.13 สรุปผลเชิงตัวเลขและการคัดเลือกแบบจำลองในแต่ละสายการทดลอง................ 48 ตารางที่ 3.14 สรุประดับการประเมินผลและตัวชี้วัดหลักที่ใช้ในงานวิจัย .................................... 51 ตารางที่ 3.15 กรอบการประเมินขั้น Retrieval และ Candidate Generation ........................ 52 ตารางที่ 3.16 กรอบการประเมินแบบจำลองในสายการทดลองหลัก .......................................... 53 ตารางที่ 3.17 กรอบการประเมินผลระดับ Calibration และ Production Workflow ............. 54 ตารางที่ 4.1 จำนวน valid profiles แยกตามแพลตฟอร์มที่ใช้จริง ............................................ 56 ตารางที่ 4.2 ผลของขั้นเตรียม candidate set ก่อนเข้าสู่ candidate scoring ........................ 57 ตารางที่ 4.3 โครงสร้างของ supervised split สำหรับ main run ............................................ 57 ตารางที่ 4.4 ผล validation leaderboard ภายใน main run ................................................. 63 ตารางที่ 4.5 ผลเปรียบเทียบ candidate-scoring configurations ที่ใช้เลือก main run ......... 64 ตารางที่ 4.6 Confusion matrix เปรียบเทียบ candidate-scoring configurations บน test split ............................................................................................................................ 65 ตารางที่ 4.7 ผลของ main run บน test split ก่อนเข้าสู่ production thresholding ............ 66 ตารางที่ 4.8 Retrieval funnel ของ Project 1 ก่อนเข้าสู่การตัดสินใจปลายทาง ..................... 71 ตารางที่ 4.9 Operating point ที่ใช้จริงในขั้น decision thresholding ................................... 72\par
% ---- source PDF page 13 ----
\clearpage
\section*{สารบัญตาราง}
ตารางที่ หน้า ตารางที่ 4.10 คุณภาพของแต่ละ tier ที่เกี่ยวข้องกับการตัดสินใจปลายทาง .............................. 73 ตารางที่ 4.11 สรุป pattern ของข้อผิดพลาดจากตัวอย่าง `fp\_top.csv` และ `fn\_top.csv` ... 75 ตารางที่ 4.12 คุณลักษณะที่สำคัญสูงสุดของ main run ............................................................ 77 ตารางที่ 5.1 สรุปเปรียบเทียบประสิทธิภาพจากทุกสายการทดลอง ........................................... 81 ตารางที่ 5.2 สรุปผลระดับ Production .................................................................................... 82\par
% ---- source PDF page 14 ----
\clearpage
\section*{สารบัญรูป}
รูปที่ หน้า รูปที่ 2.1 สถาปัตยกรรมระบบ LINKSOCIAL สำหรับการเชื่อมโยงโปรไฟล์ข้ามแพลตฟอร์ม ..... 10 รูปที่ 2.2 สถาปัตยกรรม Human-in-the-Loop Data Integration ........................................ 11 รูปที่ 2.3 ขั้นตอนการทำงานของระบบ Identity Resolution และ Lead Scoring ................. 14 รูปที่ 3.1 แสดงกรอบวิธีวิจัยที่ใช้เป็นโครงหลักของการพัฒนาระบบเชื่อมโยงอัตลักษณ์ลูกค้าและ การให้คะแนนลูกค้าเป้าหมายในงานวิจัยนี้ ................................................................................. 15 รูปที่ 3.2 ซึ่งแสดงภาพรวมของระบบในเชิงธุรกิจและเชิงเทคนิคพร้อมกัน โดยให้เห็นเส้นทาง ข้อมูลจากโปรไฟล์ต้นทาง การคัด candidate pairs การให้คะแนน การตัดสินใจ การรวมเอนทิตี และการสร้างคะแนนลูกค้า เพื่อให้ผู้อ่านเห็นตั้งแต่ต้นว่างานนี้ไม่ได้จบที่การทำนายคู่ข้อมูล แต่ไป ถึงการใช้งานจริงในระบบปลายทาง............................................................................................ 17 รูปที่ 3.3 จำนวนข้อมูลที่ไม่เป็นค่าว่างของฟิลด์สำคัญในชุดข้อมูล LinkSocial .......................... 22 รูปที่ 3.4 แสดงการไหลของข้อมูลจากแหล่งข้อมูลดิบไปสู่ชุดข้อมูลมาตรฐาน ............................. 23 รูปที่ 3.5 เพื่อแสดงความสัมพันธ์เชิงคีย์ของข้อมูลในลักษณะ Entity-Relationship Diagram .. 23 รูปที่ 3.6 Pair building pipeline การสร้างชุดคู่ข้อมูลสำหรับการฝึกแบบจำลอง ..................... 30 รูปที่ 3.7 ลำดับการลดทอนคอลัมน์และคัดเลือกคุณลักษณะ ...................................................... 49 รูปที่ 3.8 ลำดับขั้น Model Development Pipeline .............................................................. 49 รูปที่ 3.9 ความสัมพันธ์ของสายการทดลองและจุดบรรจบ ......................................................... 55 รูปที่ 3.8 Evaluation Framework: กรอบการประเมินผลของระบบในหลายระดับ .................. 55 รูปที่ 4.1 สรุปผลการเตรียมข้อมูล ................................................................................................. 58 รูปที่ 4.2 ตัวอย่างโปรไฟล์จาก 3 แพลตฟอร์มหลังผ่าน Standardization ............................... 58 รูปที่ 4.3 โครงสร้างของ main-run model pipeline สำหรับ Project 1 ....................................... 62 รูปที่ 4.4 การเปรียบเทียบ model families บน final 41-feature split ........................................ 63 รูปที่ 4.5 การเปรียบเทียบ metric ของ candidate-scoring configurations ......................... 64 รูปที่ 4.6 Confusion matrices ของ candidate-scoring configurations ทั้ง3แบบบน test split . 65 รูปที่ 4.7 Confusion matrix ของ main run บน test split ที่ threshold = 0.35 ........................ 66 รูปที่ 4.8 การกระจายของ calibrated scores บน test split แยกตาม Actual MATCH และ Actual NO\_MATCH พร้อมเส้นอ้างอิงที่ threshold = 0.35 และเส้น production thresholds ที่ 0.95 และ 0.98 ...................................................................................................................... 67 รูปที่ 4.9 พฤติกรรมของ Precision Recall และ F1 ของ main run เมื่อเปลี่ยน threshold บน test split แสดงความแตกต่างระหว่าง threshold เชิงจำแนกที่ 0.35 กับ production review threshold ที่ 0.95 .................................................................................................................... 68\par
% ---- source PDF page 15 ----
\clearpage
\section*{สารบัญรูป}
รูปที่ หน้าที่ รูปที่ 4.10 Reliability ของ calibrated probabilities หลังผ่าน Isotonic Regression เปรียบเทียบ validation set และ test set โดยใช้ขนาด marker แทนจำนวนคู่ในแต่ละ probability bin และเส้นตั้งแสดง production thresholds ที่ 0.95 และ 0.98 ..................... 69 รูปที่ 4.11 คุณภาพการจัดอันดับของ candidate scoring บน full candidate pairs โดยแสดง precision@100 precision@500 precision@1000 และ precision@5000 ร่วมกับค่า Candidate AP และ ROC-AUC ................................................................................................ 70 รูปที่ 4.12 Retrieval funnel ของ Project 1 แสดงการลด search space จาก all cross- platform pairs ไปสู่ exact-first pairs model candidates final MATCH only และ REVIEW queue ....................................................................................................................................... 71 รูปที่ 4.13 โครงสร้างของ Project 1 ในขั้น decision thresholding เริ่มจาก exact-first และ blocking ก่อนส่ง candidate pairs เข้าสู่ calibrated candidate scoring แล้วแยกผลลัพธ์เป็น MATCH REVIEW และ NO\_MATCH ........................................................................................ 71 รูปที่ 4.14 ผลของ thresholding ต่อปริมาณงานและคุณภาพการตัดสินใจ แสดงจำนวนคู่ในแต่ละ tier precision ของกลุ่ม auto-merge และ recall contribution ........................................... 73 รูปที่ 4.15 Production decision matrix ที่ operating point สุดท้าย แสดงความสัมพันธ์ระหว่าง Actual MATCH และ Actual NO\_MATCH กับผลการตัดสินใจใน tier MATCH REVIEW และ NO\_MATCH ................................................................................. 74 รูปที่ 4.16 การวิเคราะห์ข้อผิดพลาดจากตัวอย่าง false positives และ false negatives แสดงการกระจุกตัวตามคู่แพลตฟอร์มและลักษณะของสัญญาณชื่อ ภาพ และ caption ที่เกี่ยวข้องกับความผิดพลาด ...................................................................................................... 75 รูปที่ 4.17 ความสำคัญของคุณลักษณะทั้ง 41 ตัวใน main run แสดงให้เห็นว่า feature กลุ่มชื่อและ username เป็นหลักฐานหลักของระบบ ขณะที่ feature จำนวนหนึ่งมีผลต่อการตัดสินใจต่ำมากหรือเป็นศูนย์ในรอบนี้ .................................................... 77 รูปที่ 4.18 การรวม feature importance ตามตระกูลหลัก แสดงสัดส่วนของหลักฐานจากชื่อ bio style image-cross-signal และหลักฐานประเภทอื่นที่แบบจำลองใช้จริงใน main run ..... 78 รูปที่ 4.19 Coverage ของข้อมูลภาพใน main run แยกระดับ profile และ test pair แสดงให้เห็นว่าการทดลองในรอบนี้ยังเป็น partial multimodal setting มากกว่า full image- to-image matching ................................................................................................................. 78 รูปที่ 4.20 สรุปเส้นทางการพัฒนาระบบจาก Leakage diagnosis ไปสู่ Final multimodal .... 79\par
% ---- source PDF page 16 ----
\clearpage
\chapter*{บทที่ 1}
\begin{center}\Large\bfseries บทนํา\end{center}
\section*{1.1 ที่มาและความสําคัญ}
จุดเริ่มต้นของโครงงานนี้เกิดขึ้นจากการศึกษางานวิจัยด้าน cross-platform user identification ซึ่งพบว่าปัญหาหนึ่งที่พบได้บ่อยในข้อมูลโซเชียลมีเดียคือบุคคลคนเดียวกันอาจ ปรากฏอยู่บนหลายแพลตฟอร์มพร้อมกัน ไม่ว่าจะเป็น Twitter, Instagram หรือ Google+ โดย ใช้ชื่อผู้ใช้และรูปโปรไฟล์ที่แตกต่างกัน แม้เนื้อหาที่โพสต์จะบ่งชี้ว่าเป็นบุคคลเดียวกัน แต่การ พิสูจน์โดยไม่ต้องตรวจสอบทีละรายการด้วยมือนั้นยังคงเป็นโจทย์ที่ท้าทาย โครงงานนี้ใช้ Dataset-LinkSocial ที่รวบรวมและเผยแพร่โดย Sharma และ Dyreson [7] ซึ่งประกอบด้วยโปรไฟล์จาก Twitter, Instagram และ Google+ Archive รวม 36,807 โปร ไฟล์ ข้อมูลใน dataset นี้มีลักษณะไม่สมบูรณ์โดยธรรมชาติ โดยเฉพาะข้อมูลจาก Google+ ที่ปิด ตัวลงในปี 2019 ทำให้ Post ที่ยังเข้าถึงได้ ข้อมูลการมีปฏิสัมพันธ์ และ Metadata บางส่วนมี ความไม่ครบถ้วน ซึ่งจริง ๆ แล้วเป็นสถานการณ์ที่สะท้อนความเป็นจริงของงานด้านนี้ ไม่ใช่เพียง ข้อจำกัดของ dataset เพียงอย่างเดียว ในมุมขององค์กร ปัญหาเดียวกันนี้เกิดขึ้นในรูปแบบของ Data Silo เมื่อข้อมูลลูกค้าถูก เก็บแยกกันในระบบหลายระบบ ไม่ว่าจะเป็น Website, แอปพลิเคชัน, Email หรือโซเชียลมีเดีย โดยแต่ละระบบมีรูปแบบข้อมูลและ Identifier ที่แตกต่างกัน ทำให้ไม่สามารถรู้ได้ว่าลูกค้าที่ซื้อ สินค้าบน Website กับผู้ใช้ที่แสดงความเห็นใน Instagram เป็นคนเดียวกันหรือไม่ ปัญหานี้ส่งผล โดยตรงต่อคุณภาพของการวิเคราะห์พฤติกรรมลูกค้าและประสิทธิภาพของแคมเปญการตลาด โครงงานนี้จึงมีเป้าหมายในการสร้าง Proof-of-Concept ของระบบ Identity Resolution โดยนำ Dataset-LinkSocial มาใช้เป็นกรณีศึกษา กระบวนการสกัด Feature ครอบคลุมทั้งมิติข้อความ รูปภาพ และรูปแบบพฤติกรรม แล้วนำไปฝึกแบบจำลอง Ensemble ที่ ตัดสินใจว่าโปรไฟล์สองชุดใดเป็นบุคคลเดียวกัน ผล Unified Profile ที่ได้จากกระบวนการนี้ถูก นำไปใช้เป็นข้อมูลตั้งต้นสำหรับแบบจำลอง Lead Scoring เพื่อทดสอบว่าการรวมข้อมูลจาก หลายแพลตฟอร์มจะให้การประเมินโอกาส Conversion ที่แม่นยำกว่าการใช้ข้อมูลจากแหล่งเดียว หรือไม่\par
\section*{1.2 วัตถุประสงค์}
\noindent 1. เพื่อพัฒนาระบบที่สามารถรวบรวมรายชื่อลูกค้า จากทุกแหล่งข้อมูลให้อยู่ในรูปแบบโปร\par
ไฟล์เดียว เพื่อขจัดปัญหาข้อมูลซ้ำซ้อน\par
% ---- source PDF page 17 ----
\clearpage
\noindent 2. เพื่อประยุกต์ใช้เทคโนโลยี AI และ Machine Learning ในการตรวจจับความสัมพันธ์ของ\par
ข้อมูล และวิเคราะห์เส้นทางของผู้บริโภค 3. เพื่อเป็นส่วนหนึ่งของการสร้างระบบ Agentic CRM ที่ช่วยสนับสนุนการทำงานของทีม\par
ขายและการตลาดให้มีความรวดเร็ว แม่นยำ และมีข้อมูลครบถ้วนสำหรับการตัดสินใจ 4. เพื่อให้ระบบ AI สามารถให้คำแนะนำในการติดตามผลและกำหนดกิจกรรม/Script ที่ควร\par
ทำ เป็นตัวช่วยให้ทีมขาย-การตลาดทํางานได้รวดเร็ว แม่นยำ และมีข้อมูลครบถ้วน\par
\section*{1.3 วิธีการดําเนินงาน}
\noindent 1. รวบรวมและจัดเตรียมข้อมูล โดยใช้ชุดข้อมูลจาก Dataset-LinkSocial [7] ซึ่งเป็นข้อมูล\par
สาธารณะที่เผยแพร่ผ่าน GitHub เพื่อนำมาใช้เป็นข้อมูลตั้งต้นในการพัฒนาและทดสอบ ระบบ ก่อนเข้าสู่กระบวนการตรวจสอบคุณภาพและความสมบูรณ์ของข้อมูล 2. ทำความสะอาดและจัดระเบียบข้อมูลให้อยู่ในรูปแบบเป็นอันหนึ่งอันเดียวกัน ที่พร้อมใช้\par
งานสำหรับ AI โดยการรับและจัดเตรียมข้อมูลจำลอง ก่อนจะนำเข้าและการตรวจสอบ เบื้องต้น 3. ทำ Data Standardization และ Normalization และจัดเก็บข้อมูลที่ผ่านกระบวนการ\par
แล้ว 4. ทำการประยุกต์ใช้ AI/ML เพื่อวิเคราะห์ข้อมูลที่สะอาดแล้วเพื่อสร้างโปรไฟล์ลูกค้าที่\par
สมบูรณ์และเป็น Agentic 4.1. ทำการสร้าง Feature Set จากข้อมูลทั้งหมดที่มี เช่น Name , Username , Bio\par
และ URL ของ Social Media 4.2. ทำการพัฒนาแบบจำลองสำหรับจับคู่ข้อมูลลูกค้าที่อาจเป็นบุคคลเดียวกัน โดยที่\par
ระบบจะคำนวณคะแนนความคล้ายกันและแปลงเป็นเปอร์เซ็นต์ความมั่นใจออกมา 4.3. หากคะแนนความมั่นใจเกิน Threshold ที่กำหนดระบบจะแสดงรายการ\par
Candidate Profiles ที่น่าจะเป็นคนเดียวกันบนหน้าจอ UI พร้อมแสดงข้อมูล เปรียบเทียบแบบ Side-by-side และค่า \% ความมั่นใจ และผู้ตัดสินใจเลือก Merge Profile หรือไม่ใช่คนเดียวกัน จากนั้นจะบันทึกผลการตัดสินใจของผู้ใช้ เพื่อนำไปใช้ ปรับปรุงแบบจำลอง และบันทึก Audit Log ทุกครั้งที่มีการ Merge 5. นำแบบจำลองที่ได้มาวิเคราะห์พฤติกรรมของลูกค้าที่ถูกรวมโปรไฟล์เข้าด้วยกัน และทำ\par
การจัดประเภทลูกค้าว่ากำลังอยู่ช่วงใดของ Funnel 6. คำนวณคะแนนสำหรับลูกค้าแต่ละราย เพื่อหาลูกค้าที่เป็นกลุ่มเป้าหมายโดยอ้างอิงข้อมูล\par
จากหลายแหล่ง เช่น พฤติกรรม Engagement ประวัติการสื่อสาร แหล่งที่มา และความ พร้อมซื้อ\par
% ---- source PDF page 18 ----
\clearpage
\noindent 7. ประยุกต์ใช้แบบจำลองที่พัฒนาขึ้นมากับข้อมูลสถานะของลูกค้า และการให้คะแนนแก่\par
ลูกค้าเป้าหมายเพื่อแนะนำการดำเนินการที่เหมาะสมที่สุดแก่ทีมขาย\par
\section*{1.4 ขอบเขตของงาน}
\noindent 1. โครงงานใช้ชุดข้อมูลจาก Dataset-LinkSocial [7] ซึ่งเป็นข้อมูลสาธารณะ โดยข้อมูล\par
ได้รับการทำความสะอาด (Data Cleaning) และจัดให้อยู่ในรูปแบบมาตรฐาน (Data Preprocessing) ก่อนนำไปใช้กับแบบจำลอง 2. ระบบครอบคลุมการพัฒนาแบบจำลองรวมตัวตนลูกค้า (Identity Resolution) เพื่อระบุ\par
ว่าโปรไฟล์ใดเป็นบุคคลเดียวกัน และสร้างโปรไฟล์ลูกค้าที่สมบูรณ์และเป็น Agentic 3. ระบบวิเคราะห์พฤติกรรมลูกค้าและจัดประเภทเส้นทางลูกค้า (Customer Journey)\par
ตามขั้นของ Funnel โดยอาศัยข้อมูลจาก Interaction และ Engagement 4. ระบบพัฒนาแบบจำลองให้คะแนนศักยภาพลูกค้า (Lead Scoring) เพื่อประเมินระดับ\par
ความพร้อมซื้อและจัดลำดับความสำคัญของลูกค้าในระบบต้นแบบ 5. ระบบสามารถแสดง Candidate Profiles ที่มีความเป็นไปได้ว่าเป็นบุคคลเดียวกัน พร้อม\par
เปอร์เซ็นต์ความมั่นใจ และเปิดโอกาสให้ผู้ใช้ตัดสินใจรวมโปรไฟล์หรือแยกเป็นคนละ บุคคล โดยมีการบันทึกข้อมูลลงใน Audit Log 6. ระบบแสดงผลการวิเคราะห์ Journey Stage คะแนน Lead Score และคำแนะนำ\par
เบื้องต้นสำหรับการติดตามลูกค้า (Recommended Action) เพื่อสาธิตความสามารถ ของแบบจำลองและระบบ AI 7. ระบบพัฒนาในรูปแบบเว็บต้นแบบ (Web-based Prototype) เพื่อแสดงผลการทำงาน\par
ของแบบจำลองและส่วนติดต่อผู้ใช้ โดยไม่เชื่อมต่อกับระบบบริหารจัดการความสัมพันธ์ ลูกค้า (CRM) จริง\par
\section*{1.5 ประโยชน์ที่คาดว่าจะได้รับ}
\noindent 1. องค์กรมีฐานข้อมูลลูกค้าที่เป็นมาตรฐานเดียวลดความซ้ำซ้อนและความผิดพลาดของ\par
ข้อมูล 2. ทีมขายและการตลาดสามารถเข้าใจพฤติกรรมลูกค้าผ่าน Customer Journey ที่\par
วิเคราะห์โดย AI นำไปสู่กลยุทธ์ที่แม่นยำขึ้น 3. ลดระยะเวลาในการเตรียมข้อมูล ทำให้ทีมงานโฟกัสกับงานเชิงกลยุทธ์และการดูแลลูกค้า\par
ได้เต็มที่\par
% ---- source PDF page 19 ----
\clearpage
\section*{1.6 แผนการดำเนินงาน}
\begin{center}\small ตารางที่ 1.1\end{center}
การด าเนินงาน ธันวาคม 2567 มกราคม 2568 กุมภาพันธ์ 2568 มีนาคม 2568\par
1 2 3 4 1 2 3 4 1 2 3 4 1 2 3 4 1. การวางแผน\par
<····> <=> 1.2 ก ำหนดแผนกำรท ำงำน <····> < == > 2. การศึกษาและรวบรวมข้อมูล\par
\noindent 1.1 ก ำหนดวัตถุประสงค์และขอบเขต\par
<········> < ==== > 2.2 ศึกษำเทคนิค String Similarity และ Blocking\par
\noindent 2.1 ศึกษำพื้นฐำน Identity Resolution\par
<········> < ==== > 2.3 ศึกษำ Deep Learning ส ำหรับ Classification\par
<········> < === > 3. การเตรียมข้อมูล\par
<········> < ==== > 3.2 Blocking และ Entity- Aware Labeling\par
\noindent 3.1 สร้ำง Synthetic Dataset (3 แพลตฟอร์ม )\par
<········> < === > 4. การด าเนินงาน\par
<············> < ====== > 4.2 พัฒนำ Decision Thresholding และ\par
\noindent 4.1 ออกแบบและพัฒนำโมเดล IdentityMLP\par
<············> < ====== > 5. การทดสอบและประเมินผล\par
\noindent 5.1 ทดสอบโมเดลบน Test Set <············> < ========= > 5.2 เปรียบเทียบผล Thresholding รูปแบบต่ำงๆ\par
<········> < ====== > 5.3 ปรับปรุงโมเดลและ Pipeline\par
<········> < ===== > 6. สรุปผลและรายงานผล\par
\noindent 6.1 จัดท ำรำยงำนผล <········> < ===== > 6.2 น ำเสนอโครงงำน <····> <=> < ...... > = แผนที่วางไว้ <====> = ผลการดำเนินงานจริง\par
% ---- source PDF page 20 ----
\clearpage
% BEGIN CHAPTER 1 GA ADDITIONS
% ส่วนเพิ่มเติมบทที่ 1 ต่อจากหัวข้อ 1.6
\section*{1.7 ที่มาและความสำคัญของการพัฒนานโยบายการตัดสินใจเพิ่มเติม}
การพัฒนาระบบเชื่อมโยงอัตลักษณ์ในระยะแรกของโครงงานมุ่งให้ระบบสามารถนำข้อมูลโปรไฟล์จากหลายแพลตฟอร์มมาเปรียบเทียบและให้คะแนนความเป็นบุคคลเดียวกันได้ งานที่พัฒนาเพิ่มเติมในระยะต่อมามุ่งตอบคำถามว่า เมื่อมีคะแนนแล้ว ระบบควรตัดสินใจอย่างไรจึงจะเหมาะสมกับการนำข้อมูลไปใช้ใน CRM เพราะการมีแบบจำลองที่แยกคู่เหมือนและคู่ต่างได้ดี ยังไม่เพียงพอที่จะกำหนดว่าคู่ใดควรถูกรวมทันที และคู่ใดควรรอการตรวจสอบ

หัวข้อ 1.7--1.13 จึงขยายรายละเอียดจากที่มา วัตถุประสงค์ ขอบเขต และวิธีดำเนินงานเดิม โดยเน้นส่วนที่โครงงานได้พัฒนาเพิ่ม คือการค้นหานโยบายการตัดสินใจด้วยขั้นตอนวิธีเชิงพันธุกรรม การใช้หลักฐานหลายด้านอย่างเหมาะสมเมื่อข้อมูลไม่ครบ และการประเมินผลที่คำนึงถึงทั้งคุณภาพการเชื่อมโยงและภาระงานตรวจสอบ

\subsection*{1.7.1 จากปัญหาข้อมูลซ้ำสู่ความเสี่ยงของการรวมข้อมูลผิด}
ข้อมูลลูกค้าที่กระจายอยู่หลายระบบทำให้การมองเห็นภาพรวมของลูกค้าไม่ครบถ้วน แต่การพยายามรวมข้อมูลให้ได้มากที่สุดก็อาจสร้างปัญหาใหม่ หากระบบนำโปรไฟล์ของคนละคนมาอยู่ภายใต้ตัวตนเดียวกัน ข้อมูลความสนใจ ช่องทางติดต่อ หรือประวัติที่นำไปใช้ประกอบการทำงานอาจปะปนกัน ผลของการรวมผิดจึงไม่ได้จำกัดอยู่ที่คำตอบหนึ่งแถวในตาราง แต่สามารถส่งต่อไปยังกระบวนการที่ใช้โปรไฟล์รวมนั้นด้วย

ในอีกด้านหนึ่ง หากระบบระมัดระวังมากจนปฏิเสธคู่ที่เป็นคนเดียวกัน ข้อมูลของลูกค้าก็ยังคงแยกอยู่หลายรายการ การวิเคราะห์หรือจัดลำดับลูกค้าอาจใช้ข้อมูลเพียงบางส่วน การเชื่อมโยงจึงต้องพิจารณาสมดุลระหว่างการป้องกันการรวมผิดกับการลดคู่ที่ตกหล่น โดยไม่มีเหตุผลที่จะสมมติว่าความผิดพลาดสองประเภทสร้างภาระเท่ากันในทุกบริบท

ตัวอย่างเช่น คู่โปรไฟล์สองคู่ได้รับคะแนนจากแบบจำลองใกล้เคียงกัน คู่แรกมีชื่อคล้ายกันและมีลิงก์อ้างอิงสอดคล้องกัน ส่วนคู่ที่สองมีเพียงชื่อที่คล้ายกันและไม่มีหลักฐานอื่น การใช้จุดตัดเดียวกับคะแนนเพียงค่าเดียวอาจให้คำตัดสินเหมือนกันทั้งสองคู่ แม้ข้อมูลประกอบต่างกัน ตัวอย่างนี้เป็นเหตุผลในการศึกษาว่าการนำหลักฐานเพิ่มเติมเข้ามาในชั้นการตัดสินใจช่วยให้ระบบจัดการกรณีก้ำกึ่งได้เหมาะสมขึ้นหรือไม่

\subsection*{1.7.2 ข้อจำกัดของการกำหนดเกณฑ์ด้วยค่าคงที่}
ระบบเดิมใช้คะแนนและเกณฑ์สองระดับเพื่อแยกคู่เป็น MATCH, REVIEW และ NO\_MATCH แนวทางนี้ช่วยให้ไม่ต้องบังคับตัดสินทุกคู่เป็นใช่หรือไม่ใช่ทันที แต่ค่าจุดตัดยังมีผลโดยตรงต่อจำนวนคู่ที่ระบบรวมเองและจำนวนคู่ที่ต้องส่งให้คนตรวจสอบ หากตั้งเกณฑ์ MATCH สูงขึ้น จำนวนการรวมอัตโนมัติอาจลดลง ขณะที่คู่ส่วนหนึ่งย้ายไปอยู่ในคิว REVIEW หากลดเกณฑ์ลง ระบบอาจจับคู่ได้มากขึ้นพร้อมกับความเสี่ยงที่จะรวมผิดเพิ่มขึ้น

ภายใต้นโยบายที่ใช้คะแนนเดียวและเกณฑ์สองระดับ ความกว้างของช่วง REVIEW กำหนดพื้นที่ที่ระบบยอมเว้นการตัดสินใจอัตโนมัติ การขยายช่วงนี้อาจเพิ่มโอกาสให้เจ้าหน้าที่ช่วยแก้กรณียาก แต่ต้องแลกกับเวลาทำงาน การลดจำนวน REVIEW จึงเป็นเป้าหมายที่ต้องอ่านร่วมกับผลของคู่ที่ย้ายออกจากคิว ไม่ใช่นับเพียงว่าคิวสั้นลงแล้วถือว่าระบบดีขึ้น

ด้วยเหตุนี้ ผู้วิจัยจึงเพิ่มขั้นตอนค้นหาเกณฑ์จากข้อมูล validation และกำหนดต้นทุนสำหรับความผิดพลาดกับการตรวจสอบไว้ล่วงหน้า การเปรียบเทียบนโยบายบนข้อมูลเดียวกันทำให้สามารถอธิบายได้ว่าผลที่เปลี่ยนแปลงเกิดจากการปรับวิธีตัดสินใจมากน้อยเพียงใด

\subsection*{1.7.3 เหตุผลที่ใช้ GA และหลักฐานหลายด้าน}
นโยบายที่พัฒนาระยะแรกเพิ่มการใช้ความคล้ายของชื่อช่วยตัดสินคู่ในช่วง REVIEW ทำให้การค้นหามีพารามิเตอร์มากกว่าจุดตัดของคะแนนโมเดลเพียงสองค่า ต่อมาเมื่อขยายไปใช้หลักฐานหลายด้าน จำนวนพารามิเตอร์เพิ่มเป็นเกณฑ์สองระดับและน้ำหนักแปดค่า การเลือกค่าทีละตัวด้วยมืออาจไม่สะท้อนผลที่แต่ละค่ามีร่วมกัน เช่น การเพิ่มน้ำหนักชื่ออาจต้องเปลี่ยนเกณฑ์ MATCH จึงจะรักษาระดับความแม่นยำเดิมได้

GA เป็นทางเลือกหนึ่งสำหรับการค้นหาชุดพารามิเตอร์ร่วมกัน โดยสร้างนโยบายหลายชุด ประเมินผล แล้วนำชุดที่เหมาะสมไปสร้างคำตอบในรุ่นถัดไป การใช้ GA กับงานเชื่อมโยงข้อมูลมีตัวอย่างในงานของ Sun และคณะ [11] ซึ่งศึกษาการเรียนรู้กฎจับคู่ร่วมกับ active learning สำหรับโครงงานนี้ ผู้วิจัยประยุกต์ GA กับชั้นการตัดสินใจหลังได้คะแนนจากแบบจำลอง โดยตรึงคะแนนไว้ระหว่างเปรียบเทียบนโยบาย

อีกประเด็นที่ต้องพิจารณาคือความพร้อมของข้อมูล โปรไฟล์บางคู่มีประวัติย่อ ลิงก์ หรือภาพ ขณะที่บางคู่ไม่มี การไม่มีข้อมูลไม่ได้ให้ความหมายเดียวกับการมีข้อมูลที่ขัดแย้งกัน ผู้วิจัยจึงเพิ่มตัวบ่งชี้ความพร้อมของหลักฐาน และคำนวณคะแนนถ่วงน้ำหนักจากช่องที่มีข้อมูลจริงในแต่ละคู่ เพื่อศึกษานโยบายที่รองรับสภาพข้อมูลไม่ครบได้อย่างชัดเจน

\section*{1.8 วัตถุประสงค์ของงานที่พัฒนาเพิ่มเติม}
วัตถุประสงค์ในส่วนนี้ขยายจากเป้าหมายการเชื่อมโยงข้อมูลและเตรียมข้อมูลสำหรับ CRM ที่กล่าวไว้ในหัวข้อ 1.2 โดยมุ่งให้แต่ละวัตถุประสงค์มีผลลัพธ์ที่ตรวจสอบได้จากการทดลอง ดังนี้

\subsection*{1.8.1 พัฒนานโยบายการตัดสินใจที่ค้นหาพารามิเตอร์จากข้อมูล}
เพื่อพัฒนาวิธีเลือกเกณฑ์สำหรับ MATCH, REVIEW และ NO\_MATCH ด้วย GA โดยใช้ข้อมูล validation เป็นแหล่งเลือกพารามิเตอร์ และตรึงนโยบายก่อนประเมินกับ test ผลงานที่ต้องได้คือชุดพารามิเตอร์ที่นำกลับมาใช้ได้ พร้อมข้อมูลการตั้งค่าและผล validation ที่ใช้ประกอบการเลือก

\subsection*{1.8.2 ขยายการใช้หลักฐานในชั้นการตัดสินใจ}
เพื่อพัฒนานโยบายที่รวมคะแนนโมเดลกับหลักฐานด้านชื่อ ประวัติย่อ ลิงก์ ตำแหน่งที่อยู่ สัญญาณทางสังคม ภาพ และรูปแบบการเขียนตามข้อมูลที่มี โดยแยกค่าความคล้ายออกจากสถานะการมีข้อมูล ผลงานที่ต้องได้คือตารางหลักฐานรายคู่และสูตรการคำนวณที่ไม่ลดคะแนนเพียงเพราะช่องข้อมูลขาดหาย

\subsection*{1.8.3 เปรียบเทียบผลด้านคุณภาพและภาระงานอย่างเป็นระบบ}
เพื่อเปรียบเทียบระบบเดิม นโยบายคะแนนอย่างเดียว นโยบายใช้ชื่อช่วยตัดสิน และนโยบายแบบถ่วงน้ำหนัก ภายใต้คู่ข้อมูลและนิยามต้นทุนเดียวกัน โดยพิจารณา precision, recall, F1 จำนวนการรวมผิด จำนวนคู่ที่ปฏิเสธผิด และจำนวน REVIEW ร่วมกัน

\subsection*{1.8.4 จัดทำผลลัพธ์ที่ตรวจสอบย้อนกลับและส่งต่อได้}
เพื่อบันทึกแหล่งข้อมูล รุ่นนโยบาย ค่าเกณฑ์ น้ำหนัก และคำตัดสินรายคู่ให้สัมพันธ์กัน รองรับการตรวจสอบความแตกต่างจากระบบเดิม และเตรียมผลสำหรับกระบวนการรวมโปรไฟล์หรือจัดคิวตรวจสอบใน CRM โดยไม่ถือว่าการส่งออกผลเท่ากับการนำไปใช้งานจริงครบวงจรแล้ว

\section*{1.9 คำถามการศึกษาและแนวทางหาคำตอบ}
คำถามของงานเพิ่มเติมมุ่งตรวจสอบผลของการออกแบบนโยบายภายใต้ข้อมูลและข้อจำกัดที่โครงงานมีอยู่ การศึกษาไม่ได้ตั้งข้อสรุปล่วงหน้าว่า GA ต้องดีกว่าทุกแนวทาง หรือการเพิ่มหลักฐานต้องทำให้ทุกตัวชี้วัดสูงขึ้น

\begin{center}\small
\begin{tabularx}{\linewidth}{p{1.1cm}XX}
\toprule
คำถาม & ประเด็นที่ต้องการศึกษา & วิธีตรวจสอบ \\
\midrule
RQ1 & การค้นหาเกณฑ์ด้วย GA เปลี่ยนต้นทุนและภาระ REVIEW จากระบบเดิมอย่างไร & เปรียบเทียบระบบเดิมกับนโยบายคะแนนอย่างเดียวบนคู่ชุดเดียวกัน \\
RQ2 & การรวมหลักฐานหลายด้านให้ผลต่างจากการใช้ชื่อช่วยตัดสินอย่างไร & เปรียบเทียบนโยบายสี่ค่ากับนโยบายสิบค่า โดยตรึงคะแนนโมเดล \\
RQ3 & ความพร้อมของข้อมูลส่งผลต่อการใช้น้ำหนักและการอ่านผลอย่างไร & รายงาน availability, coverage และน้ำหนักของกลุ่มที่เปิดใช้จริง \\
RQ4 & นโยบายที่เลือกจาก validation ให้ผลบน test และข้อมูลรวมต่างกันอย่างไร & ตรึงพารามิเตอร์ก่อนทดสอบ และรายงานแต่ละกลุ่มข้อมูลแยกกัน \\
\bottomrule
\end{tabularx}
\par ตารางที่ 1.2 คำถามการศึกษาของงานที่พัฒนาเพิ่มเติม
\end{center}

ผลที่ใช้ตอบคำถามควรแสดงการแลกเปลี่ยนระหว่างตัวชี้วัดอย่างตรงไปตรงมา เช่น หากต้นทุนลดลงแต่ recall ลดลงด้วย ต้องอธิบายว่าสัดส่วนคู่ที่ย้ายไปเป็น NO\_MATCH หรือ REVIEW เปลี่ยนอย่างไร การเปลี่ยนแปลงดังกล่าวเป็นข้อมูลสำหรับเลือกนโยบายตามบริบท มากกว่าจะใช้ประกาศว่านโยบายหนึ่งเหมาะสมกับทุกสถานการณ์

\section*{1.10 ขอบเขตและเงื่อนไขของงานเพิ่มเติม}
\subsection*{1.10.1 ขอบเขตข้อมูลและหน่วยที่ใช้ศึกษา}
งานใช้ข้อมูลโปรไฟล์ข้ามแพลตฟอร์มที่ผ่านกระบวนการจัดเตรียมในโครงงาน โดยมีแหล่งตั้งต้นสัมพันธ์กับ Dataset-LinkSocial [7] ซึ่งเผยแพร่ข้อมูลจาก Twitter, Instagram และ Google+ หน่วยที่นโยบายตัดสินคือคู่โปรไฟล์ ไม่ใช่ลูกค้าหนึ่งรายหรือยอดขายหนึ่งรายการ จำนวนแถวในข้อมูลที่ประมวลผลแล้วจึงต้องอ้างอิงไฟล์ของรอบทดลอง และแยกจากจำนวนที่รายงานในคลังข้อมูลสาธารณะ

ในรอบข้อมูลเต็มที่ใช้เป็นหลักฐานของงานเพิ่มเติม มีคู่ 2,086,245 คู่ ประกอบด้วยคู่ที่ใช้คะแนนโมเดล 2,073,842 คู่ และคู่จาก Exact-First 12,403 คู่ การค้นหานโยบายใช้เฉพาะคู่ validation ในส่วนที่ผ่านการให้คะแนนโมเดล ส่วนคู่ Exact-First คงคำตัดสินเดิมไว้เพื่อให้ขอบเขตการเปลี่ยนแปลงชัดเจน

\subsection*{1.10.2 ขอบเขตของแบบจำลองและนโยบาย}
การทดลองนโยบายเปรียบเทียบวิธีนำคะแนนไปตัดสินใจ โดยไม่ได้ฝึกแบบจำลองใหม่ระหว่างการเปรียบเทียบแต่ละนโยบายในรอบเดียวกัน โครงสร้างแบบถ่วงน้ำหนักรองรับข้อมูลแปดช่อง แต่ช่องที่ไม่มีหลักฐานพร้อมใช้จะไม่ถูกนำมาค้นหาน้ำหนัก ในรอบข้อมูลเต็มที่รายงาน location และ social ไม่ได้ใช้งานตามสภาพของตารางข้อมูล จึงไม่อ้างว่าเป็นการใช้ครบทุกด้าน

การศึกษานี้ไม่ได้เพิ่มคู่ผู้สมัครที่หลุดจากขั้น candidate generation กลับเข้ามา ผลของ GA จึงอยู่ภายในขอบเขตคู่ที่ระบบส่งมาให้ตัดสิน ความครอบคลุมของขั้นสร้างคู่ยังเป็นข้อจำกัดที่ต้องรายงานแยกจากความสามารถของนโยบาย

\subsection*{1.10.3 ขอบเขตการประเมินต้นทุน}
ต้นทุนกำหนดจากสมมติฐานเวลาทำงานและค่าแรงของโครงงาน โดยให้การรวมผิดมีน้ำหนัก 2 การปฏิเสธคู่บวกมีน้ำหนัก 1 และ REVIEW มีน้ำหนัก 0.1 ไม่รวมรายได้ที่สูญเสีย ผลกระทบต่อชื่อเสียง หรือค่าใช้จ่ายของการแก้ไขระบบที่ไม่ได้มีข้อมูลวัดจริง การเปรียบเทียบจึงใช้ตอบว่านโยบายใดมีต้นทุนต่ำกว่าภายใต้สมมติฐานเดียวกัน

การเลือกนโยบายใช้ validation ส่วน test ใช้ประเมินหลังตรึงค่า การรายงานผลบนคู่ทั้งหมดมีไว้แสดงปริมาณงานของข้อมูลชุดนั้น และไม่ใช้แทนผลบน test นอกจากนี้ การทดสอบหลาย seed เป็นการตรวจผลจากการสุ่มค้นหา ไม่เท่ากับการทดสอบกับข้อมูลจากหลายองค์กรหรือหลายช่วงเวลา

\subsection*{1.10.4 ขอบเขตการเชื่อมต่อกับงาน CRM}
ผลที่จัดเตรียมสำหรับกระบวนการปลายทาง ได้แก่ รหัสคู่ คะแนน คำตัดสิน และข้อมูลที่ใช้ติดตามรุ่นนโยบาย ส่วนแนวคิดด้านการจัดลำดับลูกค้า การวิเคราะห์เส้นทางลูกค้า และการแนะนำกิจกรรมยังเป็นบริบทของโครงงานภาพรวม การปรับ decision policy ครั้งนี้ไม่ได้ใช้ข้อมูลการซื้อจริงมาทดสอบว่าการรวมโปรไฟล์ทำให้อัตราการซื้อเพิ่มขึ้น

หากนำผลไปใช้สร้างโปรไฟล์รวม ต้องตรวจความถูกต้องในระดับกลุ่มบุคคลเพิ่มเติม เพราะการจับคู่หลายคู่ต่อกันอาจทำให้ความผิดพลาดเพียงคู่เดียวเชื่อมกลุ่มข้อมูลที่ไม่ควรอยู่ร่วมกันได้ ผลดีในระดับคู่จึงยังไม่เพียงพอสำหรับยืนยันความถูกต้องของโปรไฟล์รวมทุกกลุ่ม

\section*{1.11 ภาพรวมวิธีดำเนินงานและการเชื่อมโยงกับงานเดิม}
งานเพิ่มเติมเริ่มจากทบทวนคะแนนและคำตัดสินของระบบเดิม เพื่อกำหนดฐานเปรียบเทียบ จากนั้นตรวจรหัสคู่และการแบ่งข้อมูล สร้างตารางหลักฐานพร้อมตัวบ่งชี้ความพร้อมใช้งาน แล้วกำหนดนโยบายและต้นทุนที่ใช้เปรียบเทียบให้ตรงกัน การค้นหาด้วย GA ดำเนินการบน validation ก่อนบันทึกพารามิเตอร์ที่เลือกและนำไปประเมิน test

หลังทดสอบ ผู้วิจัยนำพารามิเตอร์เดิมไปคำนวณคำตัดสินของคู่ทั้งหมด คงผล Exact-First ไว้ และส่งออกผลแยกตามนโยบาย การวิเคราะห์ปลายทางพิจารณาจำนวนคู่แต่ละสถานะ ความผิดพลาด และหลักฐานที่ใช้ประกอบคำตัดสิน รายละเอียดสูตรและการตั้งค่าการทดลองอธิบายไว้ในหัวข้อ 3.6--3.8

\begin{center}\small
\begin{tabularx}{\linewidth}{p{3cm}XX}
\toprule
องค์ประกอบ & งานเดิมที่ใช้ต่อ & งานที่เพิ่มในระยะนี้ \\
\midrule
ข้อมูลและคู่ผู้สมัคร & โปรไฟล์ที่จัดเตรียมแล้ว รหัสคู่ และผลสร้างคู่ & ตรวจความตรงกันของคู่กับตารางหลักฐานสำหรับนโยบาย \\
คะแนน & คะแนนจากแบบจำลองและคำตัดสินเดิม & ตรึงคะแนนเพื่อเปรียบเทียบวิธีตัดสินใจ \\
หลักฐาน & คุณลักษณะที่คำนวณไว้ในโครงงาน & จัดกลุ่มคะแนนและสถานะความพร้อมใช้งาน \\
การตัดสินใจ & MATCH, REVIEW และ NO\_MATCH & ค้นหาเกณฑ์และน้ำหนักด้วย GA \\
การประเมิน & ตัวชี้วัดคุณภาพการเชื่อมโยง & เพิ่มต้นทุน ภาระตรวจสอบ และผลแยกขอบเขตข้อมูล \\
การส่งออก & ผลจับคู่สำหรับงานปลายทาง & บันทึกรุ่นนโยบายและผลเปรียบเทียบรายคู่ \\
\bottomrule
\end{tabularx}
\par ตารางที่ 1.3 ความเชื่อมโยงระหว่างงานเดิมกับงานเพิ่มเติม
\end{center}

\section*{1.12 ประโยชน์ที่คาดว่าจะได้รับและผลงานที่ส่งมอบ}
ประโยชน์หลักคือการทำให้เหตุผลในการเลือกนโยบายชัดเจนขึ้น ผู้ใช้ผลการทดลองสามารถเห็นว่าการตั้งเกณฑ์และน้ำหนักทำให้การรวมผิด คู่ตกหล่น และงานตรวจสอบเปลี่ยนไปอย่างไร แทนการพิจารณาคะแนนจำแนกเพียงค่าเดียว นอกจากนี้ การแยกข้อมูลขาดออกจากค่าความคล้ายต่ำช่วยให้ตีความผลของหลักฐานแต่ละด้านได้รอบคอบขึ้น

ในด้านการพัฒนาระบบ การตรึงและบันทึกนโยบายช่วยให้เปรียบเทียบผลระหว่างรุ่นได้ เมื่อมีข้อมูลหรือสมมติฐานต้นทุนใหม่ สามารถนำกระบวนการเดิมกลับมาค้นหาและประเมินอีกครั้งได้โดยมีฐานอ้างอิง ผลงานที่ส่งมอบประกอบด้วยตารางหลักฐานรายคู่ ข้อมูลความครอบคลุม ผลค้นหาจากหลาย seed ไฟล์นโยบายที่ตรึงไว้ ตารางตัวชี้วัด และคำตัดสินรายคู่

ประโยชน์ด้านการลดเวลาตรวจสอบและการปรับคุณภาพข้อมูล CRM เป็นสิ่งที่คาดหวังจากการออกแบบ การยืนยันผลในองค์กรจริงยังต้องเก็บข้อมูลการตรวจสอบ การแก้ไขการรวมผิด และผลต่อกระบวนการทำงานปลายทางเพิ่มเติม จึงไม่ใช้ผลของการทดลองออฟไลน์แทนหลักฐานประโยชน์ทางธุรกิจที่เกิดขึ้นจริง

\section*{1.13 ลำดับการดำเนินงานและนิยามที่ใช้ในส่วนเพิ่มเติม}
การดำเนินงานแบ่งเป็นช่วงต่อเนื่อง ได้แก่ ตรวจฐานข้อมูลและผลเดิม เตรียมหลักฐาน ออกแบบนโยบาย ค้นหาพารามิเตอร์ ตรึงและทดสอบนโยบาย และจัดทำผลเปรียบเทียบ แต่ละช่วงต้องได้ข้อมูลที่ใช้เป็นจุดเริ่มต้นของช่วงถัดไป เช่น ต้องตรวจรหัสคู่และความพร้อมของหลักฐานก่อนค้นหาน้ำหนัก และต้องบันทึกนโยบายก่อนประเมิน test การเรียงงานเช่นนี้ทำให้ติดตามได้ว่าความคลาดเคลื่อนเกิดขึ้นในขั้นใด

ในส่วนเพิ่มเติมนี้ ``คะแนนโมเดล'' หมายถึงค่าที่แบบจำลองให้กับคู่โปรไฟล์ ส่วน ``คะแนนนโยบาย'' หมายถึงค่าที่คำนวณจากหลักฐานและน้ำหนักเพื่อใช้ตัดสินใจ ``นโยบายที่ตรึงไว้'' หมายถึงชุดเกณฑ์และน้ำหนักที่เลือกเสร็จก่อนขั้นประเมิน test และ ``ข้อมูลพร้อมใช้งาน'' หมายถึงมีหลักฐานที่สามารถคำนวณและใช้ตามนิยามของช่องนั้นได้ ไม่ได้หมายความว่าหลักฐานนั้นยืนยันความเป็นบุคคลเดียวกันเสมอไป

การใช้คำเหล่านี้อย่างแยกกันช่วยเชื่อมบทนำกับทฤษฎีในบทที่ 2 และวิธีดำเนินงานในบทที่ 3 โดยเฉพาะการแยกคุณภาพของแบบจำลองออกจากผลของนโยบาย และแยกผลการทดลองออกจากเป้าหมายการใช้งานในภาพรวมของโครงงาน

% END CHAPTER 1 GA ADDITIONS
\clearpage
\chapter*{บทที่ 2}
\begin{center}\Large\bfseries เอกสารและงานวิจัยที่เกี่ยวข้อง\end{center}
\section*{2.1 แนวคิดที่เกี่ยวข้อง}
\subsection*{2.1.1 Identity Resolution และ User Identity Linkage}
ในโลกของการจัดการข้อมูลลูกค้าผ่านแพลตฟอร์มดิจิทัลหลายช่องทาง ปัญหาพื้นฐานที่ องค์กรต้องเผชิญคือการระบุว่าบัญชีหรือระเบียนข้อมูลจากแหล่งต่าง ๆ นั้นเป็นของบุคคลเดียวกัน หรือไม่ กระบวนการนี้เรียกว่า Identity Resolution หรือ User Identity Linkage ซึ่งเป็น รากฐานสำคัญของการสร้าง Single Customer View อันเป็นเป้าหมายหลักของระบบบริหาร จัดการความสัมพันธ์ลูกค้า (CRM) สมัยใหม่ งานวิจัยในสาขานี้เริ่มเข้มข้นขึ้นตั้งแต่ราวปี 2010 และยังคงพัฒนาต่อเนื่องตามการขยายตัวของแพลตฟอร์มโซเชียลมีเดีย\par
Halimi และ Ayday [3] ได้ทบทวนเทคนิคและแนวทางสำหรับ User Identity Linkage บนเครือข่ายสังคมออนไลน์อย่างครอบคลุม พบว่าแนวทางหลักแบ่งได้เป็นสองกลุ่ม คือ Profile- based Approach ที่อาศัยข้อมูลโปรไฟล์ เช่น ชื่อ รูปภาพ และข้อมูลประชากร กับ Behavioral- based Approach ที่ใช้รูปแบบการใช้งานและพฤติกรรมออนไลน์เป็นสัญญาณในการระบุตัวตน น่าสังเกตว่างานวิจัยส่วนใหญ่ในกลุ่มนี้ถูกทดสอบบนแพลตฟอร์มที่ยังให้บริการอยู่และมีข้อมูล ครบถ้วน\par
กรณีของโครงงานนี้มีความแตกต่างที่สำคัญตรงนี้ เนื่องจาก Google+ ปิดตัวไปแล้วตั้งแต่ ปี 2019 Behavioral Data จากแพลตฟอร์มนั้นจึงไม่สามารถเก็บแบบ Real-time ได้อีกต่อไป แบบจำลองที่ออกแบบสำหรับสถานการณ์นี้จึงต้องพึ่งพา Profile-based Features เป็นหลัก และ ต้องออกแบบ Threshold Policy ให้ยอมรับความไม่แน่นอนได้สูงกว่าแบบจำลองใน สภาพแวดล้อมที่ข้อมูลครบถ้วน ซึ่งนำไปสู่แนวคิด POSSIBLE\_MATCH ที่ใช้ในการออกแบบระบบ นี้โดยตรง\par
\subsection*{2.1.2 Behavioral Modeling สำหรับการเชื่อมโยงบัญชีผู้ใช้ข้ามแพลตฟอร์ม}
Zafarani และ Liu [2] ได้เสนอแนวทาง Behavioral Modeling ซึ่งสร้างแบบจำลอง พฤติกรรมของผู้ใช้จากข้อมูลในแต่ละแพลตฟอร์ม เช่น ความถี่ในการโพสต์ คำศัพท์ที่ใช้บ่อย และ ช่วงเวลาที่ Active ของแต่ละวัน จากนั้นหาความคล้ายคลึงกันระหว่างแบบจำลองเพื่อตัดสินว่า เป็นบุคคลเดียวกันหรือไม่ ข้อได้เปรียบของแนวทางนี้คือสามารถทำงานได้แม้ข้อมูลโปรไฟล์ไม่ ครบถ้วน เพราะพฤติกรรมของผู้ใช้มักมีความสม่ำเสมอข้ามแพลตฟอร์ม\par
% ---- source PDF page 21 ----
\clearpage
Sharma และ Dyreson [1] ได้นำแนวคิดนี้ไปพัฒนาระบบ LINKSOCIAL ซึ่งผสมผสาน ทั้ง Profile-based และ Behavioral-based Features เพื่อเพิ่มความแม่นยำในการเชื่อมโยงโปร ไฟล์ข้ามแพลตฟอร์ม ผลการทดลองแสดงให้เห็นว่าการรวม Feature จากทั้งสองแหล่งให้ ประสิทธิภาพสูงกว่าการใช้แนวทางใดแนวทางหนึ่ง แม้ว่างานวิจัยนี้จะยังทดสอบบนแพลตฟอร์มที่ ข้อมูล Active อยู่ แต่หลักการของการรวม Feature ประเภทต่าง ๆ นั้นใช้ได้กับบริบทของ Archive Data ด้วยเช่นกัน\par
ในกรณีของ Google+ Archive ข้อจำกัดสำคัญคือ Behavioral Data เช่น Engagement Patterns และ Posting Frequency ไม่ครบถ้วนตามที่ควรจะเป็น ดังนั้นโครงงานนี้จึงปรับกล ยุทธ์มาเน้น Text Similarity และ Profile Metadata เป็นหลัก พร้อมทั้งนำ Temporal Features ที่ยังเหลืออยู่ใน Archive มาเสริม ซึ่งสอดคล้องกับข้อเสนอแนะของ Zafarani และ Liu ที่ระบุว่าการปรับน้ำหนัก Feature ตามความน่าเชื่อถือของข้อมูลแต่ละแหล่งเป็นสิ่งจำเป็น\par
\subsection*{2.1.3 Human-in-the-Loop สำหรับ Entity Matching}
การจับคู่ระเบียนข้อมูล (Entity Matching) โดยระบบอัตโนมัติย่อมมีข้อจำกัดในกรณีที่ ข้อมูลมีความกำกวม (Ambiguous) สูง Doan และคณะ [4] ได้รายงานความท้าทายและบทเรียน สำคัญจากการนำระบบ Human-in-the-Loop Entity Matching ไปใช้งานจริง ได้แก่ การ ออกแบบ Workflow ที่ลดภาระของผู้ตรวจสอบ การจัดการกับ Noisy Labels และการรักษา ความสม่ำเสมอในการตัดสินใจ สิ่งที่น่าสนใจคือบทเรียนเหล่านี้เกิดขึ้นจากระบบที่ใช้งานจริงใน ระดับ Production ไม่ใช่สภาพแวดล้อมทดลอง Li และคณะ [5] ได้นำเสนอสถาปัตยกรรม Human-in-the-Loop Data Integration ที่ ออกแบบมาเพื่อสมดุลระหว่างคุณภาพและต้นทุน โดยระบบจะใช้ Machine Learning ในการ ตัดสินใจโดยอัตโนมัติสำหรับกรณีที่มีความมั่นใจสูง และส่งต่อเฉพาะกรณีที่ไม่แน่ใจให้มนุษย์ พิจารณา แนวทางนี้ช่วยลดปริมาณงานที่มนุษย์ต้องตรวจสอบได้อย่างมีนัยสำคัญในขณะที่ยังคง รักษาความแม่นยำสูง แนวคิดนี้เชื่อมโยงโดยตรงกับการออกแบบระบบในโครงงาน ในขั้นตอน Merge Logic ได้ กำหนดให้แบบจำลองแบ่ง Match ออกเป็นสามระดับ ได้แก่ MATCH สำหรับกรณีที่ Confidence สูงซึ่งแบบจำลองตัดสินใจได้เอง POSSIBLE\_MATCH สำหรับกรณีกำกวมที่ต้องรอการยืนยัน และ NO\_MATCH สำหรับกรณีที่ชัดเจนว่าไม่ใช่บุคคลเดียวกัน แนวทางนี้เป็นการนำหลักการ Human- in-the-Loop มาปรับใช้ให้เหมาะสมกับบริบทของ Archive Data ที่มีข้อมูลไม่สมบูรณ์ 2.1.4 Lead Scoring และ Predictive CRM Lead Scoring เป็นกระบวนการกำหนดคะแนนให้กับลูกค้าเป้าหมาย (Lead) เพื่อ ประเมินโอกาสในการแปลงเป็นลูกค้าจริง (Conversion) Schafer, Hassan และ Muller [6] ได้ ทำการทดสอบเทคนิค Machine Learning หลายแบบสำหรับ Predictive Lead Scoring ใน\par
% ---- source PDF page 22 ----
\clearpage
บริบท B2B Sales ได้แก่ Random Forest, Gradient Boosting Machines (GBM), Logistic Regression และ Neural Networks บนชุดข้อมูลลูกค้า 50,000 ราย ผลการทดลองของ Schafer และคณะ [6] พบว่า GBM ให้ค่า AUC-ROC สูงสุดที่ 0.87 และที่สำคัญกว่านั้นคือการพิสูจน์ว่า Feature ที่รวมมาจากหลายแหล่งข้อมูล (Multi-source Identity Data) ให้ประสิทธิภาพสูงกว่า Feature จากแหล่งเดียวอย่างมีนัยสำคัญทางสถิติ ซึ่งตรง กับสมมุติฐานหลักที่โครงงานนี้ต้องการทดสอบ โครงงานนี้จึงมองส่วน Lead Scoring เป็นการทดสอบสมมุติฐานที่ว่า Unified Profile ซึ่งรวมข้อมูลจาก 3 แพลตฟอร์มจะให้ผลการพยากรณ์ที่ดีกว่าการใช้ข้อมูลจากแพลตฟอร์มเดียว ถ้าผลลัพธ์ออกมาสนับสนุนสมมุติฐานนี้ ก็จะเป็นหลักฐานเชิงประจักษ์ว่าขั้นตอน Identity Resolution ก่อนหน้าสร้างมูลค่าที่วัดได้จริง ไม่ใช่แค่การลดความซ้ำซ้อนของข้อมูลเท่านั้น\par
\section*{2.2 ทฤษฎีที่เกี่ยวข้อง}
\subsection*{2.2.1 Similarity Measures สำหรับ Record Matching}
การคำนวณความคล้ายคลึงระหว่างระเบียนข้อมูลเป็นหัวใจสำคัญของ Identity Resolution Similarity Measures ที่ใช้กันทั่วไปสามารถแบ่งออกเป็นสองกลุ่มหลัก ได้แก่ String-based Similarity เช่น Edit Distance (Levenshtein) ที่นับจำนวนการแก้ไขขั้นต่ำ และ Jaro-Winkler ที่เหมาะสำหรับชื่อบุคคล กับ Token-based Similarity เช่น Jaccard ที่วัด Overlap ของชุด Token และ TF-IDF Cosine ที่คำนึงถึงความสำคัญของคำในบริบท สูตรคำนวณหลักของแต่ละ Measure มีดังนี้ โดยสมการที่ 2.1 แสดง Edit Distance ที่ ใช้คำนวณจำนวนการแก้ไขขั้นต่ำระหว่างสองสตริง โดยที่ \ensuremath{\delta} = 0 หาก s I = t j มิฉะนั้น \ensuremath{\delta} = 1:\par
, j + 1, d i,j−1 + 1, d i−1,j−1 + \ensuremath{\delta} ) (2.1) โดยที่\par
d i,j = min( d i−1,j\par
\noindent • d i,j หมายถึง จำนวนการแก้ไขขั้นต่ำที่สุดในการแปลงข้อความ\par
\noindent • \ensuremath{\delta} (เดลตา) หมายถึง ค่าเงื่อนไข โดยมีค่าเป็น 0 หากอักขระตัวที่ ของสตริงแรก\par
( s i ) ตรงกับอักขระตัวที่ j ของสตริงที่สอง ( t j ) แต่มิฉะนั้น \ensuremath{\delta} จะมีค่าเป็น 1\par
, j + 1 หมายถึง ปฏิบัติการลบอักขระ (Deletion)\par
\noindent • d i−1,j\par
\noindent • d i,j−1 + 1 หมายถึง ปฏิบัติการแทรกอักขระ (Insertion)\par
\noindent • d i−1,j−1 + \ensuremath{\delta} หมายถึง ปฏิบัติการแทนที่อักขระ (Substitution)\par
|A\ensuremath{\cap}B|\par
|A\ensuremath{\cup}B| (2.2) โดยที่\par
J(A, B) =\par
% ---- source PDF page 23 ----
\clearpage
\noindent • J(A, B) หมายถึง ค่าความคล้ายคลึงแบบ Jaccard\par
\noindent • A, B หมายถึง ชุดของโทเค็น (Token) จากข้อความที่ต้องการเปรียบเทียบ\par
\noindent • \ensuremath{\mid}A \ensuremath{\cap}B \ensuremath{\mid} หมายถึง ขนาดของอินเตอร์เซกชัน (Intersection) หรือจำนวน\par
โทเค็นที่มีเหมือนกันระหว่างชุดข้อมูลทั้งสอง\par
\noindent • \ensuremath{\mid}A \ensuremath{\cup}B \ensuremath{\mid} หมายถึง ขนาดของยูเนียน (Union) หรือจำนวนโทเค็นทั้งหมด\par
เมื่อนำชุดข้อมูลทั้งสองมารวมกัน\par
A\ensuremath{\cdot}B \ensuremath{\Vert}A\ensuremath{\Vert}\ensuremath{\cdot}\ensuremath{\Vert}B\ensuremath{\Vert} (2.3) โดยที่\par
cos(A, B) =\par
\noindent • cos(A, B) หมายถึง ค่าความคล้ายคลึงโคไซน์ (การวัดมุมระหว่างเวกเตอร์)\par
\noindent • A, B หมายถึง เวกเตอร์ของข้อความ เช่น ค่า TF-IDF\par
\noindent • A \ensuremath{\cdot}B หมายถึง ผลคูณเชิงสเกลาร์ (Dot Product) ของเวกเตอร์ A และ B\par
\noindent • \ensuremath{\parallel}A \ensuremath{\parallel}, \ensuremath{\parallel}B \ensuremath{\parallel} หมายถึง ขนาด (Magnitude) ของเวกเตอร์ A และ B Halimi และ Ayday [3] ได้ทบทวนการนำ Similarity Measures ต่าง ๆ ไปใช้ในงาน User Identity Linkage พบว่าไม่มี Measure ใดที่ดีที่สุดในทุกสถานการณ์ แต่การรวม Measures หลายแบบ (Ensemble Approach) มักให้ประสิทธิภาพสูงกว่าการใช้ Measure เดียว โดยเฉพาะเมื่อข้อมูลมาจากแพลตฟอร์มที่มีรูปแบบข้อมูลแตกต่างกัน ตารางที่ 2.1 เปรียบเทียบ Similarity Measures ที่ใช้ในงาน Identity Resolution\par
Measure คำอธิบาย ข้อดี ข้อจำกัด\par
Edit Distance นับจำนวน Insert/Delete/Replace ขั้นต่ำ\par
ใช้ได้กับข้อความทุก รูปแบบ\par
ใช้ทรัพยากรสูง สำหรับข้อความ ยาว O(mn)\par
Jaro-Winkler วัดความคล้ายสตริงสั้น ให้น้ำหนักพิเศษกับอักขระ ต้น\par
ไม่เหมาะกับ ข้อความยาวหรือ หลายคำ Jaccard วัด Token Intersection / Union\par
เหมาะสำหรับชื่อบุคคล และรหัสสั้น\par
ไม่คำนึงถึงความถี่ และความสำคัญ ของคำ TF-IDF Cosine มุมระหว่างเวกเตอร์\par
คำนวณเร็ว ทนต่อการสลับ ลำดับคำ\par
คำนึงถึงความสำคัญของคำ ในบริบท Corpus\par
ต้องการ Corpus ขนาดใหญ่เพื่อ IDF ที่แม่นยำ\par
TF-IDF ของข้อความ\par
% ---- source PDF page 24 ----
\clearpage
\subsection*{2.2.2 Machine Learning สำหรับ Classification ใน Lead Scoring แบบจำลอง Machine Learning ที่นำมาใช้ใน Lead Scoring ส่วนใหญ่เป็นแบบจำลอง ประเภท Supervised Classification ที่ฝึกด้วยประวัติการแปลงของลูกค้าในอดีต Schafer และ คณะ [6] ได้เปรียบเทียบแบบจำลองหลายประเภทและพบว่า Gradient Boosting Machines (GBM) ให้ประสิทธิภาพสูงสุดด้วยค่า AUC-ROC 0.87 โดยมีข้อได้เปรียบสำคัญคือสามารถจัดการ Missing Values และ Imbalanced Class ได้ดีกว่า Logistic Regression ซึ่งพบบ่อยในชุดข้อมูล ลูกค้าจริง สูตรการอัปเดตของ GBM แสดงดังสมการที่ 2.4}
F m (x) = F m−1 (x) + \ensuremath{\eta} \ensuremath{\cdot}h m (x) (2.4) โดยที่\par
\noindent • F m (x) หมายถึง แบบจำลองสะสมในรอบที่ m\par
\noindent • F m−1 (x) หมายถึง แบบจำลองสะสมของรอบก่อนหน้า ( m −1 )\par
\noindent • \ensuremath{\eta} (อีตา) หมายถึง อัตราการเรียนรู้ (Learning Rate)\par
\noindent • h m (x) หมายถึง แบบจำลองการเรียนรู้แบบอ่อน (Weak Learner) เช่น\par
Decision Tree ที่ถูกฝึกด้วยค่าความคลาดเคลื่อน (Residual) ของรอบก่อน หน้า\par
\noindent • x หมายถึง ข้อมูลนำเข้า (Input Features) ของแบบจำลอง Feature Importance Analysis ในงานวิจัยของ Schafer และคณะ [6] ยังแสดงให้เห็นว่า Features ที่มาจากการรวมข้อมูลหลายแหล่ง (Merged Identity Features) มี P r e d i c t i v e P o w e r สูงกว่า F e a t u r e s จากแหล่งเดียว ซึ่งเป็นหลักฐานสำคัญที่สนับสนุนการออกแบบระบบ Identity Resolutionก่อนการสร้าง Lead Scoring Model ในโครงงานนี้\par
\subsection*{2.2.3 Active Learning และ Semi-supervised Learning}
Active Learning เป็นเทคนิคที่ให้ระบบสามารถเลือกตัวอย่างที่ต้องการ Label จาก มนุษย์ได้อย่างมีประสิทธิภาพ Li และคณะ [5] ได้นำ Active Learning มาใช้ใน Human-in-the- Loop Data Integration โดยระบบเลือกถามเฉพาะกรณีที่แบบจำลองมีความไม่แน่ใจสูงสุด (Uncertainty Sampling) กลยุทธ์นี้ช่วยลดจำนวนคำถามที่ต้องถามผู้ตรวจสอบได้อย่างมี นัยสำคัญในขณะที่ยังคงรักษาประสิทธิภาพของแบบจำลองในระดับสูง อย่างไรก็ตาม ในโครงงาน นี้ยังไม่ได้นำ Active Learning มาใช้จริง เนื่องจากข้อมูล Label ที่มีอยู่เพียงพอสำหรับการฝึก แบบ Supervised ทั่วไป และขอบเขตของโครงงานไม่ครอบคลุมถึงการพัฒนา Annotation Pipeline แบบ Online ทั้งนี้ แนวทางดังกล่าวถือเป็นแนวทางที่มีศักยภาพสำหรับการพัฒนาต่อ ในอนาคต โดยเฉพาะในกรณีที่ต้องการปรับแบบจำลองให้เข้ากับข้อมูลใหม่ที่เพิ่มขึ้นเรื่อย ๆ\par
% ---- source PDF page 25 ----
\clearpage
\section*{2.3 สถาปัตยกรรมระบบที่เกี่ยวข้อง}
\subsection*{2.3.1 LINKSOCIAL: ระบบเชื่อมโยงโปรไฟล์ข้ามแพลตฟอร์ม}
LINKSOCIAL [1] เป็นระบบที่พัฒนาโดย Sharma และ Dyreson สำหรับเชื่อมโยงโปรไฟล์ผู้ใช้งานข้ามแพลตฟอร์มสื่อสังคมออนไลน์หลายชนิด สถาปัตยกรรมของระบบประกอบด้วยสามส่วนหลัก ได้แก่ Feature Extraction ที่ดึงข้อมูลโปรไฟล์และพฤติกรรมจากแต่ละแพลตฟอร์ม, Similarity Computation ที่คำนวณความคล้ายคลึงระหว่างโปรไฟล์คู่ต่าง ๆ และ Classification ที่ตัดสินใจว่าโปรไฟล์คู่ใดเป็นบุคคลเดียวกัน ระบบ LINKSOCIAL [1] ได้รับการทดสอบบนชุดข้อมูลที่มีคู่บัญชีที่ยืนยันแล้วกว่า 7,000 คู่ ข้ามสามแพลตฟอร์ม ผลการทดลองแสดงค่า Precision สูงถึง 0.92 และ Recall ที่ 0.87 สถาปัตยกรรมนี้ให้แนวทางที่ชัดเจนสำหรับการออกแบบโมดูล Identity Resolution ของโครงงานนี้ โดยเฉพาะในด้านการจัดการ Feature จากหลายแพลตฟอร์มพร้อมกัน\par
\begin{center}\small รูปที่ 2.1 สถาปัตยกรรมระบบ LINKSOCIAL สำหรับการเชื่อมโยงโปรไฟล์ข้ามแพลตฟอลม์ [1]\end{center}
\subsection*{2.3.2 Human-in-the-Loop Data Integration Architecture Li และคณะ [5] ได้เสนอสถาปัตยกรรม Human-in-the-Loop Data Integration ที่ประกอบด้วยสี่ส่วน ได้แก่ ( 1 ) M a c h i n e P r o c e s s o r สำหรับกรณีที่ระบบมีความมั่นใจสูงและดำเนินการอัตโนมัติ (2) Human Processor ที่รับงานที่มีความไม่แน่ใจสูงและต้องการการตัดสินใจจากมนุษย์ (3) Task Manager ที่จัดสรรงานระหว่างสองส่วนแรกตามค่า Confidence Score และ (4) Data Integration System ที่รวบรวมผลลัพธ์เป็นฐานข้อมูลที่สมบูรณ์ Doan และคณะ [4] ได้เพิ่มเติมข้อค้นพบสำคัญเกี่ยวกับการออกแบบ Interface ที่มี ประสิทธิภาพสำหรับผู้ตรวจสอบโดยเน้นว่าควรแสดงข้อมูลที่เกี่ยวข้องที่สุดในรูปแบบที่เข้าใจง่าย พร้อมทั้งมีระบบ Audit Log ที่บันทึกการตัดสินใจทุกครั้งเพื่อความโปร่งใสและตรวจสอบได้ สถาปัตยกรรมนี้เป็นแบบอย่างในการออกแบบ Human Review Module ของโครงงานนี้}
% ---- source PDF page 26 ----
\clearpage
\begin{center}\small รูปที่ 2.2 สถาปัตยกรรม Human-in-the- Loop Data Integration [5]\end{center}
\section*{2.4 งานวิจัยที่เกี่ยวข้อง}
\subsection*{2.4.1 LINKSOCIAL: Linking User Profiles Across Multiple Social Media Platforms [1]}
งานวิจัยของ Sharma และ Dyreson [1] ซึ่งนำเสนอในการประชุม IEEE International C o n f e r e n c e o n B i g K n o w l e d g e ( I C B K ) ปี 2 0 1 8 ถือเป็นงานวิจัยหลักที่เกี่ยวข้องโดยตรงกับโครงงานนี้ ระบบ LINKSOCIAL ได้พัฒนาวิธีการดึง Feature จากโปรไฟล์ที่หลากหลาย ทั้งชื่อผู้ใช้ รูปโปรไฟล์ คำอธิบายตัวเอง ตำแหน่งที่ตั้ง และประวัติกิจกรรม จากนั้นใช้ Classifier ตัดสินว่าโปรไฟล์สองชุดเป็นบุคคลเดียวกันหรือไม่ จุดเด่นของงานวิจัยนี้ที่นำมาประยุกต์ใช้ในโครงงานได้แก่ กระบวนการ Feature E n g i n e e r i n g แบบ M u l t i - s o u r c e และการออกแบบ P i p e l i n e ที่รองรับข้อมูลจากหลายแพลตฟอร์มพร้อมกัน โครงงานนี้ขยายแนวคิดจาก Social Media Profile Linking ไปสู่ Customer Data Linking ในบริบทของ CRM Platform ที่มีระบบ Transaction, Behavioral และ Demographic Data\par
\subsection*{2.4.2 Connecting Users Across Social Media Sites: A Behavioral-Modeling Approach [2]}
งานวิจัยของ Zafarani และ Liu [2] ที่นำเสนอใน ACM SIGKDD ปี 2013 เป็นงานบุกเบิกที่แสดงให้เห็นว่าพฤติกรรมออนไลน์ของผู้ใช้มีความสม่ำเสมอข้ามแพลตฟอร์ม ซึ่งเปิดทางสู่การระบุตัวตนได้แม้ว่าชื่อผู้ใช้และข้อมูลโปรไฟล์จะแตกต่างกัน งานนี้สร้าง Behavioral Profile จากข้อมูลเชิงสถิติของการใช้งาน แล้วนำมาเปรียบเทียบโดยใช้ Machine Learning ข้อมูลเชิงลึกจากงานวิจัยนี้มีผลโดยตรงต่อการออกแบบ Feature Engineering ของโครงงาน โดยเฉพาะในส่วนของ Behavioral Signal เช่น เวลาในการซื้อ\par
% ---- source PDF page 27 ----
\clearpage
ความถี่ในการเข้าเว็บไซต์ และรูปแบบการเลือกสินค้า ซึ่งสามารถใช้เป็น Soft Signal สำหรับยืนยันการ Matching เมื่อ Hard Identifier เช่น Email หรือหมายเลขโทรศัพท์ ไม่สามารถใช้ได้\par
\subsection*{2.4.3 User Identity Linkage on Social Networks: A Review [3] งานวิจัยของ Halimi และ Ayday [3] ที่ตีพิมพ์ใน IEEE Access ปี 2024 เป็นการทบทวนวรรณกรรมที่ครอบคลุมที่สุดในสาขา User Identity Linkage ในช่วงไม่กี่ปีที่ผ่านมา งานนี้จำแนกเทคนิคออกเป็น 4 กลุ่ม ได้แก่ Profile-based, Behavioral- b a s e d , N e t w o r k - b a s e d และ H y b r i d A p p r o a c h พร้อมทั้งวิเคราะห์จุดแข็งและจุดอ่อนของแต่ละแนวทาง ผลการทบทวนของ Halimi และ Ayday [3] ระบุว่า Hybrid Approach ที่รวม Feature จากหลายมิติให้ประสิทธิภาพสูงที่สุดโดยเฉลี่ย แต่มีความซับซ้อนในการ Implement สูงกว่า ข้อค้นพบนี้สนับสนุนการตัดสินใจออกแบบของโครงงานนี้ที่เลือกใช้ Ensemble Approach ที่รวม Multiple Similarity Measures และ Behavioral Signals}
\subsection*{2.4.4 Human-in-the-Loop Challenges for Entity Matching [4] และ Data Integration [5]}
กลุ่มงานวิจัย Human-in-the-Loop ของ Doan และคณะ [4] และ Li และคณะ [5] ได้วางรากฐานสำคัญสำหรับการนำมนุษย์เข้ามาอยู่ในกระบวนการ Data Integration อย่างมีประสิทธิภาพ Doan และคณะ [4] เน้นที่ด้าน User Experience และ Workflow Design ในขณะที่ Li และคณะ [5] เน้นที่สถาปัตยกรรมระบบและการผสานรวมระหว่าง Machine Learning กับการตัดสินใจของมนุษย์ ทั้งสองงานวิจัยมีผลโดยตรงต่อการออกแบบโมดูล Human Review Interface และ Audit Log ของโครงงานนี้ โดยเฉพาะในแง่ของการออกแบบ Review Queue ที่เรียงลำดับตาม Confidence Score การแสดงข้อมูล Context ที่เพียงพอสำหรับการตัดสินใจ และการบันทึก Decision Log เพื่อใช้ในการ Retrain แบบจำลองในอนาคต\par
\subsection*{2.4.5 Predictive Lead Scoring Using Machine Learning Techniques [6] งานวิจัยของ Schafer, Hassan และ Muller [6] ที่ตีพิมพ์ใน Arabian Journal of Science and Engineering ปี 2022 เป็นงานวิจัยที่มีความเกี่ยวข้องโดยตรงกับโมดูล Lead Scoring ของโครงงานนี้ งานนี้ทดสอบแบบจำลอง Machine Learning บนชุดข้อมูล B2B Sales จริง ครอบคลุมทั้ง Feature Engineering, Model Selection, Evaluation Metrics และ Business Impact Assessment ข้อค้นพบสำคัญที่โครงงานนี้นำมาประยุกต์ใช้คือ การพิสูจน์ว่าการรวม Features จากหลายแหล่งข้อมูล (Identity Data Fusion) เพิ่มประสิทธิภาพของ Lead Scoring Model}
% ---- source PDF page 28 ----
\clearpage
อย่างมีนัยสำคัญ นอกจากนี้ยังได้นำ GBM มาเป็น Baseline Model เริ่มต้น และใช้ AUC-ROC เป็น Primary Evaluation Metric ตามที่งานวิจัยนี้แนะนำ\par
\begin{center}\small ตารางที่ 2.2 สรุปงานวิจัยที่เกี่ยวข้องกับโครงงาน\end{center}
อ้างอิง (ปี) หัวข้อ เทคนิคหลัก ความเกี่ยวข้องกับ\par
โครงงาน [1] Sharma \& Dyreson (2018)\par
สถาปัตยกรรม Multi-source Identity Resolution [2] Zafarani \& Liu (2013)\par
LINKSOCIAL Profile \& Behavioral Features, Classifier\par
Feature Engineering สำหรับ Behavioral Matching [3] Halimi \& Ayday (2024)\par
Behavioral- Modeling User Linking\par
Behavioral Profile, ML Classification\par
Framework การเลือก Similarity Measures [4] Doan et al. (2017)\par
User Identity Linkage Review\par
Hybrid Approach, Multi-dimensional Features\par
ออกแบบ Review Interface \& Audit Log [5] Li et al. (2017)\par
Human-in-the- Loop Entity Matching\par
Crowdsourcing, Workflow Design\par
สถาปัตยกรรม Hybrid Machine- Human Processing [6] Schafer et al. (2022)\par
Human-in-the- Loop Data Integration\par
Active Learning, Task Manager\par
Lead Scoring Model \& Feature Importance\par
Predictive Lead Scoring\par
GBM, AUC-ROC, Multi-source Features\par
จากการทบทวนงานวิจัยข้างต้น จะเห็นได้ว่างานส่วนใหญ่ถูกพัฒนาและทดสอบบนข้อมูล จากแพลตฟอร์มที่ยังให้บริการอยู่ ซึ่งสามารถดึง Behavioral Signal แบบ Real-time ได้ครบถ้วน สิ่งที่ยังขาดอยู่ในวรรณกรรมคือแนวทางที่รองรับกรณีที่หนึ่งในแพลตฟอร์มปิดตัวไปแล้วและมี เพียง Archive ที่ไม่สมบูรณ์เหลืออยู่ โครงงานนี้จึงพยายามเติมเต็มช่องว่างดังกล่าวโดยออกแบบ Feature Engineering และ Threshold Policy ที่ทำงานได้กับ Profile ที่ขาดข้อมูลบางส่วน\par
% ---- source PDF page 29 ----
\clearpage
และนำเสนอระบบ POSSIBLE\_MATCH เป็นกลไกจัดการความไม่แน่นอนที่เกิดขึ้นจากข้อจำกัด เชิงโครงสร้างของข้อมูล\par
\begin{center}\small รูปที่ 2.3 ขั้นตอนการทำงานของระบบ Identity Resolution และ Lead Scoring ที่พัฒนาในโครงงานนี้\end{center}
% ---- source PDF page 30 ----
\clearpage
% BEGIN CHAPTER 2 GA ADDITIONS
% ส่วนเพิ่มเติมบทที่ 2 ต่อจากหัวข้อ 2.4
\section*{2.5 กรอบแนวคิดของการให้คะแนนและการตัดสินใจเชื่อมโยงข้อมูล}
ทฤษฎีในส่วนเพิ่มเติมนี้ใช้รองรับการพัฒนานโยบายด้วย GA ที่อธิบายไว้ในบทที่ 3 โดยเชื่อมจากปัญหาการเชื่อมโยงอัตลักษณ์เดิมไปสู่การตัดสินใจภายใต้ข้อมูลไม่ครบและต้นทุนที่แตกต่างกัน เนื้อหาครอบคลุมบทบาทของการสร้างคู่ การให้คะแนน การตัดสินใจสามทาง การคำนวณต้นทุน การค้นหาพารามิเตอร์ และการประเมินผล เพื่อให้เห็นเหตุผลของการออกแบบแต่ละส่วนก่อนพิจารณาวิธีดำเนินงานจริง

\subsection*{2.5.1 ความสัมพันธ์ระหว่าง Candidate Generation, Pair Scoring และ Decision Policy}
กระบวนการเชื่อมโยงข้อมูลสามารถพิจารณาเป็นสามหน้าที่ต่อเนื่องกัน ขั้นแรกคือการเลือกคู่ที่มีโอกาสเกี่ยวข้องกันเพื่อส่งไปเปรียบเทียบ ขั้นที่สองคือการให้คะแนนแต่ละคู่จากคุณลักษณะที่มี และขั้นที่สามคือการกำหนดการดำเนินการจากคะแนนและหลักฐาน หน้าที่ทั้งสามเกี่ยวข้องกันแต่ตอบคำถามต่างกัน การสร้างคู่ตอบว่าควรเปรียบเทียบใครกับใคร การให้คะแนนตอบว่าคู่นั้นมีแนวโน้มเป็นบุคคลเดียวกันเพียงใด และนโยบายตอบว่าระบบควรทำอะไรต่อ

Papadakis และคณะ [9] จัดกลุ่มแนวทางเพิ่มประสิทธิภาพของ entity resolution เป็น blocking, filtering และแนวทางผสม โดย blocking ลดการเปรียบเทียบให้เหลือคู่ที่มีแนวโน้มตรงกัน ส่วน filtering ใช้เงื่อนไขเพื่อตัดคู่ที่ไม่น่าผ่านเกณฑ์ความคล้าย แนวคิดนี้รองรับการแยกขั้นคัดคู่จากขั้นตัดสินใจ และทำให้ต้องพิจารณาความครอบคลุมของคู่ควบคู่กับจำนวนคู่ที่ลดลง

สำหรับโครงงานนี้ ความสัมพันธ์ดังกล่าวมีผลต่อการอ่านผลโดยตรง หากคู่จริงไม่ผ่านเข้ามาในชุดผู้สมัคร นโยบายปลายทางไม่มีโอกาสตัดสินคู่นั้น แม้ GA จะเลือกพารามิเตอร์ที่เหมาะสมกับคู่ที่ได้รับก็ตาม ในทางกลับกัน หากคู่ถูกส่งเข้ามาแล้วแต่ตัดสินผิด ปัญหานั้นอยู่ในขอบเขตการให้คะแนนหรือการตัดสินใจ การแยกสองกรณีช่วยระบุว่าควรปรับปรุงระบบที่ขั้นใด

\subsection*{2.5.2 คะแนนความคล้าย คะแนนโมเดล และคำตัดสิน}
คะแนนความคล้ายเป็นค่าที่คำนวณจากข้อมูลเฉพาะด้าน เช่น ชื่อหรือข้อความ คะแนนสูงจึงสะท้อนความสอดคล้องตามวิธีวัดนั้น แต่ไม่ได้ยืนยันความเป็นบุคคลเดียวกันโดยลำพัง ชื่อที่เหมือนกันอาจเป็นชื่อทั่วไป ขณะที่ชื่อที่ต่างกันอาจเป็นชื่อย่อหรือนามแฝงของบุคคลเดียวกัน

คะแนนโมเดลเป็นผลที่ได้จากการพิจารณาคุณลักษณะหลายตัวร่วมกัน ส่วนคำตัดสินเป็นผลจากกฎที่นำคะแนนไปใช้ ความแตกต่างนี้สำคัญเพราะคะแนนเดิมสามารถนำไปสู่คำตัดสินต่างกันได้เมื่อเปลี่ยนเกณฑ์หรือต้นทุน การประเมินโมเดลจึงไม่เท่ากับการประเมินนโยบายที่ใช้งานโมเดลนั้น

ในนโยบายที่พัฒนาขึ้น คะแนนโมเดลเป็นหลักฐานหนึ่งที่นำมารวมกับข้อมูลด้านอื่น แม้คะแนนรวมสุดท้ายอยู่ระหว่าง 0 ถึง 1 ก็ยังเป็นคะแนนสำหรับตัดสินใจตามสูตรที่กำหนด การตีความเป็นโอกาสถูกต้องร้อยละใดต้องมีการตรวจสอบการปรับเทียบรองรับอีกขั้นหนึ่ง

\section*{2.6 การตัดสินใจสามทางและบทบาทของการตรวจสอบโดยมนุษย์}
\subsection*{2.6.1 แนวคิดการยอมรับ ปฏิเสธ และรอการตรวจสอบ}
การตัดสินใจสามทางเปิดทางให้ระบบแยกกรณีที่ยังไม่แน่ชัดออกจากกรณีที่ตัดสินอัตโนมัติได้ Ouyang และคณะ [10] ศึกษา record linkage โดยใช้ granular descriptors อธิบายบริเวณข้อมูลที่ไม่แน่นอน และส่งกรณีเหล่านั้นให้ตรวจสอบเพิ่มเติม แนวคิดนี้แสดงว่าการมีสถานะรอตรวจสอบเป็นองค์ประกอบของวิธีจัดการความไม่แน่นอนในงานเชื่อมโยงข้อมูล ไม่ใช่เพียงผลลัพธ์ที่ระบบประมวลผลไม่สำเร็จ

ในโครงงานนี้กำหนดชื่อผลลัพธ์เป็น MATCH, REVIEW และ NO\_MATCH โดยใช้เกณฑ์บนและเกณฑ์ล่างของคะแนน ถ้าให้ $z_i$ เป็นคะแนนที่ใช้ตัดสินคู่ที่ $i$ กฎพื้นฐานเขียนได้ดังนี้
\begin{equation*}
D_i=\begin{cases}
\mathrm{MATCH}, & z_i\geq t_m,\\
\mathrm{REVIEW}, & t_r\leq z_i<t_m,\\
\mathrm{NO\_MATCH}, & z_i<t_r,
\end{cases}\qquad 0\leq t_r\leq t_m\leq1.
\end{equation*}

เมื่อ $t_m$ และ $t_r$ ห่างกันมากขึ้น พื้นที่สำหรับ REVIEW จะกว้างขึ้น แต่จำนวนคู่ที่ย้ายสถานะจริงยังขึ้นอยู่กับการกระจายคะแนน หากไม่มีคู่ใดอยู่ในช่วงที่เปลี่ยนเกณฑ์ จำนวนคำตัดสินอาจไม่เปลี่ยนเลย การประเมินจึงต้องคำนวณกับข้อมูลจริง แทนการพิจารณาระยะห่างของเกณฑ์เพียงอย่างเดียว

\subsection*{2.6.2 REVIEW เป็นภาระงานที่ต้องมีปลายทางรองรับ}
สถานะ REVIEW หมายถึงระบบยังไม่ยืนยันการรวมและยังไม่ปฏิเสธคู่ในขั้นอัตโนมัติ หากนำไปใช้งาน เจ้าหน้าที่ต้องได้รับข้อมูลที่เพียงพอสำหรับตรวจสอบ เช่น โปรไฟล์ทั้งสองฝั่ง หลักฐานที่มี คะแนน และเหตุผลที่ถูกส่งเข้าคิว รวมถึงมีการบันทึกผลเมื่อพิจารณาเสร็จแล้ว

การส่ง REVIEW มากขึ้นอาจเปิดโอกาสให้คนช่วยแก้คู่ยาก แต่ไม่ได้แปลว่าคู่เหล่านั้นถูกแก้ไขถูกต้องแล้ว ในการทดลองที่ยังไม่มีผลตรวจจากเจ้าหน้าที่ คู่ REVIEW จึงยังเป็นงานค้าง และไม่ควรถูกรวมเป็นจำนวนคู่ที่ระบบจับคู่สำเร็จโดยอัตโนมัติ การแยกสถานะนี้ทำให้รายงานความสามารถของระบบกับภาระของผู้ปฏิบัติงานได้ตรงกัน

อีกประเด็นหนึ่งคือความสามารถในการตรวจสอบภายในเวลาที่กำหนด หากคิวมีคู่ 10,000 คู่และสมมติให้ใช้เวลาคู่ละ 1 นาที จะคิดเป็นเวลาทำงานประมาณ 166.7 ชั่วโมง ตัวอย่างนี้เป็นการคำนวณประกอบความเข้าใจ ไม่ใช่ผลที่วัดจากเจ้าหน้าที่ของโครงงาน และยังไม่รวมเวลาพัก การประสานงาน หรือกรณีที่ซับซ้อนกว่าปกติ

\subsection*{2.6.3 ความต่างระหว่าง Human Review และ Active Learning}
Human review เป็นขั้นตอนให้คนช่วยตัดสินหรือยืนยันข้อมูล ส่วน active learning เป็นวิธีเลือกตัวอย่างที่จะขอฉลากเพิ่มเติมเพื่อช่วยการเรียนรู้ของโมเดล สองแนวคิดอาจทำงานร่วมกันได้ แต่การมีคิว REVIEW เพียงอย่างเดียวไม่ได้แสดงว่าระบบทำ active learning แล้ว

สำหรับนโยบายที่ศึกษาในรอบนี้ การค้นหาใช้ฉลากและคะแนนที่เตรียมไว้ การนำผลตรวจของคนกลับไปเลือกตัวอย่างหรือฝึกโมเดลใหม่เป็นขั้นตอนต่อยอดที่ต้องออกแบบและประเมินแยกต่างหาก จึงต้องแยกสถานะงานที่ทำแล้วออกจากความสามารถที่คาดว่าจะเพิ่มในอนาคต

\section*{2.7 การตัดสินใจโดยคำนึงถึงต้นทุนความผิดพลาด}
\subsection*{2.7.1 เหตุผลที่ความผิดพลาดแต่ละประเภทมีน้ำหนักต่างกัน}
Elkan [12] อธิบายการเรียนรู้และการตัดสินใจเมื่อความผิดพลาดแต่ละประเภทมีบทลงโทษต่างกัน โดยชี้ให้เห็นความสำคัญของการกำหนดต้นทุนให้สอดคล้องกับความหมายทางเศรษฐศาสตร์ แนวคิดนี้ช่วยให้การเลือกคำตอบพิจารณาผลของการกระทำเพิ่มเติมจากโอกาสจำแนกถูกเพียงอย่างเดียว

ในการเชื่อมโยงโปรไฟล์ การรวมคนละคนกับการไม่รวมคนเดียวกันอาจใช้เวลาแก้ไขต่างกัน ขณะเดียวกัน การส่งให้คนตรวจสอบมีต้นทุนแม้สุดท้ายจะได้คำตอบถูกต้อง การออกแบบนโยบายของโครงงานจึงกำหนดน้ำหนักให้การกระทำทั้งสามประเภทอย่างเปิดเผย แทนการซ่อนสมมติฐานดังกล่าวไว้ในจุดตัดของคะแนน

\subsection*{2.7.2 ตารางต้นทุนที่ใช้ในโครงงาน}
กำหนดให้ $y=1$ หมายถึงคู่ที่เป็นบุคคลเดียวกัน และ $y=0$ หมายถึงคู่ที่เป็นคนละคน ตารางต้นทุนอย่างง่ายสำหรับนโยบายของโครงงานเป็นดังตารางที่ 2.3 โดยถือว่าการตัดสินอัตโนมัติที่ถูกต้องไม่มีต้นทุนความผิดพลาดในฟังก์ชันนี้

\begin{center}\small
\begin{tabularx}{\linewidth}{p{3.2cm}XX}
\toprule
คำตัดสิน & คู่จริงเป็นคนเดียวกัน & คู่จริงเป็นคนละคน \\
\midrule
MATCH & 0 & $c_{FP}$ \\
REVIEW & $c_R$ & $c_R$ \\
NO\_MATCH & $c_{FN}$ & 0 \\
\bottomrule
\end{tabularx}
\par ตารางที่ 2.3 โครงสร้างต้นทุนการตัดสินใจสามทางของโครงงาน
\end{center}

ต้นทุนรวมจึงเขียนได้เป็น $C=c_{FP}FP+c_{FN}FN+c_RR$ โดย $FP$ นับคู่ลบที่เป็น MATCH, $FN$ นับคู่บวกที่เป็น NO\_MATCH และ $R$ นับทุกคู่ที่เป็น REVIEW ตารางนี้เป็นแบบจำลองต้นทุนของการตัดสินใจขั้นปัจจุบัน ยังไม่ได้คำนวณความผิดพลาดที่เจ้าหน้าที่อาจทำภายหลังหรือผลกระทบทางธุรกิจทั้งหมด

เมื่อใช้สมมติฐานค่าแรง 100 บาทต่อชั่วโมง เวลาตรวจ 1 นาที เวลาแก้การรวมผิด 20 นาที และเวลาจัดการคู่ตกหล่น 10 นาที จะได้อัตราส่วนต้นทุน $c_{FP}:c_{FN}:c_R=2:1:0.1$ หลังปรับหน่วยให้ต้นทุนคู่ตกหล่นเท่ากับ 1 ดังนั้นฟังก์ชันที่ใช้ค้นหาในโครงงานคือ $C=2FP+FN+0.1R$ อัตราส่วนนี้มาจากสมมติฐานของโครงงาน ไม่ใช่ค่ามาตรฐานที่กำหนดโดยงานอ้างอิง [12]

\subsection*{2.7.3 ความสัมพันธ์ระหว่างความน่าจะเป็นกับการเลือกการกระทำ}
เพื่ออธิบายหลักการ สมมติว่า $p$ เป็นความน่าจะเป็นที่คู่หนึ่งเป็นบุคคลเดียวกันและสะท้อนข้อมูลได้เหมาะสม ภายใต้ตารางต้นทุนข้างต้น ต้นทุนคาดหมายของการกระทำแต่ละแบบเขียนได้เป็น
\begin{equation*}
L_{\mathrm{MATCH}}=c_{FP}(1-p),\qquad
L_{\mathrm{NO\_MATCH}}=c_{FN}p,\qquad
L_{\mathrm{REVIEW}}=c_R.
\end{equation*}

ตัวอย่างเมื่อ $p=0.60$ และใช้อัตราส่วน 2:1:0.1 จะได้ต้นทุนคาดหมายของ MATCH เท่ากับ 0.80 ของ NO\_MATCH เท่ากับ 0.60 และของ REVIEW เท่ากับ 0.10 การรอตรวจจึงมีต้นทุนคาดหมายต่ำที่สุดภายใต้สมมติฐานนี้ ตัวอย่างดังกล่าวแสดงเหตุผลว่าทำไมการเลือกการกระทำอาจต่างจากการเลือกคลาสที่มีโอกาสมากที่สุด

ในนโยบายแบบถ่วงน้ำหนักของโครงงาน คะแนนรวมไม่ได้ถูกยืนยันว่าเป็นความน่าจะเป็นที่ปรับเทียบแล้ว จึงไม่ได้แทนค่าคะแนนรวมนั้นลงในสูตรต้นทุนคาดหมายเพื่อสรุปเกณฑ์เชิงทฤษฎีโดยตรง แต่ใช้ GA ประเมินต้นทุนที่เกิดจากคำตัดสินบน validation วิธีอธิบายสองแบบนี้จึงต้องแยกกัน คือสูตรข้างต้นใช้ทำความเข้าใจหลักการ ส่วนวิธีในบทที่ 3 เป็นการค้นหาจากผลบนข้อมูล

\subsection*{2.7.4 ข้อควรพิจารณาเมื่อเปรียบเทียบต้นทุน}
ต้นทุนรวมเพิ่มตามจำนวนคู่ จึงเปรียบเทียบโดยตรงได้เมื่อใช้นโยบายต่างกันบนคู่ชุดเดียวกัน หากเปรียบเทียบคนละชุด ควรรายงานจำนวนคู่ สัดส่วนคู่บวก และต้นทุนเฉลี่ยต่อคู่ร่วมด้วย แม้ต้นทุนเฉลี่ยเท่ากัน สัดส่วนความผิดพลาดอาจต่างกันได้

การเปลี่ยนน้ำหนักต้นทุนยังอาจทำให้เลือกนโยบายคนละชุด หากต้องการศึกษาความไวต่อสมมติฐาน ต้องกำหนดสถานการณ์ต้นทุนใหม่และเลือกนโยบายบน validation อีกครั้ง ส่วนการนำพารามิเตอร์เดิมไปคำนวณต้นทุนใหม่ตอบได้เพียงว่านโยบายเดิมมีค่าใช้จ่ายเท่าใดภายใต้สมมติฐานใหม่ ไม่ได้ยืนยันว่าเป็นคำตอบที่เหมาะสมที่สุดของสถานการณ์นั้น

\section*{2.8 ขั้นตอนวิธีเชิงพันธุกรรมสำหรับค้นหานโยบาย}
\subsection*{2.8.1 แนวคิดและองค์ประกอบพื้นฐาน}
GA เป็นวิธีค้นหาที่ปรับปรุงคำตอบหลายชุดผ่านกระบวนการคัดเลือกและสร้างคำตอบรุ่นถัดไป แนวคิดพื้นฐานและการประยุกต์กับปัญหาการค้นหาสามารถศึกษาได้จาก Mitchell [18] สำหรับโครงงานนี้ คำตอบหนึ่งชุดแทนนโยบายหนึ่งนโยบาย และคุณภาพของคำตอบวัดด้วยต้นทุนบน validation รายละเอียดต่อไปนี้อธิบายความหมายขององค์ประกอบ GA ตามการออกแบบที่ใช้ในโครงงาน

โครโมโซม (Chromosome) คือชุดพารามิเตอร์ของนโยบาย ยีน (Gene) คือพารามิเตอร์แต่ละตำแหน่ง ประชากร (Population) คือกลุ่มนโยบายที่ประเมินในรอบเดียวกัน และรุ่น (Generation) คือรอบการประเมินและสร้างประชากรใหม่ คำว่า fitness ใช้เรียกค่าประเมินคุณภาพของคำตอบ ในงานนี้ใช้ต้นทุนเป็นเกณฑ์จัดอันดับโดยตรงและให้ค่าต่ำเป็นคำตอบที่ดีกว่า

\subsection*{2.8.2 การแทนคำตอบเป็นค่าจริงและการกำหนดพื้นที่ค้นหา}
พารามิเตอร์ของนโยบายประกอบด้วยเกณฑ์และน้ำหนักซึ่งเป็นค่าจริง จึงสามารถแทนโครโมโซมเป็นเวกเตอร์ตัวเลขได้โดยตรง นโยบายแบบถ่วงน้ำหนักมีโครงสร้างเป็น
\begin{equation*}
\boldsymbol{g}=[t_m,t_r,w_1,w_2,\ldots,w_8].
\end{equation*}

การแทนคำตอบต้องมาพร้อมเงื่อนไข เช่น เกณฑ์ REVIEW ไม่สูงกว่าเกณฑ์ MATCH น้ำหนักไม่ติดลบ และผลรวมน้ำหนักเท่ากับ 1 หากกลุ่มข้อมูลใดไม่มีหลักฐานพร้อมใช้ จะกำหนดน้ำหนักเป็นศูนย์ เงื่อนไขเหล่านี้จำกัดการค้นหาให้ได้คำตอบที่นำไปใช้กับกฎการตัดสินใจได้จริง

ตัวอย่างเช่น น้ำหนักดิบ 2, 3 และ 5 สามารถปรับเป็น 0.2, 0.3 และ 0.5 โดยหารด้วยผลรวม 10 ขั้นตอนนี้ทำให้ตีความน้ำหนักเป็นสัดส่วนได้ แต่ไม่ได้ทำให้ข้อมูลที่สัมพันธ์กันกลายเป็นหลักฐานอิสระ จึงยังต้องระวังการตีความว่าน้ำหนักแต่ละตัวบอกคุณค่าของข้อมูลด้านนั้นโดยแยกขาดจากด้านอื่น

\subsection*{2.8.3 การคัดเลือก การผสม และการกลายพันธุ์}
การคัดเลือกทำให้คำตอบที่มีต้นทุนต่ำมีบทบาทในการสร้างประชากรรุ่นถัดไป ในตัวค้นหาของโครงงาน ระบบเก็บกลุ่มอันดับต้นไว้ แล้วเลือกพ่อแม่จากกลุ่มนั้นเพื่อสร้างลูก การเก็บคำตอบที่ดีไว้โดยตรงเรียกว่า elitism และช่วยไม่ให้คำตอบที่ดีที่สุดที่พบในรอบก่อนหายไปทั้งหมดเมื่อสร้างประชากรใหม่

การผสมแบบ uniform crossover เลือกแต่ละยีนจากพ่อหรือแม่ด้วยโอกาสเท่ากัน ทำให้ลูกสามารถได้รับเกณฑ์จากคำตอบหนึ่งและน้ำหนักบางส่วนจากอีกคำตอบหนึ่ง หลังผสมแล้วอาจได้ค่าที่ต้องปรับให้อยู่ในเงื่อนไข เช่น เกณฑ์สลับลำดับหรือผลรวมน้ำหนักไม่เท่ากับ 1

การกลายพันธุ์เพิ่มการเปลี่ยนแปลงเล็กน้อยให้ยีนบางตัว เพื่อให้ประชากรสำรวจบริเวณที่ต่างจากพ่อแม่ ในโครงงานใช้ค่ารบกวนแบบปกติกับยีนที่ถูกเลือกให้กลายพันธุ์ แล้วตรวจขอบเขตและปรับน้ำหนักอีกครั้ง การกลายพันธุ์มากหรือน้อยเกินไปอาจทำให้ลักษณะการค้นหาต่างกัน จึงต้องบันทึกการตั้งค่าไว้ร่วมกับผลทดลอง

\subsection*{2.8.4 ความเหมาะสมกับฟังก์ชันต้นทุนของนโยบาย}
ต้นทุนของนโยบายคำนวณจากจำนวนคู่ที่ถูกตัดสินเป็นแต่ละสถานะ เมื่อขยับเกณฑ์เพียงเล็กน้อยแต่ยังไม่มีคู่ใดเปลี่ยนสถานะ ต้นทุนจะไม่เปลี่ยน แต่เมื่อเกณฑ์ข้ามคะแนนของคู่บางคู่ ต้นทุนจะกระโดดตามจำนวนและชนิดของคู่ที่เปลี่ยนคำตัดสิน ฟังก์ชันนี้จึงมีลักษณะเป็นช่วงคงที่และเปลี่ยนแบบไม่ต่อเนื่องในหลายบริเวณ

GA ประเมินคำตอบจากค่าต้นทุนได้โดยไม่ต้องหาอนุพันธ์ของฟังก์ชัน และค้นหาเกณฑ์กับน้ำหนักร่วมกันได้ จึงเป็นทางเลือกที่นำมาศึกษาในงานนี้ อย่างไรก็ตาม การค้นหาแบบกริด การสุ่ม หรือวิธีอื่นอาจเหมาะกับปัญหาบางรูปแบบได้เช่นกัน การใช้ GA ในโครงงานไม่ได้เป็นหลักฐานว่า GA มีประสิทธิภาพเหนือทุกวิธี หากยังไม่ได้เปรียบเทียบภายใต้งบประมาณการค้นหาที่เท่ากัน

\subsection*{2.8.5 การหยุดค้นหาและความแปรปรวนจากการสุ่ม}
โครงงานกำหนดจำนวนรุ่นเป็นเงื่อนไขหยุด แล้วเลือกคำตอบที่ดีที่สุดจากประชากรที่ประเมิน การใช้จำนวนรุ่นแน่นอนช่วยควบคุมขอบเขตการทดลอง แต่จำนวนรุ่นเท่ากันไม่ได้ทำให้เวลาใช้จริงเท่ากันเสมอ เพราะแต่ละนโยบายมีต้นทุนการคำนวณต่อคำตอบต่างกัน

ผล GA ยังขึ้นกับ seed ที่กำหนดประชากรเริ่มต้นและการสุ่มในแต่ละรุ่น การทำซ้ำหลาย seed จึงช่วยให้เห็นความเปลี่ยนแปลงของต้นทุนและพารามิเตอร์ หากหลายชุดพารามิเตอร์ให้ผลใกล้กัน ไม่ควรอธิบายน้ำหนักของคำตอบเพียงชุดเดียวว่าเป็นค่าที่จำเป็นตายตัวของปัญหา และการพบคำตอบเดิมหลายครั้งก็ยังไม่ใช่ข้อพิสูจน์ว่าพบคำตอบดีที่สุดทั่วพื้นที่ค้นหา

\section*{2.9 การรวมหลักฐานหลายด้านและการจัดการข้อมูลที่ขาดหาย}
\subsection*{2.9.1 ความคล้ายเชิงตัวอักษรและเชิงความหมาย}
ข้อมูลชื่อและข้อความอาจให้หลักฐานต่างลักษณะกัน ความคล้ายเชิงตัวอักษรช่วยจับรูปสะกดที่ใกล้กัน ส่วนความคล้ายเชิงความหมายช่วยพิจารณาข้อความที่ใช้คำต่างกันแต่สื่อความใกล้เคียง Reimers และ Gurevych [14] เสนอ Sentence-BERT เพื่อสร้างเวกเตอร์ประโยคที่นำมาเปรียบเทียบด้วย cosine similarity ได้ แนวคิดนี้รองรับการเพิ่มหลักฐานจากข้อความสั้น เช่น ประวัติย่อในโปรไฟล์

ในโครงงาน การที่ประวัติย่อมีความหมายคล้ายกันเป็นหลักฐานประกอบการจับคู่ เพราะคนละคนอาจเขียนถึงอาชีพ งานอดิเรก หรือสถานที่เดียวกันได้ ตัวอย่างเช่น โปรไฟล์สองคนระบุว่าเป็นนักออกแบบและชอบถ่ายภาพ ข้อความอาจคล้ายกันมากแต่ยังไม่เพียงพอสำหรับยืนยันตัวตน การใช้หลายด้านร่วมกันจึงต้องอ่านทั้งความสอดคล้องและข้อจำกัดของข้อมูลแต่ละประเภท

\subsection*{2.9.2 Missingness และ Availability Mask}
ค่าความคล้ายศูนย์อาจหมายถึงข้อมูลสองฝั่งไม่ตรงกัน แต่ช่องที่ไม่มีข้อมูลไม่สามารถให้ข้อสรุปเดียวกันได้ โครงงานจึงใช้ availability mask แยกว่าคะแนนด้านนั้นมีหลักฐานพร้อมใช้งานหรือไม่ โดยให้ $m_{ik}=1$ เมื่อช่องที่ $k$ ของคู่ที่ $i$ ใช้ได้ และให้เป็น 0 เมื่อใช้ไม่ได้

คำว่าใช้ได้ต้องมีนิยามตามชนิดข้อมูล เช่น ความคล้ายของประวัติย่ออาจต้องมีข้อความที่ใช้งานได้ทั้งสองฝั่ง ส่วนกลุ่มภาพอาจอาศัยภาพหรือคำบรรยายภาพตามที่ตารางคุณลักษณะรองรับ การใช้ mask จึงไม่ได้แก้ปัญหาข้อมูลหายทุกชนิดโดยอัตโนมัติ แต่ทำให้สูตรตัดสินใจรู้ว่าค่าใดควรถูกนำมาคิด

\subsection*{2.9.3 การถ่วงน้ำหนักตามหลักฐานที่มีในแต่ละคู่}
ให้ $s_{ik}$ เป็นคะแนนหลักฐาน และ $w_k$ เป็นน้ำหนักของนโยบาย โครงงานคำนวณคะแนนรวมเฉพาะช่องที่พร้อมใช้ด้วยสูตร
\begin{equation*}
S_i=\frac{\sum_{k=1}^{8}w_km_{ik}s_{ik}}{\sum_{k=1}^{8}w_km_{ik}}.
\end{equation*}

เมื่อไม่มีช่องที่มีน้ำหนักพร้อมใช้งาน ผลรวมในตัวหารเท่ากับศูนย์และระบบกำหนด REVIEW แทนการคำนวณคะแนน สูตรนี้เป็นข้อกำหนดของนโยบายที่พัฒนาขึ้นในโครงงาน ส่วนค่าของน้ำหนักและเกณฑ์ค้นหาจาก validation ตามวิธีในบทที่ 3

ตัวอย่างสมมติมีน้ำหนักคะแนนโมเดล 0.2 ชื่อ 0.5 และประวัติย่อ 0.3 คู่หนึ่งมีคะแนนโมเดล 0.9 คะแนนชื่อ 0.8 แต่ไม่มีประวัติย่อ คะแนนที่ใช้ตัดสินจะเท่ากับ $(0.2\times0.9+0.5\times0.8)/0.7$ หรือประมาณ 0.829 หากแทนประวัติย่อที่หายด้วยศูนย์และยังหารด้วย 1 จะเหลือ 0.58 ซึ่งลดลงจากการไม่มีข้อมูล ตัวอย่างนี้แสดงผลของการแยก mask ออกจากคะแนนอย่างเป็นรูปธรรม

แม้วิธีนี้ไม่ลงโทษข้อมูลขาดด้วยศูนย์โดยตรง แต่คะแนนอาจพึ่งหลักฐานเพียงไม่กี่ด้านเมื่อข้อมูลเหลือน้อย จึงต้องตรวจความผิดพลาดแยกตามความพร้อมใช้งานด้วย การปรับตัวหารเป็นการแก้ความหมายของคะแนน ไม่ใช่การรับประกันว่าคู่ที่มีข้อมูลน้อยจะตัดสินได้แม่นยำเท่าคู่ที่มีข้อมูลครบ

\subsection*{2.9.4 Coverage, Support และน้ำหนักที่มีผลจริง}
Coverage ของกลุ่มข้อมูลเป็นสัดส่วนคู่ที่มีหลักฐานกลุ่มนั้นพร้อมใช้งาน ส่วน positive support เป็นจำนวนคู่บวกที่มีหลักฐานดังกล่าว ทั้งสองค่าช่วยบอกขนาดข้อมูลที่รองรับการประเมินกลุ่มนั้น หลักฐานที่มีความแม่นยำสูงในตัวอย่างเพียงไม่กี่คู่ยังไม่เพียงพอสำหรับสรุปความน่าเชื่อถือของทั้งกลุ่ม

น้ำหนักที่มีผลจริงต่อคู่หนึ่งคำนวณได้เป็น
\begin{equation*}
\widetilde{w}_{ik}=\frac{w_km_{ik}}{\sum_jw_jm_{ij}}.
\end{equation*}

เมื่อข้อมูลครบ น้ำหนักที่มีผลจริงตรงกับน้ำหนักรวมที่ปรับผลรวมแล้ว แต่เมื่อบางช่องหาย สัดส่วนของช่องที่เหลือจะเพิ่มขึ้น การรายงานเฉพาะน้ำหนักรวมจึงยังไม่อธิบายการตัดสินใจของทุกคู่ และน้ำหนักเป็นศูนย์ต้องแยกว่าเกิดจากการปิดช่องเพราะไม่มีข้อมูล หรือเกิดจาก GA เลือกไม่ใช้ช่องที่มีข้อมูล

\subsection*{2.9.5 น้ำหนักนโยบายกับความสำคัญของคุณลักษณะ}
น้ำหนักนโยบายบอกสัดส่วนที่ใช้ในสูตรตัดสินใจ ส่วน feature importance อธิบายบทบาทของคุณลักษณะในแบบจำลองตามวิธีวิเคราะห์นั้น ทั้งสองอย่างอาจไม่ตรงกัน เพราะคะแนนโมเดลที่นำมาเป็นหนึ่งช่องอาจรวมผลของคุณลักษณะหลายตัวไว้แล้ว

หากคะแนนโมเดลและคะแนนชื่อให้ข้อมูลสัมพันธ์กัน GA อาจลดน้ำหนักช่องหนึ่งและเพิ่มอีกช่องหนึ่งโดยยังให้ผลตัดสินใกล้เดิม จึงไม่สามารถอ่านน้ำหนักว่าเป็นสัดส่วนสาเหตุของการเป็นบุคคลเดียวกัน หรือสรุปว่ากลุ่มที่น้ำหนักต่ำไม่มีประโยชน์ในทุกบริบท การอธิบายผลควรใช้ทั้งความครอบคลุม ความแปรปรวนของน้ำหนัก และการทดลองตัดกลุ่มข้อมูลออกตามขอบเขตที่มีหลักฐาน

\section*{2.10 การปรับเทียบคะแนนและการแยกข้อมูลเพื่อประเมินผล}
\subsection*{2.10.1 ความหมายของ Probability Calibration}
การปรับเทียบพิจารณาว่าค่าความน่าจะเป็นที่โมเดลให้สอดคล้องกับความถี่ของเหตุการณ์จริงหรือไม่ Guo และคณะ [13] ศึกษาประเด็นนี้ในโครงข่ายประสาทเทียมสมัยใหม่ และแสดงว่าความแม่นยำของการจำแนกกับคุณภาพการปรับเทียบเป็นคนละประเด็น งานดังกล่าวจึงรองรับการตรวจคะแนนก่อนนำไปตีความเป็นความมั่นใจ แต่ไม่ได้ยืนยันว่าเทคนิคเดียวกันให้ผลเหมือนกันกับข้อมูลโปรไฟล์ในโครงงาน

เอกสาร scikit-learn [15] อธิบายว่า หากกลุ่มตัวอย่างได้รับคะแนนใกล้ 0.8 และโมเดลปรับเทียบได้เหมาะสม สัดส่วนตัวอย่างบวกในกลุ่มนั้นควรใกล้ร้อยละ 80 การตรวจด้วย reliability diagram เปรียบเทียบค่าเฉลี่ยคะแนนกับสัดส่วนบวกในแต่ละช่วง และต้องดูจำนวนตัวอย่างประกอบ โดยเฉพาะช่วงที่มีข้อมูลน้อย

โครงงานเดิมมีขั้นปรับเทียบคะแนนก่อนตัดสินใจ แต่เมื่อเพิ่มสูตรคะแนนถ่วงน้ำหนัก คะแนนที่ได้เป็นตัวแปรใหม่ การปรับเทียบคะแนนโมเดลเดิมจึงไม่ถ่ายทอดมาเป็นคุณสมบัติของคะแนนรวมโดยอัตโนมัติ ประเด็นนี้เป็นเหตุผลให้ใช้คำว่า ``คะแนนนโยบาย'' แยกจาก ``ความน่าจะเป็นของโมเดล'' ในรายงาน

\subsection*{2.10.2 บทบาทของ Train, Calibration, Validation และ Test}
ข้อมูลแต่ละส่วนมีหน้าที่ต่างกัน Train ใช้เรียนรู้แบบจำลอง Calibration ใช้เรียนรู้การแปลงคะแนนเมื่อมีขั้นดังกล่าว Validation ใช้เลือกพารามิเตอร์หรือเปรียบเทียบนโยบาย และ Test ใช้ประเมินหลังการเลือกสิ้นสุด เอกสาร scikit-learn [15] ระบุว่าข้อมูลสำหรับเรียนรู้ตัวปรับเทียบควรเป็นอิสระจากข้อมูลที่ใช้ฝึกตัวจำแนก เพื่อหลีกเลี่ยงความมั่นใจที่เอนเอียงจากการประเมินข้อมูลที่โมเดลเคยเห็นแล้ว

การแบ่งบทบาทไม่ได้แปลว่าทุกรอบทดลองในโครงงานมีไฟล์สี่ชุดที่แยกกันเหมือนกันทั้งหมด วิธีที่ใช้จริงต้องตรวจจากกระบวนการสร้างคะแนนและข้อมูลแบ่งกลุ่มของรอบนั้น สำหรับการค้นหา GA รอบข้อมูลเต็ม ข้อกำหนดที่ใช้คือเลือกนโยบายจาก validation และบันทึกก่อนขั้นอ่าน test ส่วนความเป็นอิสระของคะแนนต้นทางต้องตรวจสอบย้อนกลับไปยังการฝึกโมเดลด้วย

\subsection*{2.10.3 การรั่วไหลของข้อมูลในงานที่หนึ่งบุคคลมีหลายโปรไฟล์}
การรั่วไหลเกิดขึ้นเมื่อข้อมูลที่ไม่ควรมีในช่วงฝึกหรือเลือกพารามิเตอร์เข้ามามีผลต่อการสร้างระบบ ทำให้ผลประเมินดูดีเกินกว่าการใช้งานกับข้อมูลใหม่ เอกสารแนวปฏิบัติ scikit-learn [16] เน้นการแยกข้อมูลก่อนเรียนรู้กระบวนการแปลงและไม่ใช้ test ในการ fit ขั้นตอนของระบบ

สำหรับโครงงานเชื่อมโยงตัวตน ต้องพิจารณาว่าคู่หลายคู่สามารถมีโปรไฟล์หรือบุคคลเดียวกันร่วมกันได้ การสุ่มแยกแถวคู่จึงอาจทำให้ข้อมูลของบุคคลเดียวกันอยู่ทั้งสองชุด หากต้องการทดสอบการทำงานกับบุคคลที่ไม่เคยเห็น ต้องแบ่งกลุ่มตามตัวตนและตรวจคู่ที่ข้ามกลุ่มด้วย ส่วนนี้เป็นเงื่อนไขของการตีความผล ไม่ใช่สิ่งที่ยืนยันได้จากชื่อคอลัมน์ train หรือ test เพียงอย่างเดียว

การใช้ผล test หลายครั้งเพื่อเลือกโครงสร้างนโยบายก็ทำให้ข้อมูลทดสอบมีอิทธิพลต่อการออกแบบ แม้จะไม่ได้ส่ง test เข้าไปในฟังก์ชัน GA โดยตรง การตรึงพารามิเตอร์ก่อนทดสอบหนึ่งรอบจึงเป็นหลักฐานของลำดับการเลือกในรอบนั้น ส่วนการอ้างว่า test ไม่เคยถูกใช้ในการพัฒนามาก่อนต้องมีประวัติการทดลองรองรับเพิ่มเติม

\section*{2.11 ตัวชี้วัดสำหรับนโยบายสามทาง}
\subsection*{2.11.1 การนับคำตัดสินให้ครอบคลุมสถานะ REVIEW}
การประเมินนโยบายสามทางต้องแยกผลจริงสองกลุ่มออกจากคำตัดสินสามสถานะ ตารางที่ 2.4 แสดงชื่อจำนวนที่ใช้ในส่วนเพิ่มเติมนี้ เพื่อไม่ให้สับสนระหว่างคู่บวกที่ปฏิเสธแล้วกับคู่บวกที่ยังรอตรวจสอบ

\begin{center}\small
\begin{tabularx}{\linewidth}{p{3.2cm}XX}
\toprule
ผลจริง & MATCH / REVIEW & NO\_MATCH \\
\midrule
เป็นคนเดียวกัน & $TP$ / $R_+$ & $FN_N$ \\
เป็นคนละคน & $FP$ / $R_-$ & $TN$ \\
\bottomrule
\end{tabularx}
\par ตารางที่ 2.4 สัญลักษณ์การนับผลจริงและคำตัดสินสามทาง
\end{center}

กำหนดให้ $R_+$ เป็นคู่บวกใน REVIEW, $R_-$ เป็นคู่ลบใน REVIEW และ $FN_N$ เป็นคู่บวกที่ถูกตัดสิน NO\_MATCH ดังนั้นจำนวนคู่บวกทั้งหมดในชุดประเมินคือ $TP+R_++FN_N$ และจำนวน REVIEW คือ $R_++R_-$ นิยามนี้เป็นการขยายการนับสำหรับนโยบายของโครงงานจากตัวชี้วัดการจำแนกพื้นฐาน [17]

\subsection*{2.11.2 Precision, Recall และ F1 ของการรวมอัตโนมัติ}
Precision ของการรวมอัตโนมัติพิจารณาว่าคู่ที่ระบบเลือกเป็น MATCH ถูกต้องเท่าใด ส่วน recall พิจารณาว่าคู่บวกทั้งหมดถูกนำไปเป็น MATCH แล้วเท่าใด คู่บวกที่อยู่ใน REVIEW ยังไม่ถือว่ารวมสำเร็จ สูตรที่ใช้จึงเป็น
\begin{equation*}
P=\frac{TP}{TP+FP},\qquad
R=\frac{TP}{TP+R_++FN_N},\qquad
F1=\frac{2PR}{P+R}.
\end{equation*}

ตัวอย่างมีคู่บวก 100 คู่ แบ่งเป็น MATCH 70 คู่ REVIEW 20 คู่ และ NO\_MATCH 10 คู่ จะได้ recall ของการรวมอัตโนมัติเท่ากับร้อยละ 70 ไม่ใช่ $70/(70+10)$ เพราะยังมีคู่บวกที่รอตรวจอยู่ 20 คู่ ขณะที่ในฟังก์ชันต้นทุน คู่ REVIEW ถูกคิดเป็นต้นทุนตรวจและไม่คิดซ้ำเป็น $FN_N$ การระบุนิยามทำให้ต้นทุนกับ recall ใช้ร่วมกันได้โดยไม่ขัดแย้ง

หากไม่มีคู่ MATCH ค่า precision ไม่มีตัวอย่างรองรับ หากไม่มีคู่บวก ค่า recall ก็ไม่มีตัวหารที่มีความหมาย โปรแกรมอาจเติมศูนย์เพื่อให้ตารางคำนวณต่อได้ แต่รายงานต้องบอกข้อจำกัดของกลุ่มนั้นด้วย

\subsection*{2.11.3 ภาระตรวจสอบและความครอบคลุมของระบบ}
Review rate คำนวณจากจำนวน REVIEW หารด้วยจำนวนคู่ทั้งหมดในชุดที่ประเมิน ส่วนอัตราการตัดสินอัตโนมัติคำนวณจากจำนวน MATCH รวมกับ NO\_MATCH หารด้วยจำนวนคู่ทั้งหมด ตัวชี้วัดนี้บอกสัดส่วนงานที่ระบบจัดการเอง แต่ต้องอ่านพร้อมความผิดพลาด เพราะระบบที่ตัดสินทุกคู่เองอาจมีอัตราอัตโนมัติสูงแม้คุณภาพต่ำ

ในระดับต้นทาง ให้ $N_+^{all}$ เป็นจำนวนคู่บวกในขอบเขตอ้างอิงทั้งหมด และ $N_+^{cand}$ เป็นคู่บวกที่เข้าสู่ชุดผู้สมัคร ความครอบคลุมคู่บวกของขั้นสร้างคู่คือ $N_+^{cand}/N_+^{all}$ หากระบบปลายทางตัดสินได้เฉพาะคู่ผู้สมัคร recall ของการรวมอัตโนมัติเมื่อวัดกับคู่บวกทั้งหมดจะไม่เกินสัดส่วนนี้ ภายใต้ขอบเขตอ้างอิงและนิยามคู่ที่ตรงกัน

การประเมินบน candidate pairs กับการประเมินตลอดกระบวนการจึงใช้ตัวหารต่างกัน การเพิ่มความแม่นยำของนโยบายในชุด candidate ไม่ได้แสดงว่าระบบแก้คู่ที่หลุดจากการสร้างคู่แล้ว ต้องรายงานทั้งสองระดับแยกกัน

\subsection*{2.11.4 การอ่านผลหลายตัวชี้วัดร่วมกัน}
AP และ ROC-AUC ใช้ศึกษาคุณภาพการจัดอันดับของคะแนนตามนิยามของแต่ละตัวชี้วัด [17] ส่วนจำนวน REVIEW และต้นทุนขึ้นกับนโยบายที่เลือกใช้คะแนน การเปลี่ยน threshold บนคะแนนเดิมจึงอาจเปลี่ยนต้นทุนและ F1 โดยไม่เปลี่ยนลำดับคะแนน หากเปลี่ยนเป็นคะแนนรวมแบบถ่วงน้ำหนัก ลำดับอาจเปลี่ยนด้วยและต้องประเมินใหม่

เมื่อรายงานผลควรระบุชื่อชุดข้อมูล จำนวนคู่ จำนวนคู่บวก และว่ารวม Exact-First หรือไม่ หากรวมคู่ Exact-First ผลที่ได้สะท้อนทั้งกฎเดิมและนโยบายที่พัฒนาใหม่ การแยกผลบนคู่ที่โมเดลให้คะแนนช่วยให้เห็นส่วนที่ GA มีอิทธิพลได้ชัดกว่า ส่วนผลรวมช่วยอธิบายผลลัพธ์ของกระบวนการที่นำไปใช้กับข้อมูลทั้งหมด

\section*{2.12 งานวิจัยเพิ่มเติมและความเกี่ยวข้องกับโครงงาน}
\subsection*{2.12.1 งานด้านการลดจำนวนคู่เปรียบเทียบ}
งานสำรวจของ Papadakis และคณะ [9] ให้กรอบจัดหมวดวิธีลดภาระการเปรียบเทียบใน entity resolution ความเกี่ยวข้องกับโครงงานคือการใช้กรอบนี้แยกคำถามด้านการคัดคู่ให้ครอบคลุมออกจากคำถามด้านการตัดสินคู่ที่คัดมาแล้ว งานดังกล่าวเป็นพื้นฐานเชิงระบบ ไม่ได้เป็นผลเปรียบเทียบโดยตรงกับ GA ที่ใช้ในโครงงาน

\subsection*{2.12.2 งานด้านการตัดสินใจสามทาง}
Ouyang และคณะ [10] เสนอวิธีใช้ granular descriptors เพื่อจำแนกบริเวณที่ไม่แน่นอนและส่งตรวจเพิ่มเติม จุดที่นำมาสนับสนุนงานนี้คือการให้สถานะไม่ตัดสินทันทีมีบทบาทชัดเจนใน record linkage ส่วนวิธีของโครงงานใช้คะแนนและเกณฑ์ที่ GA ค้นหาภายใต้ต้นทุนที่กำหนด จึงไม่ได้ทำซ้ำกลไก granular computing ของงานอ้างอิง

\subsection*{2.12.3 งานด้าน GA สำหรับ Entity Resolution}
Sun และคณะ [11] ศึกษาการใช้ GA เรียนรู้กฎจับคู่ที่ผสานการเปรียบเทียบหลายคุณลักษณะและเกณฑ์ที่เหมาะสม พร้อมใช้ active learning ลดภาระการติดฉลาก งานนี้แสดงตัวอย่างการประยุกต์ evolutionary search กับ entity resolution สำหรับโครงงานปัจจุบัน ความต่างอยู่ที่การค้นหานโยบายหลังได้คะแนนจากโมเดลที่ตรึงไว้ และยังไม่ได้ใช้วงจรขอฉลากใหม่แบบ active learning ในการค้นหานโยบายรอบนี้

\subsection*{2.12.4 งานด้านต้นทุน การปรับเทียบ และข้อความเชิงความหมาย}
Elkan [12] ให้พื้นฐานเรื่องการตัดสินใจเมื่อความผิดพลาดมีต้นทุนต่างกัน Guo และคณะ [13] ศึกษาคุณภาพของค่าความมั่นใจของโมเดล และ Reimers กับ Gurevych [14] พัฒนาเวกเตอร์ประโยคเพื่อเปรียบเทียบความหมาย งานทั้งสามเกี่ยวข้องกับคนละส่วนของระบบ ได้แก่ วัตถุประสงค์ของการตัดสินใจ ความหมายของคะแนน และการสร้างหลักฐานจากข้อความ

การนำแนวคิดเหล่านี้มาประกอบกันในโครงงานยังต้องตรวจสอบผลกับข้อมูลที่ใช้จริง งานอ้างอิงไม่ได้ให้ค่าต้นทุน 2:1:0.1 ไม่ได้ยืนยันผลการปรับเทียบของคะแนนรวม และไม่ได้รับรองว่าประวัติย่อที่มีความหมายคล้ายกันเป็นบุคคลเดียวกัน ข้อสรุปเหล่านี้ต้องมาจากการกำหนดขอบเขตและการทดลองของโครงงานเอง

\begin{center}\small
\begin{tabularx}{\linewidth}{p{2.8cm}XX}
\toprule
งานอ้างอิง & ประเด็นที่ใช้รองรับ & สิ่งที่โครงงานต้องศึกษาเอง \\
\midrule
Papadakis และคณะ [9] & การแยกขั้นลดคู่จากการเปรียบเทียบคู่ & ขอบเขตคู่ที่ GA เปลี่ยนคำตัดสินได้ \\
Ouyang และคณะ [10] & การจัดการความไม่แน่นอนด้วยการตัดสินใจสามทาง & เกณฑ์และภาระ REVIEW ของข้อมูลโครงงาน \\
Sun และคณะ [11] & การค้นหากฎเชื่อมโยงด้วย GA & นโยบายบนคะแนนที่ตรึงไว้และผลของการถ่วงน้ำหนัก \\
Elkan [12] & การพิจารณาต้นทุนความผิดพลาด & สมมติฐานต้นทุนและผลแลกเปลี่ยนของนโยบาย \\
Guo และคณะ [13] & การแยกความแม่นยำจากการปรับเทียบ & คุณภาพคะแนนของโมเดลและคะแนนนโยบายที่ใช้จริง \\
Reimers และ Gurevych [14] & หลักฐานความคล้ายเชิงความหมายของข้อความ & ประโยชน์ของข้อความต่อการระบุตัวตนในชุดข้อมูลนี้ \\
\bottomrule
\end{tabularx}
\par ตารางที่ 2.5 ความเชื่อมโยงระหว่างงานวิจัยเพิ่มเติมกับโครงงาน
\end{center}

\section*{2.13 การสังเคราะห์แนวคิดสู่กรอบการพัฒนาของโครงงาน}
จากทฤษฎีและงานที่ทบทวน การตัดสินใจสามทางและการใช้ GA ใน entity resolution มีพื้นฐานจากงานก่อนหน้าแล้ว จุดที่โครงงานศึกษาต่อคือวิธีประกอบคะแนนโมเดล หลักฐานที่มีไม่ครบ ต้นทุนที่กำหนด และกระบวนการเลือกนโยบายให้ทำงานร่วมกันในข้อมูลโปรไฟล์ที่เตรียมไว้ จึงวางข้อเสนอของงานเป็นการออกแบบและประเมินนโยบายในขอบเขตนี้ โดยไม่อ้างว่าเป็นผู้ริเริ่ม GA หรือการตัดสินใจสามทางในงานเชื่อมโยงข้อมูล

กรอบการพัฒนาประกอบด้วยข้อมูลนำเข้าเป็นคู่ผู้สมัครและคะแนนที่ตรึงไว้ ส่วนประมวลผลเป็นการรวมหลักฐานตามความพร้อมใช้งานและค้นหาเกณฑ์กับน้ำหนักบน validation ส่วนผลลัพธ์เป็นคำตัดสินสามทางพร้อมต้นทุนและตัวชี้วัด หลังเลือกเสร็จจึงตรึงนโยบายและประเมิน test ก่อนรายงานปริมาณงานบนคู่ทั้งหมด

ข้อเชื่อมโยงสำคัญมีสามประการ ประการแรก การคงคะแนนและชุดคู่เดิมช่วยให้แยกผลของนโยบายได้ ประการที่สอง การแยก mask ออกจากคะแนนทำให้ความหมายของข้อมูลขาดชัดเจน และประการที่สาม การแยกผล validation, test และข้อมูลรวมช่วยไม่ให้ผลเลือกพารามิเตอร์ถูกนำไปตีความเป็นผลกับข้อมูลใหม่ หลักการเหล่านี้นำไปสู่ขั้นตอนปฏิบัติและผลส่งออกที่อธิบายไว้ในหัวข้อ 3.6--3.8

% END CHAPTER 2 GA ADDITIONS
\clearpage
\chapter*{บทที่ 3}
\begin{center}\Large\bfseries วิธีการดําเนินวิจัย\end{center}
งานวิจัยนี้เลือกใช้กรอบ CRISP-DM (Cross-Industry Standard Process for Data\par
Mining) เป็นหลักในการพัฒนาระบบเชื่อมโยงอัตลักษณ์และการจัดลำดับลูกค้าเป้าหมายสำหรับ\par
ระบบบริหารจัดการความสัมพันธ์ลูกค้า (CRM) เนื่องจากครอบคลุมครบวงจรตั้งแต่การทำความ\par
เข้าใจธุรกิจไปจนถึงการนำไปใช้งานจริง ซึ่งตอบโจทย์ระบบที่มีความซับซ้อนทั้งส่วน Retrieval,\par
Classification และ Entity Consolidation\par
การประยุกต์ใช้เป็นลักษณะ Iterative (วนซ้ำ) ไม่ได้เป็นลำดับเส้นตรง ช่วยให้สามารถ ย้อนกลับไปปรับปรุงขั้นตอนก่อนหน้าได้ทันที โดยในบทนี้จะอธิบายรายละเอียดผ่าน 5 ขั้นตอนหลัก และเพิ่มการพัฒนานโยบายการตัดสินใจด้วย GA การประเมินนโยบาย และขั้นตอนดำเนินงานต่อยอดในหัวข้อ 3.6--3.8 ดังนี้\par
\noindent 1. การทำความเข้าใจธุรกิจ (Business Understanding)\par
\noindent 2. การทำความเข้าใจข้อมูล (Data Understanding)\par
\noindent 3. การเตรียมข้อมูล (Data Preparation)\par
\noindent 4. การสร้างแบบจำลอง (Modeling)\par
\noindent 5. การประเมินผล (Evaluation)\par
\begin{center}\small รูปที่ 3.1 แสดงกรอบวิธีวิจัยที่ใช้เป็นโครงหลักของการพัฒนาระบบเชื่อมโยงอัตลักษณ์ลูกค้าและ\end{center}
การให้คะแนนลูกค้าเป้าหมายในงานวิจัยนี้\par
% ---- source PDF page 31 ----
\clearpage
\section*{3.1 การทำความเข้าใจธุรกิจ (Business Understanding)}
\subsection*{3.1.1 ปัญหาทางธุรกิจ}
ปัญหาตั้งต้นของงานวิจัยนี้เกิดจากข้อจำกัดที่พบได้จริงในระบบบริหารจัดการ\par
ความสัมพันธ์ลูกค้า (CRM) คือ ข้อมูลของลูกค้าคนเดียวกันไม่ได้ถูกรวบรวมไว้ภายใต้ระเบียน\par
(Record) เดียว แต่กระจายอยู่หลายแหล่งและหลายแพลตฟอร์ม เช่น ทวิตเตอร์ (Twitter)\par
อินสตาแกรม (Instagram) และกูเกิลพลัส (Google+) โดยแต่ละแหล่งใช้โครงสร้างข้อมูลต่างกัน\par
มีความสมบูรณ์ของข้อมูลไม่เท่ากัน และผู้ใช้เองก็อาจใช้ชื่อ รูปแบบการเขียน หรือช่องทางการ\par
แสดงตัวตนต่างกันในแต่ละบัญชี ส่งผลให้บุคคลเดียวกันอาจถูกเก็บเป็นหลายระเบียนที่ไม่\par
เชื่อมโยงกันอย่างถูกต้อง เมื่อระบบไม่สามารถรวมระเบียนเหล่านี้ให้เป็นเอนทิตี (Entity) เดียวได้\par
ระบบ CRM จะเห็นลูกค้าเพียงบางส่วน ทำให้การวิเคราะห์พฤติกรรม การจัดกลุ่มลูกค้า การ\par
นำเสนอแบบเฉพาะบุคคล (Personalization) การจัดลำดับความสำคัญของลูกค้าเป้าหมาย และ\par
การวางแผนการสื่อสารทางการตลาดคลาดเคลื่อนได้\par
ในทางปฏิบัติ สิ่งที่ธุรกิจต้องการไม่ใช่เพียงการบอกว่าคู่โปรไฟล์ (Profile) ใดคล้ายกัน แต่\par
ต้องการระบบที่สามารถระบุได้ว่าระเบียนใดเป็นบุคคลเดียวกัน แล้วรวมข้อมูลเหล่านั้นให้\par
กลายเป็นข้อมูลลูกค้าแบบเบ็ดเสร็จ (Unified Customer Profile) ที่พร้อมใช้ในงานปลายทาง\par
เช่น การคัดเลือกลูกค้าเป้าหมาย การวางแผนการติดตามลูกค้า และการประเมินคุณค่าทางธุรกิจ\par
ของลูกค้าแต่ละราย\par
\subsection*{3.1.2 การแปลงเป็นปัญหาเชิงข้อมูล}
เมื่อแปลงโจทย์ทางธุรกิจให้อยู่ในรูปปัญหาเชิงข้อมูล ได้มีการกำหนดให้แกนหลักของ\par
ระบบเป็นงาน Identity Resolution โดยนิยามว่า ระบบต้องตัดสินว่าโปรไฟล์สองรายการเป็น\par
บุคคลเดียวกันหรือไม่ จากนั้นผลลัพธ์ที่ยืนยันแล้วจะถูกนำไปรวมเป็นเอนทิตีลูกค้าเดียวเพื่อใช้ใน\par
ขั้น Lead scoring ต่อไป อย่างไรก็ตาม งานนี้ไม่ใช่ปัญหา Classification เพียงอย่างเดียว เพราะ\par
ก่อนที่ระบบจะตัดสินว่าคู่ใด match ระบบต้องหา Candidate pairs ที่ควรพิจารณาให้พบก่อน\par
ข้อจำกัดสำคัญคือจำนวนคู่เปรียบเทียบทั้งหมดมีขนาดใหญ่มาก หากใช้โปรไฟล์ทั้งหมด\par
36,804 รายการ ซึ่งหากนับทุกคู่แบบไม่เรียงลำดับจะได้ 677,248,806 คู่ และหากพิจารณา\par
% ---- source PDF page 32 ----
\clearpage
เฉพาะคู่ข้ามแพลตฟอร์มตามนิยามของการทดลองหลัก จะยังมีมากถึง 449,149,239 คู่ ดังนั้น\par
การนำทุกคู่เข้าสู่แบบจำลองโดยตรงจึงไม่เหมาะสมทั้งในแง่เวลาและทรัพยากร\par
จากข้อจำกัดดังกล่าว จึงออกแบบกระบวนการแก้ปัญหาออกเป็น 3 ลำดับชั้น ที่ทำงาน สอดประสานกันอย่างเป็นระบบ ได้แก่\par
\noindent 1) ชั้นการดึงข้อมูลและคัดกรอง (Retrieval Layer) ทำหน้าที่กรอง Candidate Pairs\par
ออกจากพื้นที่ค้นหาขนาดใหญ่ เพื่อลดปริมาณข้อมูลขยะ โดยในสายการประมวลผลหลัก\par
ได้เลือกใช้กลยุทธ์ Exact-first และ Blocking เป็นกลไกสำคัญ\par
\noindent 2) ชั้นการให้คะแนน (Scoring Layer) ทำหน้าที่ประเมินคู่ Candidate ที่ผ่านการคัด\par
กรองจากชั้นแรก โดยจะคำนวณค่าความน่าจะเป็นในการเป็นบุคคลเดียวกัน ผ่านการ\par
คำนวณคุณลักษณะความคล้ายคลึง และการใช้ Multimodal Evidence\par
\noindent 3) ชั้นการตัดสินใจ (Decision Layer) ทำหน้าที่ประเมินผลลัพธ์จากคะแนนที่ได้ เพื่อฟัน\par
ธงว่าคู่โปรไฟล์ดังกล่าวควรจับคู่หรือไม่ และส่งต่อผลลัพธ์ไปยังกระบวนการควบรวม\par
ข้อมูลปลายทางต่อไป\par
\begin{center}\small รูปที่ 3.2 ซึ่งแสดงภาพรวมของระบบในเชิงธุรกิจและเชิงเทคนิคพร้อมกัน โดยให้\end{center}
% ---- source PDF page 33 ----
\clearpage
เห็นเส้นทางข้อมูลจากโปรไฟล์ต้นทาง การคัด candidate pairs การให้คะแนน การ ตัดสินใจ การรวมเอนทิตี และการสร้างคะแนนลูกค้า เพื่อให้ผู้อ่านเห็นตั้งแต่ต้นว่างานนี้ไม่ได้จบที่ การทำนายคู่ข้อมูล แต่ไปถึงการใช้งานจริงในระบบปลายทาง\par
\subsection*{3.1.3 ขอบเขตเชิงธุรกิจกับขอบเขตการทดลอง}
แม้เป้าหมายเชิงธุรกิจของระบบจะครอบคลุมการระบุว่าโปรไฟล์ใดเป็นบุคคลเดียวกันโดย\par
ไม่จำกัดแพลตฟอร์ม แต่การทดลองหลักของโครงการนี้มุ่งเน้นกรณี Cross-platform identity\par
linkage เป็นหลัก เนื่องจากข้อจำกัดของชุดข้อมูลกำกับที่มีอยู่จริง ในเชิงนิยามของปัญหา โปร\par
ไฟล์ทั้ง 36,804 รายการสามารถสร้างคู่แบบไม่เรียงลำดับได้ทั้งหมด 677,248,806 คู่ โดยแบ่งเป็น\par
คู่ข้ามแพลตฟอร์ม 449,149,239 คู่ และคู่ในแพลตฟอร์มเดียวกัน 228,099,567 คู่ อย่างไรก็ตาม\par
เมื่อใช้ user\_folder เป็นตัวแทนเอนทิตีต้นทางเพื่อสร้าง Ground truth จะพบว่าคู่เชิงบวกแบบ\par
Cross-platform มี 29,243 คู่ ขณะที่คู่เชิงบวกแบบ Same-platform มีเพียง 4 คู่เท่านั้น\par
ข้อจำกัดนี้ทำให้กรณี Same-platform ยังไม่มีจำนวนตัวอย่างเชิงบวกมากพอสำหรับการ\par
ฝึกและประเมินผลอย่างมั่นคงในระดับที่ใช้สรุปเชิงวิชาการได้ ด้วยเหตุนี้ เพื่อให้การสร้าง\par
positive pairs และการประเมินผลหลักยึดกับส่วนของข้อมูลที่มีหลักฐานรองรับเพียงพอที่สุด ผล\par
การทดลองหลักของงาน เช่น จำนวน Ground truth Cross-platform 29,243 คู่ Exact\par
Matches 12,403 คู่ Candidate pairs 2,073,842 คู่ และ retrieval coverage 88.67\% จึงควร\par
ถูกตีความว่าเป็นผลลัพธ์สำหรับกรณี Cross-platform เป็นสำคัญ ไม่ใช่ผลยืนยันประสิทธิภาพ\par
ของ All-platform identity resolution\par
ระบบมีขอบเขตเชิงธุรกิจที่กว้างกว่า Implementation เชิงทดลองในรายงานฉบับนี้ โดย\par
ขอบเขตธุรกิจมุ่งสู่การรวม Identity ในทุกกรณี ขณะที่ขอบเขตการทดลองหลักถูกจำกัดให้เหมาะ\par
กับข้อเท็จจริงของข้อมูลที่มีอยู่\par
\subsection*{3.1.4 เป้าหมายของระบบ}
จากการวิเคราะห์ข้างต้น เป้าหมายของระบบในงานนี้จึงไม่ได้จำกัดอยู่ที่การเพิ่มค่าความ แม่นยำของแบบจำลองอย่างเดียว แต่ครอบคลุมเป้าหมายเชิงระบบ 4 ประการพร้อมกัน\par
\noindent 1) ระบบต้องค้นหาคู่ที่น่าจะเป็นบุคคลเดียวกันได้ครอบคลุมเพียงพอ กล่าวคือ\par
retrieval ต้องไม่ทำให้คู่จริงหลุดจาก pipeline มากเกินไป\par
% ---- source PDF page 34 ----
\clearpage
\noindent 2) ระบบต้องลดจำนวนคู่เปรียบเทียบลงจนสามารถคำนวณได้จริงในระดับ production\par
โดยไม่ต้องสร้าง feature และ score ทุกคู่ในจักรวาลข้อมูล\par
\noindent 3) ระบบต้องให้คะแนนที่สามารถนำไปกำหนด Threshold สำหรับการใช้งานจริงได้\par
เพื่อให้สามารถแบ่งผลลัพธ์เป็น MATCH, REVIEW และ NO\_MATCH อย่างเป็น\par
ระบบ\par
\noindent 4) ผลลัพธ์สุดท้ายต้องสามารถแปลงเป็น Unified customer profiles และ Lead\par
tiers ที่พร้อมใช้งานในระบบบริหารจัดการความสัมพันธ์ลูกค้า (CRM) ได้โดยตรง\par
ดังนั้น ขั้น Business Understanding เป็นขั้นที่ทำให้ตัดสินใจว่า Pipeline จำเป็นต้อง\par
ประกอบด้วย Retrieval, Scoring, Decision tiers, Review workflow, Entity merge และ\par
Lead scoring อยู่ในระบบเดียวกัน ไม่ใช่การพัฒนา classifier แยกเดี่ยวเพียงตัวเดียวแล้วจบ\par
กระบวนการ\par
\section*{3.2 การทำความเข้าใจข้อมูล (Data Understanding)}
ข้อมูลที่ใช้ในงานนี้อ้างอิงจากชุดข้อมูล LinkSocial ซึ่งเป็นชุดข้อมูลมาตรฐานสำหรับงาน เชื่อมโยงอัตลักษณ์ผู้ใช้ข้ามแพลตฟอร์ม ข้อมูลต้นทางถูกจัดเก็บในรูปแบบไฟล์ JSON ของโปรไฟล์ สาธารณะจากหลายแพลตฟอร์มสังคมออนไลน์ และในงานวิจัยนี้เลือกใช้ 3 แพลตฟอร์มหลัก ได้แก่ Twitter, Google+ และ Instagram ในขั้นตอนการเตรียมข้อมูล เริ่มจากไล่อ่าน โฟลเดอร์ 1.profile.data 2.profile.datและ 3.profile.data แล้วรวมโปรไฟล์ทั้งหมดเป็นตาราง เดียว โดยดึง user\_folder, source\_folderและ platform มาจากโครงสร้าง Directory และชื่อ ไฟล์ต้นทาง การรวมข้อมูลในขั้นนี้ทำให้ข้อมูลกึ่งโครงสร้างหลายแหล่งถูกเปลี่ยนเป็นรูปแบบตาราง ที่พร้อมตรวจสอบและประมวลผลต่อได้เพื่อให้สะดวกต่อการทำความสะอาด และนำไปใช้ใน ขั้นตอนถัดไป หลังจากนั้นจึงเพิ่มข้อมูลกำกับที่จำเป็น เพื่อให้สามารถติดตามที่มาของข้อมูลและ เชื่อมโยงข้อมูลระหว่างขั้นตอนต่าง ๆ ของการประมวลผลได้อย่างถูกต้อง\par
เมื่อผ่านการเตรียมข้อมูลเบื้องต้นแล้ว ข้อมูลจะถูกทำความสะอาดและปรับให้อยู่ใน\par
รูปแบบมาตรฐานสำหรับใช้ในสถาปัตยกรรมการประมวลผลหลักของงานวิจัย โดยชุดข้อมูล\par
สุดท้ายที่ใช้ประกอบด้วยข้อมูลทั้งหมด ดังแสดงในตารางที่ 3.1\par
% ---- source PDF page 35 ----
\clearpage
\begin{center}\small ตารางที่ 3.1 จำนวนข้อมูลแยกตามแพลตฟอร์มในชุดข้อมูลมาตรฐานหลัก\end{center}
Platform Number of Profiles\par
Twitter 13,960\par
Google+ 11,890\par
Instagram 10,957\par
รวม 36,807\par
เมื่อพิจารณาโครงสร้างของข้อมูลที่รวมได้แล้ว จะพบว่าตารางกลางมีคอลัมน์สำคัญ ทั้งหมด 11 คอลัมน์ ได้แก่ userName, fullName, bio, location, externalUrl, pictureURL, user\_folder, outputProfileName, platform, source\_folder และ bigrams อย่างไรก็ตาม ไม่ใช่ทุกคอลัมน์จะถูกนำเข้าสู่แบบจำลองโดยตรง บางคอลัมน์ทำหน้าที่เป็นข้อมูลดิบสำหรับสร้าง คุณลักษณะ บางคอลัมน์ทำหน้าที่เป็นคีย์อ้างอิงของระบบ และบางคอลัมน์เป็นข้อมูลทางผ่านที่ใช้ เฉพาะในช่วง reprocessing เท่านั้น\par
\begin{center}\small ตารางที่ 3.2 ฟิลด์อ้างอิงหลักและบทบาทในการใช้งานจริงของระบบ\end{center}
กลุ่มข้อมูล คอลัมน์\par
คำอธิบายโดยย่อ บทบาทหลักในระบบ\par
อ้างอิง\par
Core Identity Attributes\par
userName ชื่อบัญชีผู้ใช้หรือ\par
ใช้ใน exact matching, prefix blocking\par
handle\par
และ string similarity\par
fullName ชื่อที่แสดงบนโปร\par
ใช้เปรียบเทียบชื่อแสดงผลและเสริมการจับคู่\par
ไฟล์\par
externalUrl ลิงก์ภายนอกที่\par
ใช้ใน URL/domain normalization,\par
ผู้ใช้ระบุ\par
exact-first และ domain overlap\par
matching\par
% ---- source PDF page 36 ----
\clearpage
\begin{center}\small ตารางที่ 3.2 (ต่อ) ฟิลด์อ้างอิงหลักและบทบาทในการใช้งานจริงของระบบ\end{center}
กลุ่มข้อมูล คอลัมน์อ้างอิง คำอธิบายโดยย่อ บทบาทหลักในระบบ\par
Contextual \& Multimodal Signals\par
bio คำอธิบายตัวตนหรือ\par
ใช้สร้าง semantic และ\par
ข้อความโปรไฟล์\par
text features\par
location สถานที่ที่ผู้ใช้ระบุ ใช้ใน location\par
normalization และ\par
location similarity\par
pictureURL URL ของรูปโปรไฟล์\par
ใช้เป็นจุดตั้งต้นของสายภาพ\par
ต้นทาง\par
และ multimodal\par
features\par
Linkage Keys \& ERD Identifiers\par
profile\_id รหัสอ้างอิงโปรไฟล์\par
ใช้เป็นคีย์อ้างอิงมาตรฐาน\par
หลัง normalize\par
ของโปรไฟล์ใน pipeline\par
profile\_row\_id รหัสลำดับแถวของ\par
ใช้เชื่อม feature matrix,\par
โปรไฟล์\par
scoring outputs และ\par
CRM artifacts\par
platform ชื่อแพลตฟอร์มของ\par
ใช้กำหนดขอบเขต cross-\par
โปรไฟล์\par
platform pairing และ\par
สร้าง platform\_pair\_code\par
Validation \& Audit Trail\par
user\_folder ตัวระบุเอนทิตีจาก\par
ใช้สร้าง ground truth และ\par
แหล่งต้นทาง\par
positive pairs\par
source\_folder ตัวระบุแหล่งที่มาของ\par
ใช้ตรวจสอบ provenance\par
ข้อมูล\par
และ audit trail\par
Processing Artifacts\par
outputProfileName ตัวระบุชั่วคราวในขั้น\par
ใช้เป็นฐานสำหรับการ\par
preprocessing\par
สร้าง profile\_id\par
% ---- source PDF page 37 ----
\clearpage
\begin{center}\small ตารางที่ 3.2 (ต่อ) ฟิลด์อ้างอิงหลักและบทบาทในการใช้งานจริงของระบบ\end{center}
กลุ่มข้อมูล คอลัมน์อ้างอิง คำอธิบายโดยย่อ บทบาทหลักในระบบ\par
Processing Artifacts\par
bigrams ข้อมูลตัวอักษรคู่ที่\par
ไม่ได้อยู่ใน selected\par
สร้างในสาย\par
features ของแบบจำลอง\par
preprocessing รุ่น\par
หลัก\par
ต้น\par
จากตารางที่ 3.2 จะเห็นได้ว่าฟิลด์สำคัญของระบบไม่ได้มีบทบาทเหมือนกันทั้งหมด ฟิลด์\par
กลุ่ม Core Identity Attributes เช่น userName, fullName และ externalUrl เป็นแกนหลัก\par
ของการพิสูจน์อัตลักษณ์ ขณะที่ bio, location และ pictureURL ทำหน้าที่เป็นข้อมูลบริบทและ\par
ข้อมูลเสริมที่ช่วยเพิ่มความแม่นยำในกรณีที่สัญญาณหลักไม่เพียงพอ ส่วน profile\_id,\par
profile\_row\_id, user\_folder และ source\_folder เป็นฟิลด์เชิงโครงสร้างที่ทำให้ระบบ\par
สามารถเชื่อมโยงข้อมูล ประเมินผล และตรวจสอบย้อนกลับได้อย่างเป็นระบบ\par
เมื่อพิจารณาความพร้อมใช้งานของข้อมูล พบว่าฟิลด์อย่าง userName, fullName และ externalUrl มีข้อมูลค่อนข้างสมบูรณ์ จึงเหมาะสำหรับใช้เป็นคุณลักษณะหลัก ในขณะที่ location และ pictureURL มีข้อมูลขาดหายมากกว่า จึงเหมาะสำหรับใช้เป็นข้อมูลสนับสนุน มากกว่าการใช้คุณลักษณะหลักเพียงลำพัง เพื่อให้เห็นภาพรวมของความสมบูรณ์ของข้อมูลได้ ชัดเจนยิ่งขึ้น จึงแสดงจำนวนข้อมูลที่ไม่เป็นค่าว่างของแต่ละฟิลด์ในรูปแบบกราฟแท่ง ดังรูปที่ 3.3\par
\begin{center}\small รูปที่ 3.3 จำนวนข้อมูลที่ไม่เป็นค่าว่างของฟิลด์สำคัญในชุดข้อมูล LinkSocial\end{center}
% ---- source PDF page 38 ----
\clearpage
ในแง่ของการจัดการข้อมูล ขั้นตอนนี้ไม่ได้มีเพียงการรวมข้อมูลหรือตรวจสอบจำนวน\par
คอลัมน์เท่านั้น แต่เป็นการเตรียมข้อมูลให้พร้อมสำหรับ Pipeline ในระยะถัดไป ข้อมูลดิบจาก\par
หลายแพลตฟอร์มซึ่งอยู่ในรูปแบบกึ่งโครงสร้างจะถูกแปลงให้อยู่ในรูปแบบมาตรฐาน เพื่อให้พร้อม\par
นำไปใช้งานต่อในขั้นตอน Pair Building, Feature Engineering และ Model Training\par
\begin{center}\small รูปที่ 3.4 แสดงการไหลของข้อมูลจากแหล่งข้อมูลดิบไปสู่ชุดข้อมูลมาตรฐาน\end{center}
\begin{center}\small รูปที่ 3.5 เพื่อแสดงความสัมพันธ์เชิงคีย์ของข้อมูลในลักษณะ Entity-Relationship\end{center}
Diagram\par
% ---- source PDF page 39 ----
\clearpage
\section*{3.3 การเตรียมข้อมูล (Data Preparation)}
ขั้นตอนการเตรียมข้อมูลนับเป็นหนึ่งในกระบวนการที่มีความสำคัญอย่างยิ่ง เนื่องจากเป็น ขั้นตอนที่เปลี่ยนผ่านข้อมูลดิบที่มีความหลากหลายและไม่สม่ำเสมอสูง ให้กลายเป็นชุดข้อมูล มาตรฐานที่พร้อมใช้งานสำหรับกระบวนการสร้างชุดคู่ข้อมูล (Pair Building) การทำวิศวกรรม คุณลักษณะ (Feature Engineering) และการฝึกสอนตัวแบบ (Model Training) ทั้งนี้ การ เตรียมข้อมูลในโครงงานนี้ไม่ได้จำกัดอยู่เพียงการจัดการค่าว่างหรือการแปลงข้อความขั้นพื้นฐาน เท่านั้น แต่ยังครอบคลุมถึงกระบวนการเชิงลึกอย่างครบวงจร ตั้งแต่การทำความสะอาดข้อมูลเชิง ข้อความ การสร้างคีย์อ้างอิงกลาง การปรับมาตรฐาน (Normalize) คุณลักษณะสำคัญ เช่น ข้อมูล สถานที่ การจัดเตรียมฐานข้อมูลภาพ ไปจนถึงการสร้างชุดคู่ข้อมูลสำหรับการฝึกสอน และการ สร้างเมทริกซ์คุณลักษณะ (Feature Matrix) แบบแบ่งส่วน (Chunked) เพื่อรองรับการแบ่งชุด ข้อมูลฝึกสอน (Train), ชุดข้อมูลตรวจสอบ (Validation) และชุดข้อมูลทดสอบ (Test) อย่างมี ประสิทธิภาพ\par
\subsection*{3.3.1 การทำความสะอาดข้อมูลและสร้างรหัสอ้างอิง (Data Cleaning \& Canonical key) กระบวนการนี้มุ่งเน้นไปที่การเตรียมความพร้อมของข้อมูลดิบให้มีความถูกต้องและเป็น มาตรฐานเดียวกัน โดยมีขั้นตอนหลักดังต่อไปนี้}
\noindent 1) การจัดการค่าที่ขาดหาย (Missing Value Handling)\par
\noindent • กลุ่มข้อมูลระบุตัวตนหลัก ได้แก่ userName, fullName, bio, location,\par
externalUrl และ pictureURL จะถูกแทนที่ค่าที่ขาดหายด้วยค่าว่าง เพื่อรักษา\par
สะพานการเชื่อมต่อข้อมูล ให้ระบบสามารถประมวลผลในขั้นต่อไปได้อย่างสม่ำเสมอ\par
โดยไม่เกิดข้อผิดพลาด\par
\noindent • กลุ่มข้อมูลเพื่อการเชื่อมโยง ในกรณีที่ outputProfileName ไม่มีค่า ระบบจะดึง\par
ข้อมูลจาก user\_folder มาเติมทดแทน เพื่อให้มั่นใจว่าทุกโปรไฟล์มีตัวอ้างอิงตั้งต้น\par
ที่สามารถใช้งานได้ทันที\par
\noindent 2) การปรับมาตรฐานข้อความ (Text Normalization)\par
ข้อมูลในแต่ละฟิลด์จะถูกทำความสะอาดด้วยชุดกฎที่สอดคล้องกับธรรมชาติของข้อมูล นั้น ๆ เพื่อเพิ่มความแม่นยำในการเปรียบเทียบ ดังนี้\par
% ---- source PDF page 40 ----
\clearpage
\noindent • userName ตัดสัญลักษณ์ @ นำหน้า แปลงเป็นตัวพิมพ์เล็กทั้งหมด ลบอีโมจิและ\par
สัญลักษณ์ที่ไม่จำเป็น รวมถึงยุบรวมตัวคั่นที่พิมพ์ซ้ำซ้อน (เช่น ... หรือ \_\_\_) ให้เหลือ\par
เพียงตัวเดียว\par
\noindent • fullName แปลงเป็นตัวพิมพ์เล็ก ตัดคำนำหน้าชื่อที่ระบบไม่ต้องการ ลบอักขระ\par
พิเศษ และจัดระเบียบช่องว่างให้เป็นมาตรฐาน เพื่อช่วยให้ชื่อที่สะกดต่างรูปแบบกัน\par
แต่มีความหมายเดียวกันสามารถนำมาเปรียบเทียบได้ง่ายขึ้น\par
\noindent • bio เนื่องจากเป็นฟิลด์ที่มีข้อมูลหลากหลายรูปแบบซ่อนอยู่ กระบวนการจึง\par
ครอบคลุมทั้งการลบข้อมูลรบกวนและการสกัดข้อมูลเฉพาะ โดยจะดึง URL การ\par
กล่าวถึง (@mention) และแฮชแท็ก (\#---) แยกออกมา เพื่อนำไปสร้างเป็น\par
คุณลักษณะเชิงบริบทในขั้นตอนถัดไป\par
\noindent 3) การสร้างคีย์อ้างอิงกลาง (Canonical Key Generation) เพื่อให้ระบบสามารถ ตรวจสอบย้อนกลับ ได้ตั้งแต่ต้นน้ำจนถึงปลายน้ำ ได้ออกแบบระบบคีย์อ้างอิงแบบ 2 ระดับ ได้แก่\par
\noindent • profile\_id (Entity Key) รหัสหลักที่ได้จากการปรับมาตรฐานข้อมูล ทำหน้าที่เป็น\par
ศูนย์กลางในการรวบรวมและเชื่อมโยงโปรไฟล์ทั้งหมดของบุคคลเดียวกันเข้าด้วยกัน\par
\noindent • profile\_row\_id (Row Key) รหัสกำกับข้อมูลรายบรรทัดที่ใช้จัดการระบบหลัง\par
บ้าน เพื่อให้การส่งต่อข้อมูลไปยังขั้นตอน Pair Building, Scoring และระบบบริหาร\par
จัดการความสัมพันธ์ลูกค้า (CRM) ดำเนินไปอย่างราบรื่น\par
\begin{center}\small ตารางที่ 3.3 ตัวอย่างการแปลงข้อมูลหลักในขั้นทำความสะอาดและปรับมาตรฐาน\end{center}
ข้อมูลตั้งต้น ตัวอย่างก่อนประมวลผล ฟิลด์ที่สร้างใหม่ ค่าผลลัพธ์ที่ได้\par
userName @A\_A.ron userName\_clean a\_a.ron\par
fullName Mr. Aaron Bird fullName\_clean aaron bird\par
bio\_clean now designer...\par
bio Now designer @Uber...\par
bio\_mentions uber\par
% ---- source PDF page 41 ----
\clearpage
\begin{center}\small ตารางที่ 3.3 (ต่อ) ตัวอย่างการแปลงข้อมูลหลักในขั้นทำความสะอาดและปรับมาตรฐาน\end{center}
ข้อมูลตั้งต้น ตัวอย่างก่อนประมวลผล ฟิลด์ที่สร้างใหม่ ค่าผลลัพธ์ที่ได้\par
location San Francisco, CA location san francisco, ca\par
linkedin.com/in/aar\par
externalUrl\_clean\par
externalUrl HTTPS//www.LinkedIn.c om /in/Aaron/?utm=1\par
on\par
external\_domain linkedin.com\par
outputProfileName aaronbird profile\_id aaronbird\par
\subsection*{3.3.2 การสกัดคุณลักษณะเชิงบริบทและสัญญาณสนับสนุน (Contextual Feature}
Extraction and Supportive Signals)\par
หลังจากผ่านกระบวนการทำความสะอาดข้อมูลเบื้องต้น ขั้นตอนต่อไปคือการสร้าง ฟิลด์ คุณลักษณะสนับสนุนจากข้อมูลบริบท เพื่อแปลงข้อมูลให้เป็นสัญญาณที่มีความละเอียดและ โครงสร้างชัดเจนยิ่งขึ้น โดยมีการจัดการข้อมูลเฉพาะในแต่ละฟิลด์ดังนี้\par
\noindent • bio นอกจากการทำความสะอาดข้อความแล้ว ระบบยังสกัดข้อมูลเชิงโครงสร้าง เช่น\par
รายการ/จำนวน URL และการกล่าวถึง (@mention) เพื่อใช้ประเมินทั้งความหมายและ\par
พฤติกรรมการพิมพ์ของผู้ใช้\par
\noindent • location ระบบจะจำแนกประเภทสถานที่ ได้แก่ ข้อความ พิกัดตัวเลข หรือข้อมูลขยะ\par
และสร้างฟิลด์กำกับ เช่น location\_type, latitude, longitude เพื่อให้สามารถเลือก\par
ระหว่างการเทียบข้อความ หรือการเทียบพิกัดภูมิศาสตร์ได้อย่างเหมาะสม\par
\noindent • externalUrl จัดรูปแบบลิงก์โดยตัดส่วนเกิน (Protocol, www., Query String) ออก\par
แล้วสกัดเฉพาะโดเมนหลักแยกออกมา เพื่อให้ลิงก์เป็นสัญญาณจับคู่ได้แม่นยำขึ้น\par
\noindent • pictureURL ทำหน้าที่เป็นตัวอ้างอิงไปยังฐานข้อมูลภาพที่ประมวลผลไว้ล่วงหน้า เพื่อดึง\par
คุณลักษณะต่าง ๆ เช่น Metadata, Image Hash และคำบรรยายภาพ มาใช้งานได้ทันที\par
\begin{center}\small ตารางที่ 3.4 ตัวอย่างการสร้างฟิลด์สนับสนุนจากข้อมูลบริบท\end{center}
% ---- source PDF page 42 ----
\clearpage
ฟิลด์ตั้งต้น ตัวอย่างก่อนประมวลผล ฟิลด์สนับสนุนที่ สร้างขึ้น\par
ตัวอย่างผลลัพธ์ที่ ได้\par
bio Now designer @Uber... bio\_mentions uber\par
Co-Founder... http//bizible.com\par
bio\_urls http//bizible.co m bio\_url\_count 1\par
location San Francisco, CA location\_type text\_location\par
location\_valid 1 40.7128, -74.0060 latitude 40.7128 longitude -74.0060 location\_type coordinates externalUrl https//www.twitch.tv/Weapo nII?ref=profile\par
externalUrl\_clean twitch.tv/weap\par
onii external\_domain twitch.tv url\_count 1 pictureURL https://.../profile.jpg ข้อมูลภาพที่เตรียมไว้ ล่วงหน้า\par
Metadata, Quality features, Hash similarity, Captions\par
จากตัวอย่างในตารางจะเห็นได้ว่าฟิลด์สนับสนุนไม่ได้เป็นเพียงการเปลี่ยนชื่อคอลัมน์ แต่ เป็นการแยกข้อมูลที่ซ่อนอยู่ในฟิลด์ดิบออกมาให้อยู่ในรูปที่ระบบสามารถนำไปใช้งานได้โดยตรง เช่น bio หนึ่งบรรทัดสามารถแตกออกเป็นข้อความที่สะอาด รายการ mention และรายการ URL ขณะที่ location หนึ่งค่าก็สามารถถูกตีความต่อได้ว่าเป็นข้อความสถานที่ทั่วไปหรือเป็นพิกัด เชิงภูมิศาสตร์จริง การเตรียมฟิลด์ลักษณะนี้ทำให้ขั้น feature engineering ในภายหลังมี สัญญาณหลายมิติมากกว่าการพึ่งข้อความดิบเพียงอย่างเดียว\par
\subsection*{3.3.3 การสร้างชุดข้อมูลคู่เปรียบเทียบเพื่อการเรียนรู้ (Training Pair Formulation)}
% ---- source PDF page 43 ----
\clearpage
เมื่อกระบวนการเตรียมข้อมูลระดับโปรไฟล์เสร็จสมบูรณ์ ขั้นตอนถัดมาคือการ ปรับเปลี่ยนโครงสร้างปัญหา จากการวิเคราะห์ระดับโปรไฟล์เดี่ยว (Profile-level) ไปสู่การ วิเคราะห์ระดับคู่เปรียบเทียบ (Pair-level) เพื่อให้แบบจำลองสามารถเรียนรู้และแยกแยะ ความสัมพันธ์ได้ว่าโปรไฟล์คู่ใดคือบุคคลเดียวกัน และคู่ใดไม่ใช่\par
\noindent • ประเภทของคู่ข้อมูล ได้ทำการสังเคราะห์คู่ข้อมูลออกเป็น 3 ประเภทหลัก ได้แก่\par
\noindent 1) Positive Pairs โปรไฟล์บุคคลเดียวกันจริง โดยจะใช้เฉพาะข้อมูลข้ามแพลตฟอร์ม\par
เพื่อให้สอดคล้องกับขอบเขตของงานวิจัย\par
\noindent 2) Random Negative Pairs สุ่มจับคู่โปรไฟล์ที่เป็นคนละคนกันอย่างชัดเจน\par
\noindent 3) Hard Negative Pairs โปรไฟล์ที่เป็นคนละคนกัน แต่มีข้อมูลหรือชื่อคล้ายคลึงกัน\par
มาก หากแบบจำลองฝึกสอนด้วยคู่ลบที่แยกแยะได้ง่ายเพียงอย่างเดียว จะเกิด\par
แนวโน้มการเรียนรู้แบบผิวเผิน และนำไปสู่ความผิดพลาดประเภทผลบวกลวง (False\par
Positive) ได้ง่ายเมื่อเผชิญกับกรณีที่ชื่อหรือข้อมูลกำกวม การนำคู่ลบแบบยากเข้ามา\par
เป็นส่วนหนึ่งของชุดฝึกสอน จึงเป็นการบังคับให้แบบจำลองเรียนรู้ที่จะแยกแยะกรณี\par
เหล่านี้ออกได้อย่างแม่นยำ\par
\noindent • กระบวนการสกัดและสัดส่วนของชุดคู่ข้อมูล\par
\noindent 1) การกำหนดฐานจากคู่ Positive Pairs ด้วยการรวบรวมโปรไฟล์ที่มี user\_folder\par
เดียวกันเป็นตัวแทนของ Entityเดียวกัน จากนั้นนำมาจับคู่แบบสองตัวเลือก โดยคัด\par
กรองเฉพาะกรณีที่อยู่คนละแพลตฟอร์ม ทำให้ได้คู่บวกตั้งต้นจำนวน 4,060,200 คู่\par
ซึ่งตัวเลขนี้ถูกนำไปใช้เป็นฐานในการคำนวณคู่ลบในลำดับถัดไป\par
\noindent 1) สร้างคู่ Random Negatives ทั่วไป เพื่อสร้างคลาสลบให้มีขนาดใหญ่พอที่จะ\par
สะท้อนความเป็นจริง ได้ตั้งเป้าหมายอัตราส่วนไว้ที่ 5 เท่าของคู่ Positive Pairs\par
(4,060,200 × 5) ซึ่งเท่ากับ 20,301,000 คู่ คู่ข้อมูลเหล่านี้เกิดจากการสุ่มจับคู่โปร\par
ไฟล์ที่อยู่ต่าง Entity กัน และมีการตรวจสอบเพื่อตัดคู่ที่ซ้ำกับคู่ Positive Pairs ออก\par
เสมอ ซึ่งผลลัพธ์ในขั้นตอนนี้สามารถสร้างข้อมูลได้ครบถ้วนตามเป้าหมายที่ตั้งไว้\par
\noindent 2) สกัดคู่ Hard Negatives โดยมีการตั้งเป้าหมายอัตราส่วนไว้ที่ 2 เท่าของคู่บวก\par
(8,120,400 คู่) โดยมีเงื่อนไขคือ ต้องเป็นโปรไฟล์ต่างเอนทิตีที่มีค่า userName\par
คล้ายคลึงกันและผ่านเกณฑ์ประเมินที่ระดับ 0.65 ขึ้นไป รวมทั้งต้องไม่ซ้ำซ้อนกับชุด\par
ข้อมูลที่สร้างไปแล้ว อย่างไรก็ตาม เมื่อดำเนินการกับชุดข้อมูลจริง พบว่ามีคู่ที่ผ่าน\par
เกณฑ์อันเข้มงวดนี้เพียง 302,433 คู่ นัยสำคัญของตัวเลขที่ลดลงนี้ สะท้อนให้เห็นว่า\par
% ---- source PDF page 44 ----
\clearpage
ชุดข้อมูลในงานวิจัยนี้เป็นผลลัพธ์ที่ได้จากการคัดกรองตามความคล้ายคลึงที่เกิดขึ้น\par
จริง\par
\begin{center}\small ตารางที่ 3.5 โครงสร้างชุดคู่ข้อมูลสำหรับการฝึกแบบจำลอง\end{center}
Pair Type จำนวนคู่ บทบาทในการฝึก\par
Positive pairs 4,060,200 ตัวอย่างคู่ที่เป็นบุคคลเดียวกันจริง\par
Random\par
20,301,000 ตัวอย่างคู่ลบทั่วไปเพื่อให้สะท้อน Distribution ของโลก\par
negatives\par
จริง\par
Hard negatives 302,433 ตัวอย่างคู่ลบที่คล้ายกันมาก เพื่อกด False Positives\par
รวมทั้งหมด 24,663,633 ชุด Labeled Pairs ที่ใช้ในขั้น Feature Engineering\par
และ Modeling\par
% ---- source PDF page 45 ----
\clearpage
\begin{center}\small รูปที่ 3.6 Pair building pipeline การสร้างชุดคู่ข้อมูลสำหรับการฝึกแบบจำลอง\end{center}
\subsection*{3.3.4 การสร้าง split artifacts และ feature matrix แบบ chunked}
หลังจากได้ชุดคู่ข้อมูล (Labeled Pairs) ขั้นตอนถัดมาคือการแปลงคู่ข้อมูลเหล่านี้ให้เป็น\par
ข้อมูลเชิงตัวเลขและเวกเตอร์ เพื่อเตรียมความพร้อมสำหรับการฝึกสอนแบบจำลอง กระบวนการ\par
ในขั้นตอนนี้ประกอบด้วยการแบ่งชุดข้อมูลและการสกัดคุณลักษณะ ซึ่งมีรายละเอียดดังต่อไปนี้\par
\noindent 1. การแบ่งชุดข้อมูลระดับเอนทิตี (Entity-Level Data Splitting)\par
\noindent • วัตถุประสงค์ เพื่อป้องกันปัญหาข้อมูลรั่วไหล (Data Leakage) ที่ทำให้ผลประเมิน\par
สูงเกินจริง โดยเลี่ยงการสุ่มรายโปรไฟล์หรือรายคู่ แล้วเปลี่ยนมาแบ่งตามกลุ่มบุคคล\par
แทน\par
\noindent • เกณฑ์การแบ่งสัดส่วน ระบบใช้ user\_folder เป็นตัวแทนของ Entity และทำการ\par
สุ่มแบ่งเอนทิตีทั้งหมดออกเป็น 3 ส่วน ในสัดส่วนประมาณ 70 : 15 : 15 โดย\par
กำหนดให้เป็นชุดฝึกสอน (Training Set) 70\% ส่วนอีก 30\% ที่เหลือถูกแบ่งครึ่งเป็น\par
ชุดตรวจสอบ (Validation Set) และชุดทดสอบ (Test Set) ภายใต้เงื่อนไขบังคับคือ\par
โปรไฟล์ทั้งหมดของบุคคลเดียวกันต้องอยู่ในชุดข้อมูลเดียวกัน\par
% ---- source PDF page 46 ----
\clearpage
\noindent • การจัดการคู่ข้อมูลข้ามชุด เพื่อป้องกัน Data Leakage อย่างเด็ดขาด ระบบจะตัดคู่\par
ข้อมูลที่มีโปรไฟล์อยู่ต่างชุดกันทิ้งทั้งหมด (เช่น ฝั่งหนึ่งอยู่ Train อีกฝั่งอยู่ Test) โดย\par
เก็บไว้เฉพาะคู่ที่ทั้งสองฝั่งอยู่ในชุดเดียวกันเท่านั้น\par
\noindent 2. สถิติและขนาดของชุดข้อมูล (Dataset Statistics)\par
จากหลักการแบ่งและคัดกรองข้างต้น ทำให้ได้ชุดข้อมูลผลลัพธ์ที่พร้อมใช้งาน\par
จำนวน 36,807 โปรไฟล์ และ 15,107,705 คู่ข้อมูล รายละเอียดดังตารางที่ 3.6\par
\begin{center}\small ตารางที่ 3.6 ขนาดของชุดข้อมูลที่ได้หลังแบ่ง train, validation และ test\end{center}
ชุดข้อมูล (Split) จำนวนโปรไฟล์ (Profiles)\par
จำนวนคู่ข้อมูล (Pairs)\par
Training Set 15,848 10,133,486 Validation Set 17,595 4,511,311 Test Set 3,364 462,908 รวมทั้งหมด 36,807 15,107,705\par
หมายเหตุ : จำนวนคู่ข้อมูลหลังการแบ่งชุดข้อมูล (15,107,705 คู่) มีจำนวนลดลงจากชุด\par
Labeled Pairs ตั้งต้น (24,663,633 คู่) ส่วนต่างจำนวน 9,555,928 คู่ที่หายไป เกิดจากการ\par
บังคับใช้เงื่อนไขตัดคู่ข้อมูลที่เชื่อมข้ามระหว่างชุดออกเพื่อป้องกันข้อมูลรั่วไหล\par
\noindent 3. กลยุทธ์การประมวลผลข้อมูลขนาดใหญ่ (Chunked Processing)\par
เนื่องจากปริมาณคู่ข้อมูลมีขนาดใหญ่ถึง 15 ล้านคู่ การประมวลผลทั้งหมดใน\par
หน่วยความจำหลัก (RAM) ในคราวเดียวจึงไม่สามารถทำได้ ระบบจึงประยุกต์ใช้เทคนิค\par
การประมวลผลแบบแบ่งส่วนย่อย (Chunked Processing) โดยอ่านข้อมูลและสร้าง\par
คุณลักษณะทีละ 100,000 คู่ (Batch Size) แล้วจึงนำผลลัพธ์มารวมกันเป็นตารางข้อมูล\par
ขนาดใหญ่ในขั้นตอนสุดท้าย วิธีการนี้ช่วยให้ระบบสามารถจัดการข้อมูลระดับ Big Data\par
ได้อย่างมีประสิทธิภาพ โดยไม่ก่อให้เกิดปัญหาหน่วยความจำเกินพิกัด\par
\noindent 4. การสกัดคุณลักษณะ (Feature Engineering)\par
% ---- source PDF page 47 ----
\clearpage
ในแต่ละรอบการประมวลผลส่วนย่อย (Chunk) ระบบจะดำเนินการสร้าง\par
คุณลักษณะ (Features) โดยนำข้อมูลจากโปรไฟล์ทั้งสองฝั่งมาคำนวณร่วมกัน เพื่อแปลง\par
ให้เป็นเวกเตอร์ตัวเลข 1 แถวต่อ 1 คู่ข้อมูล ซึ่งประกอบด้วยมิติต่าง ๆ ดังนี้\par
\noindent 1. ความคล้ายคลึงของข้อความ (Textual Similarity) การเปรียบเทียบ userName\par
และ fullName\par
\noindent 2. การวิเคราะห์ประวัติส่วนตัว (Bio Analysis) การวัดความคล้ายคลึงของ bio ทั้งใน\par
เชิงคำศัพท์ผ่านเทคนิค TF-IDF และในเชิงความหมายผ่าน Word Embedding\par
\noindent 3. ข้อมูลบริบทและสัญญาณอื่น ๆ (Contextual Signals) สัญญาณความเชื่อมโยง\par
จาก URL, Domain, สถานที่ตั้ง (location), รวมถึงข้อมูลบริบทจากการใช้\par
Mentions และ Hashtags\par
\noindent 4. ข้อมูลพหุรูปแบบ (Multimodal Features) การนำสัญญาณจากข้อมูลภาพมา\par
วิเคราะห์ร่วมด้วย\par
\noindent 5. ข้อควรระวังเชิงระเบียบวิธีวิจัย (Methodological Rigor) ในการสร้างตัวแทนเชิง\par
ข้อความ เช่น การสร้างพจนานุกรมคำศัพท์สำหรับ TF-IDF ระบบจะดำเนินการ\par
ฝึกสอน (Fit) บนข้อมูล bio ในชุดข้อมูลฝึกสอน (Training Set) เท่านั้น จากนั้นจึง\par
นำพจนานุกรมดังกล่าวไปแปลงข้อมูล (Transform) ในชุดข้อมูลตรวจสอบ\par
(Validation Set) และชุดข้อมูลทดสอบ (Test Set) วิธีนี้ช่วยป้องกันไม่ให้ข้อมูลจาก\par
ชุดประเมินไหลย้อนกลับมาส่งผลต่อชุดข้อมูลฝึกสอน\par
\noindent 5. การผสานข้อมูลและวิวัฒนาการของชุดคุณลักษณะ (Feature Merging \&\par
Evolution)\par
หลังการคำนวณแบบ Chunk เสร็จสิ้น ระบบจะรวมผลลัพธ์เป็น Merged\par
Feature Matrix เพื่อส่งต่อจากขั้น Data Preparation ไปยังขั้น Modeling ซึ่ง\par
ความหมายของ จำนวนคอลัมน์ตามลำดับ Pipeline เป็น 3 ระดับ มีดังนี้\par
\noindent • 48 คอลัมน์ (Merged Feature Matrix) ตารางข้อมูลหลังกระบวนการ Merge ซึ่ง\par
ยังคงรวมคอลัมน์ระบุตัวคู่และคอลัมน์กำกับไว้ด้วย ได้แก่ profile\_id\_a,\par
profile\_id\_b, label, และ pair\_type\par
\noindent • 44 คุณลักษณะ (Feature Pool) ชุดคุณลักษณะเชิงตัวเลขที่เหลืออยู่หลังจากตัด\par
คอลัมน์กำกับออก เป็นกลุ่มคุณลักษณะตั้งต้นที่พร้อมป้อนเข้าสู่แบบจำลอง\par
% ---- source PDF page 48 ----
\clearpage
\noindent • 41 คุณลักษณะ (Final Feature Set) ชุดคุณลักษณะที่ถูกคัดเลือกไปใช้งานจริงใน\par
รอบการทดลองหลักของสายการวิเคราะห์แบบ Multimodal\par
แต่สะท้อนให้เห็นลำดับของกระบวนการ กล่าวคือ ขั้น Data Preparation ทำ\par
หน้าที่สร้างตารางให้พร้อมใช้ จากนั้นตัดคอลัมน์ที่ไม่จำเป็นออก และในขั้น Modeling\par
จึงทำการคัดเลือกเฉพาะคุณลักษณะที่สะท้อนสัญญาณได้ดีที่สุดไปใช้งานจริง\par
\begin{center}\small รูปที่ 3.7 ลำดับการลดทอนคอลัมน์และคัดเลือกคุณลักษณะก่อนเข้าสู่กระบวนการ\end{center}
ฝึกสอนตัวแบบ\par
\noindent 6. รูปแบบการคำนวณและการแปลงรูปทางคณิตศาสตร์ (Mathematical\par
Formulations)\par
ในเชิงวิธีการคำนวณ ระบบไม่ได้ป้อนข้อมูลดิบของโปรไฟล์เข้าสู่แบบจำลอง\par
โดยตรง แต่ใช้การเปรียบเทียบข้อมูลจากโปรไฟล์ทั้งสองฝั่งทีละมิติ แล้วแปลงผลลัพธ์ให้\par
เป็นค่าตัวเลขหนึ่งค่าในแต่ละคุณลักษณะ หากกำหนดให้โปรไฟล์สองฝั่งคือ A และ B\par
เวกเตอร์คุณลักษณะของคู่ข้อมูลจะอยู่ในรูปแบบ\par
x A,B = [f 1 (A, B), f 2 (A, B), … , f k (A, B)]\par
% ---- source PDF page 49 ----
\clearpage
โดย f i (A, B) คือฟังก์ชันที่ใช้วัดความคล้ายคลึง ความแตกต่าง อัตราส่วน หรือ\par
สถานะของโปรไฟล์คู่นั้น ในทางปฏิบัติ คุณลักษณะที่ใช้ในงานนี้สามารถอธิบายด้วย\par
รูปแบบการคำนวณหลัก 7 รูปแบบ ดังนี้\par
\noindent 1. การเทียบสถานะแบบตรงกัน (Exact Match) ใช้เมื่อต้องการตรวจสอบความตรงกัน\par
โดยสมบูรณ์ เช่น domain\_exact\_match\par
f(A, B) = \{ 1, ถ้าค่าตรงกันและไม่ว่าง\par
0, ในกรณีอื่น ๆ\par
\noindent 2. การวัดความคล้ายคลึงของข้อความ (String Similarity) ใช้กับชื่อและตำแหน่งที่ตั้ง\par
เพื่อรองรับกรณีที่พิมพ์ไม่ตรงกันแบบตัวต่อตัว โดยมีสมการที่ใช้ ได้แก่\par
\noindent • Jaro-Winkler : f(A, B) = JW(a, b)\par
Lev(a,b)\par
\noindent • Normalized Levenshtein : f(A, B) = 1 −\par
max(|a|,|b|, 1 )\par
TokenSort(a,b)\par
\noindent • Token-sort Ratio : f(A, B) =\par
100\par
\noindent 3. การวัดความทับซ้อนของชุดข้อมูล (Set Overlap) ใช้กับข้อมูลที่มีลักษณะเป็นกลุ่มคำ\par
เช่น โดเมน, URL, Mentions และ Hashtags โดยประยุกต์ใช้ Jaccard Similarity\par
|S A \ensuremath{\cap}S B |\par
\noindent 4. f(A, B) =\par
|s A \ensuremath{\cup}s B |\par
\noindent 5. การวัดความคล้ายคลึงเชิงความหมาย (Semantic Similarity) ใช้เปรียบเทียบเนื้อหาใน\par
ระดับความหมาย สำหรับข้อความ Bio และ Caption ภาพ ผ่านการหาค่า Cosine\par
Similarity ของเวกเตอร์ (เช่น TF-IDF หรือ SBERT)\par
\noindent 6. f(A, B) = cos (v A , v B )f(A, B) = cos (v A , v B )\par
% ---- source PDF page 50 ----
\clearpage
\noindent 7. การวัดความแตกต่างและอัตราส่วน (Difference \& Ratio) ใช้กับคุณลักษณะเชิงสถิติ\par
เช่น ความยาวข้อความ, ความสว่างภาพ, และขนาดไฟล์\par
\noindent • ความแตกต่าง ( Absolute Difference) : f(A, B) = |x A − x A |\par
min (x A ,x B )\par
\noindent • อัตราส่วน ( Min-Max Ratio) : f(A, B) =\par
max (x A ,x B ,\ensuremath{\in})\par
\noindent 8. การวัดความใกล้เคียงเชิงภูมิศาสตร์ (Geographic Proximity) เมื่อข้อมูลตำแหน่งถูก\par
แปลงเป็นพิกัด ระยะทาง (d) จะถูกคำนวณผ่านสูตร Haversine และแปลงเป็นค่าความ\par
คล้ายคลึง\par
f(A, B) = e −d\par
100\par
\noindent 9. การเชื่อมสัญญาณข้ามสื่อ (Cross-modal Token Matching) เป็นการนำข้อมูลข้อความ\par
อธิบายภาพจากฝั่งหนึ่ง ไปตรวจสอบความสอดคล้องกับข้อความพื้นฐานของโปรไฟล์อีก\par
ฝั่งหนึ่งโดยตรง เพื่อเพิ่มมิติการวิเคราะห์ข้ามรูปแบบ ซึ่งวิธีการนี้ช่วยให้ระบบสามารถใช้\par
สัญญาณจากภาพ เพื่อยืนยันความสอดคล้องกับชื่อหรือข้อความของโปรไฟล์อีกฝั่งหนึ่งได้\par
ในงานวิจัยนี้ คุณลักษณะกลุ่มดังกล่าวถูกประเมินด้วย การเปรียบเทียบแบบสองทิศทาง\par
(Bidirectional Comparison) คือ ระบบจะเปรียบเทียบคำบรรยายภาพของฝั่ง A กับ\par
ข้อความของฝั่ง B และในขณะเดียวกันก็เปรียบเทียบคำบรรยายภาพของฝั่ง B กับ\par
ข้อความของฝั่ง A จากนั้นจึงเลือกค่าผลประเมินที่สูงกว่ามาเป็นค่าคุณลักษณะสุดท้าย\par
\noindent 10.\par
\noindent • สำหรับคุณลักษณะ Feature image\_caption\_fullname\_token\_cross สามารถ\par
เขียนสมการได้ดังนี้\par
% ---- source PDF page 51 ----
\clearpage
f cap−fullname (A, B) = max ( TokenSort(c A , n B )\par
100 , TokenSort(c B , n A )\par
100 )\par
\noindent 11. โดยที่ c A และ c B คือคำบรรยายภาพจากฝั่ง A และ B ตามลำดับ ส่วน n A\par
และ n B คือ fullName ที่ผ่านการ Normalize แล้วของแต่ละฝั่ง\par
\noindent • สำหรับคุณลักษณะ image\_caption\_username\_token\_cross จะใช้โครงสร้าง\par
สมการเดียวกัน แต่เปลี่ยนชุดข้อมูลอ้างอิงเป็น userName\par
f cap−username (A, B) = max ( TokenSort(c A , u B )\par
100 , TokenSort(c B , u A )\par
100 )\par
โดยที่ u A และ u B คือข้อมูล userName ที่ผ่านการปรับมาตรฐานแล้ว ค่าผลลัพธ์ที่\par
ได้จะอยู่ในช่วง 0 ถึง 1 ซึ่งค่าที่ยิ่งสูง จะยิ่งสะท้อนให้เห็นว่าคำบรรยายภาพของโปร\par
ไฟล์ฝั่งหนึ่ง มีถ้อยคำที่ใกล้เคียงกับชื่อหรือชื่อบัญชีของอีกฝั่งหนึ่งอย่างมีนัยสำคัญ\par
\noindent 7. โครงสร้างข้อมูลรายคู่ ( Pairwise Feature Matrix Row)\par
Feature Matrix มีขนาด 1 แถวต่อ 1 คู่ข้อมูล หมายความว่า ข้อมูลใน 1 แถว\par
จะประกอบไปด้วยองค์ประกอบ 3 ส่วนหลัก\par
\noindent 12. คอลัมน์ระบุตัวคู่ (Identifier Columns) เช่น profile\_id\_a และ profile\_id\_b\par
\noindent 13. คอลัมน์กำกับ (Label Columns) เช่น label และ pair\_type\par
\noindent 14. ชุดคุณลักษณะเชิงเปรียบเทียบ คือชุดตัวแปรที่ผ่านการคำนวณจากทั้งสองฝั่ง\par
\noindent 15. ระบบในงานวิจัยนี้ไม่ได้ป้อนข้อมูลดิบของโปรไฟล์แต่ละฝั่งเข้าสู่แบบจำลอง\par
โดยตรง แต่เป็นการป้อน ผลลัพธ์จากการเปรียบเทียบความสัมพันธ์ระหว่างสองโปรไฟล์\par
ในรูปแบบเวกเตอร์ตัวเลขแทน ดังรายละเอียดที่แจกแจงในตารางที่ 3.7\par
\noindent 16.\par
% ---- source PDF page 52 ----
\clearpage
\begin{center}\small ตารางที่ 3.7 แจกแจงคุณลักษณะ 41 ตัวที่ถูกใช้จริงในแบบจำลองหลัก\end{center}
กลุ่ม\par
จำนวน วิธีการคำนวณโดย\par
รายชื่อคุณลักษณะ (Features)\par
คุณลักษณะ\par
ย่อ\par
Identity\par
6 วัดความคล้ายของ\par
username\_jaro, username\_lev,\par
string\par
userName และ\par
username\_token\_sort, fullname\_jaro,\par
fullName ด้วย\par
fullname\_lev, fullname\_token\_sort\par
String Similarity\par
หลายรูปแบบ\par
Bio text\par
2 วัดความคล้ายของ\par
bio\_tfidf\_cosine, bio\_sbert\_cosine\par
semantics\par
bio ทั้งเชิงคำและ\par
เชิงความหมาย\par
URL /\par
4 วัดการซ้อนทับของ\par
domain\_jaccard, url\_jaccard,\par
domain\par
URL และโดเมน\par
domain\_count\_a, domain\_count\_b\par
พร้อมนับจำนวน\par
โดเมน\par
Location 2 วัดความคล้ายของ\par
location\_jaro, location\_token\_sort\par
ตำแหน่งที่ตั้งหลัง\par
การปรับมาตรฐาน\par
Mention /\par
4 วัดการซ้อนทับของ\par
mention\_jaccard, hashtag\_jaccard,\par
hashtag\par
Mentions และ\par
hashtag\_count\_a, hashtag\_count\_b\par
Hashtags ใน bio\par
Stylometric 4 วัดความต่างของ\par
style\_caps\_diff, style\_avgword\_diff,\par
รูปแบบการเขียน\par
style\_biolen\_ratio, style\_punct\_diff\par
เช่น การใช้ตัวพิมพ์\par
ใหญ่ ความยาวคำ\par
และวรรคตอน\par
% ---- source PDF page 53 ----
\clearpage
\begin{center}\small ตารางที่ 3.7 (ต่อ) แจกแจงคุณลักษณะ 41 ตัวที่ถูกใช้จริงในแบบจำลองหลัก\end{center}
กลุ่ม\par
จำนวน วิธีการคำนวณโดย\par
รายชื่อคุณลักษณะ (Features)\par
คุณลักษณะ\par
ย่อ\par
Platform /\par
1 เข้ารหัสชนิดของคู่\par
platform\_pair\_code\par
metadata\par
แพลตฟอร์มที่กำลัง\par
เปรียบเทียบ\par
Image\par
3 ตรวจสอบว่าคู่โปร\par
image\_any\_local, image\_both\_local,\par
availability\par
ไฟล์มีข้อมูลภาพ\par
image\_one\_local\_only\par
ส่วนตัว (Local\par
Image) พร้อมใช้\par
หรือไม่\par
Image\par
11 วัดความคล้ายของ\par
image\_phash\_sim, image\_dhash\_sim,\par
statistics\par
ภาพจาก Hash\par
image\_brightness\_diff,\par
และความแตกต่าง\par
image\_contrast\_diff,\par
ของคุณภาพภาพ\par
image\_entropy\_diff, image\_rgb\_l1,\par
เชิงสถิติ\par
image\_filesize\_ratio,\par
image\_face\_count\_diff,\par
image\_face\_area\_diff, image\_blur\_diff,\par
image\_metadata\_any\par
Image-\par
4 ใช้ข้อความอธิบาย\par
image\_caption\_any,\par
caption\par
ภาพ (Caption)\par
image\_caption\_bio\_sbert\_cross,\par
context\par
เชื่อมกับข้อความ\par
image\_caption\_fullname\_token\_cross,\par
พื้นฐานของอีกฝั่ง\par
image\_caption\_username\_token\_cross\par
หนึ่งข้ามสื่อ\par
% ---- source PDF page 54 ----
\clearpage
\begin{center}\small ตารางที่ 3.8 ขนาดของ Feature Matrices ที่พร้อมใช้ในการฝึกสอนและประเมินผล\end{center}
ชุดข้อมูล จำนวนแถว จำนวน\par
บทบาทในการทดลอง\par
คอลัมน์\par
ชุดฝึกสอน (Training Set)\par
10,133,486 48 ใช้ฝึกสอนให้แบบจำลองเรียนรู้รูปแบบ\par
ชุดตรวจสอบ (Validation Set)\par
4,511,311 48 ใช้ประเมินเพื่อเลือกแบบจำลองที่ดีที่สุดและ กำหนด Threshold\par
ชุดทดสอบ (Test Set)\par
462,908 48 ใช้ประเมินประสิทธิภาพการทำงานในขั้น สุดท้าย (Final Evaluation)\par
\section*{3.4 การสร้างแบบจำลอง (Modeling)}
หลังจากได้ feature matrices แล้ว ขั้น Modeling ของโครงการนี้ไม่ได้เลือกแบบจำลอง\par
เพียงตัวเดียวตั้งแต่ต้น แต่ใช้การทดลองเปรียบเทียบหลายสายเพื่อให้ตอบคำถามวิจัย 3 ข้อได้แก่\par
\noindent • แบบจำลองกลุ่มใดจัดการข้อมูลแบบตารางของงานนี้ได้ดีที่สุด?\par
\noindent • การเพิ่มข้อมูลภาพและบริบทข้ามสื่อช่วยแยกแยะโปรไฟล์ที่กำกวมได้จริงหรือไม่?\par
\noindent • เมื่อเทียบกับแบบจำลอง Deep Learning แล้ว แบบจำลองที่ตีความง่ายกว่าและ\par
ดูแลง่ายกว่าสามารถให้ผลลัพธ์ที่เพียงพอหรือดีกว่าหรือไม่?\par
\subsection*{3.4.1 โครงสร้างสายการทดลองหลัก (Experimental Tracks)}
การสร้างแบบจำลองในงานวิจัยนี้ถูกแบ่งออกเป็น 3 สายการทดลองที่แยกขาดจากกัน\par
เพื่อให้แต่ละคำถามวิจัยมีกรอบการพิสูจน์ที่ชัดเจน โดยเน้นการควบคุมตัวแปรไม่ให้ปะปนกัน\par
หลักการนี้ช่วยให้สามารถเปรียบเทียบผลลัพธ์ระหว่างสายได้ และมีเหตุผลรองรับในการคัดเลือก\par
แบบจำลองสุดท้าย ดังสรุปในตารางที่ 3.9\par
% ---- source PDF page 55 ----
\clearpage
\begin{center}\small ตารางที่ 3.9 สรุปสายการทดลองหลักในขั้นตอน Modeling\end{center}
สายการทดลอง แนวทางที่ใช้ สิ่งที่ต้องการพิสูจน์ บทบาทใน งานวิจัย Classical Leakage-safe\par
ฝึกแบบจำลองหลาย ตระกูลบนข้อมูลชุด เดียวกัน ภายใต้การ ควบคุม Data Leakage\par
แบบจำลองกลุ่มใด เหมาะกับข้อมูลเชิง ตารางของงานนี้มาก ที่สุด\par
ใช้เป็นฐาน เปรียบเทียบของ ระบบ\par
Multimodal Suite\par
ทยอยเพิ่มข้อมูลจากภาพ เข้าไปทีละระดับบนแกน ทดลองเดียวกัน\par
ข้อมูลภาพช่วยเพิ่ม ประสิทธิภาพการ แยกแยะคู่กำกวมได้จริง หรือไม่\par
ใช้กำหนด ชุด คุณลักษณะ สำหรับ Final Run\par
ใช้ยืนยันความ เหมาะสมของ คำตอบสุดท้าย\par
Neural Reference\par
ฝึกโครงข่ายประสาทเทียม (Neural Network) บน ชุดข้อมูลเดียวกัน\par
แบบจำลองเชิงลึกให้ ประโยชน์คุ้มค่ากับความ ซับซ้อนที่เพิ่มขึ้นหรือไม่\par
\begin{center}\small ตารางที่ 3.10 สรุปเกณฑ์ที่ใช้คัดเลือกแบบจำลองในขั้น Modeling\end{center}
สายการทดลอง ตัวชี้วัดที่ใช้คัดเลือก วิธีตัดสิน\par
Classical\par
validation AP, validation ROC-\par
เลือก model family ที่ดีที่สุดบน\par
leakage-safe\par
AUC, validation F1\par
validation set แล้วค่อยหา\par
threshold\par
Multimodal\par
composite score, validation\par
เปรียบเทียบแต่ละชุด feature\par
suite\par
AP, validation F1, validation\par
ภายใต้เงื่อนไขเดียวกัน แล้วเลือก\par
ROC-AUC\par
best run\par
Neural\par
validation AP, early stopping\par
ใช้เป็น reference ไม่ได้บังคับให้\par
reference\par
result, validation F1 หลัง\par
ต้องชนะทุกตัวชี้วัด\par
calibration\par
% ---- source PDF page 56 ----
\clearpage
\subsection*{3.4.2 สายการทดลอง Classical Leakage-Safe (การหาฐานเปรียบเทียบจากข้อมูลตาราง)}
แม้ปัญหางานวิจัยนี้จะเป็นการจำแนกประเภท Binary Classification เพื่อระบุว่าโปร\par
ไฟล์สองรายการเป็นบุคคลเดียวกันหรือไม่ แต่ข้อมูลนำเข้าได้ถูกแปลงให้อยู่ในรูปของเมทริกซ์\par
คุณลักษณะเชิงตาราง ที่ผ่านการคำนวณความคล้ายคลึงและความแตกต่างมาแล้ว ลักษณะข้อมูล\par
เช่นนี้จึงเหมาะสมกับ แบบจำลองกลุ่มคลาสสิก (Classical Models) เนื่องจากสามารถเรียนรู้เส้น\par
แบ่งการตัดสินใจได้ดี มีความเสถียร ตีความผลลัพธ์ได้ง่ายและรองรับการปรับเทียบได้อย่างเป็น\par
ระบบ\par
เพื่อพิสูจน์ว่าโครงสร้างข้อมูลที่เตรียมไว้เหมาะกับแบบจำลองกลุ่มใดมากที่สุด\par
กระบวนการในสายการทดลองนี้จึงถูกกำหนดกรอบการทำงานไว้ 4 ประเด็นหลัก ดังนี้\par
\noindent 1) การควบคุมการรั่วไหลของข้อมูล\par
\noindent • การแบ่งชุดข้อมูล ยึดหลักการแบ่งข้อมูลระดับเอนทิตี (Entity-level Split) เพื่อ\par
รับประกันว่าจะไม่มีโปรไฟล์ของบุคคลเดียวกันหลุดรอดไปอยู่ต่างชุดข้อมูล (Train,\par
Validation, Test)\par
\noindent • การปรับสเกลข้อมูล (Feature Scaling) ระบบจะคำนวณค่าสถิติจากชุด Training\par
Set เพียงชุดเดียวเท่านั้น แล้วจึงนำค่าพารามิเตอร์ดังกล่าวไปแปลงข้อมูลในชุด\par
Validation และ Test วิธีนี้ช่วยให้ทุกแบบจำลองถูกฝึกสอนบนฐานข้อมูลเดียวกัน\par
ภายใต้เงื่อนไขเดียวกัน\par
\noindent 2) การคัดเลือกและเปรียบเทียบตระกูลแบบจำลอง\par
กำหนดขอบเขตการเปรียบเทียบไว้ที่แบบจำลอง 3 รูปแบบหลัก ซึ่งเป็นตัวแทนของ\par
แนวคิดการเรียนรู้ที่แตกต่างกันอย่างชัดเจน\par
\noindent • แบบจำลองเชิงเส้น (Linear Models) ใช้เป็นเส้นฐานอ้างอิงที่เรียบง่ายที่สุด\par
(Simple Baseline) เพื่อทดสอบว่าปัญหาเบื้องต้นนี้จำเป็นต้องใช้กลไกที่ซับซ้อน\par
หรือไม่\par
\noindent • แบบจำลองต้นไม้แบบ Boosting (Boosting Trees) ใช้ทดสอบประสิทธิภาพใน\par
การจับความสัมพันธ์แบบ Non-linear และ Feature Interactions ซึ่งเป็นจุดเด่น\par
ของชุดข้อมูลประเภทความคล้ายคลึง\par
% ---- source PDF page 57 ----
\clearpage
\noindent • แบบจำลองต้นไม้แบบ Bagging (Bagging Trees) ใช้ทดสอบแนวทางที่เน้นความ\par
เสถียรสูง (Robustness) และทนทานต่อ Noise ผ่านการเฉลี่ยคำตอบจากหลาย\par
แบบจำลอง\par
\noindent 3) การกำหนดหน้าที่ของชุดข้อมูล\par
เพื่อให้กระบวนการประเมินผลมีความน่าเชื่อถือเชิงวิชาการ ชุดข้อมูลทั้งสามจึงถูกกำหนด\par
หน้าที่แยกขาดจากกัน\par
\noindent • Training Set ใช้สำหรับเรียนรู้พารามิเตอร์ของแบบจำลอง\par
\noindent • Validation Set ใช้ประเมินผลด้วยเมตริกที่รองรับปัญหาข้อมูลไม่สมดุล (Class\par
Imbalance) เพื่อจัดอันดับและคัดเลือกแบบจำลองที่ดีที่สุด นอกจากนี้ยังใช้สำหรับ\par
ทำ Probability Calibration และค้นหา Threshold ที่เหมาะสมที่สุดสำหรับการใช้\par
งานจริง\par
\noindent • Test Set ใช้ประเมินประสิทธิภาพขั้นสุดท้ายด้วย Threshold ที่ได้จาก Validation\par
Set โดยข้อมูลชุดนี้จะไม่มีส่วนในการตัดสินใจเลือกระหว่างทางอย่างเด็ดขาด\par
\noindent 4. บทสรุปและบทบาทของสายการทดลอง\par
เป้าหมายสูงสุดของสายการทดลอง Classical Leakage-Safe ไม่ได้หยุดอยู่แค่การหาว่า\par
แบบจำลองใดได้คะแนนสูงสุด แต่ทำหน้าที่เป็น ฐานอ้างอิงหลัก (Primary Baseline) ของทั้ง\par
โครงการ เพื่อยืนยันให้ชัดเจนว่า การจำกัดการทดลองไว้แค่ 3 ตระกูลแบบจำลองหลัก ช่วย\par
ลดความซ้ำซ้อนจากการฝึกสอนแบบจำลองฟที่การทำงานคล้ายกัน ทำให้ได้ข้อสรุปที่ตรง\par
ประเด็นและคุมต้นทุนการคำนวณได้ อีกทั้งเป็นการหาคำตอบให้ชัดเจนก่อนว่า ถ้ายึดแค่\par
คุณลักษณะข้อความและแอตทริบิวต์ โครงสร้างข้อมูลตารางที่เตรียมมานั้น ทำงานได้ดีกับ\par
แบบจำลองตระกูลไหนที่สุด ซึ่งแบบจำลองที่ชนะจากขั้นนี้ จะถูกยกไปเป็นแกนหลักในสาย\par
Multimodal ทันที เพื่อใช้วัดว่า คุณลักษณะภาพช่วยเพิ่มความแม่นยำแค่ไหนในขั้นตอน\par
ต่อไป โดยไม่ต้องเริ่มทดลองใหม่\par
\subsection*{3.4.3 สายการทดลอง Multimodal Suite (การประเมินคุณลักษณะภาพ)}
หลังจากได้ฐานเปรียบเทียบจากสายแรกแล้ว ระบบได้นำแบบจำลองที่ดีที่สุดมาทดลอง\par
เพิ่มคุณลักษณะจากภาพอย่างเป็นลำดับขั้น เพื่อตรวจสอบว่าข้อมูลภาพช่วยยกระดับความแม่นยำ\par
ได้จริงหรือไม่ โดยยังคงเงื่อนไขการแบ่งชุดข้อมูลไว้ดังเดิม เพื่อให้ผลต่างที่เกิดขึ้นสะท้อนถึง\par
% ---- source PDF page 58 ----
\clearpage
ประสิทธิภาพที่เปลี่ยนไปโดยเกิดจากภาพอย่างแท้จริง โดยมีโครงสร้างการทำงานแบ่งออกเป็น 4\par
ส่วนสำคัญ ดังนี้\par
\noindent 1) การเตรียมข้อมูลส่วนกลาง ระบบจะโหลดข้อมูลที่ต้องใช้ร่วมกันในทุกชุดการทดลอง\par
เช่น SBERT Embeddings และ Image Descriptors มาเก็บไว้ในหน่วยความจำ\par
(Cache) วิธีนี้ช่วยลดความคลาดเคลื่อนจากการคำนวณซ้ำ และลดระยะเวลาในการ\par
ประมวลผลได้อย่างมาก\par
\noindent 2) การระบุระดับของคุณลักษณะ ระบบจะเพิ่มคุณลักษณะตามระดับการทดลอง 3 รูปแบบ\par
ได้แก่\par
\noindent • text\_attr\_hybrid เส้นฐานที่ใช้เฉพาะข้อความและแอตทริบิวต์\par
\noindent • image\_stats เพิ่มคุณลักษณะภาพระดับสถิติ เช่น Perceptual Hash, ความสว่าง,\par
ขนาดไฟล์\par
\noindent • image\_context เพิ่มคุณลักษณะบริบทข้ามสื่อ เช่น การนำข้อความอธิบายภาพไป\par
เทียบกับ SBERT Embeddings ของประวัติส่วนตัวฝั่งตรงข้าม\par
\noindent 3) การฝึกสอนและประเมินผล ทุกชุดการทดลองจะถูกฝึกสอนภายใต้ค่ากำหนดเดียวกัน\par
เพื่อให้เปรียบเทียบกันได้อย่างยุติธรรม ผลลัพธ์จะถูกจัดอันดับด้วยค่าคะแนนรวม ซึ่ง\par
คำนวณจากสมการ\par
Composite Score = 0.45 \ensuremath{\ast} Average Precision + 0.35 \ensuremath{\ast} F1\par
−Score + 0.20 \ensuremath{\ast} ROC −AUC โดยมีหลักการและเหตุผลของการถ่วงน้ำหนักคือการป้องกันความเอนเอียงจาก\par
การตัดสินแบบจำลองด้วยตัวชี้วัดเพียงค่าเดียว งานวิจัยนี้จึงกระจายน้ำหนักการประเมิน\par
ดังนี้\par
\noindent • Average Precision (AP) [น้ำหนัก 45\%] ให้น้ำหนักสูงสุด เนื่องจากชุดข้อมูลมี\par
ปัญหา Class Imbalance สูง ตัวชี้วัดนี้จึงสะท้อนคุณภาพการจัดอันดับคู่ข้อมูลที่มี\par
แนวโน้มเป็นบุคคลเดียวกันได้ดีที่สุด\par
\noindent • F1-score [น้ำหนัก 35\%] ใช้สะท้อนความสมดุลระหว่างความแม่นยำ (Precision)\par
และความครอบคลุม (Recall) ณ จุดตัด (Threshold) ที่เลือก ซึ่งการันตีว่าโมเดล\par
พร้อมนำไปใช้งานจริง\par
% ---- source PDF page 59 ----
\clearpage
\noindent • ROC-AUC [น้ำหนัก 20\%] ให้น้ำหนักน้อยที่สุด เพราะแม้จะใช้วัดความสามารถใน\par
การแยกแยะคลาสในภาพรวมได้ แต่ในชุดข้อมูลที่ Imbalance สูง ค่านี้มักจะมี\par
แนวโน้มสูงเกินจริงและไม่สะท้อนคุณภาพการจัดอันดับที่แท้จริง\par
\noindent 4) สรุปผลการทดลอง จากรายงานผล พบว่าชุดการทดลองแบบ image\_context รวม 41\par
คุณลักษณะ ให้ผลลัพธ์ดีที่สุด โดยมีความสมดุลระหว่างประสิทธิภาพและความสามารถ\par
ในการตีความ ซึ่งช่วยยืนยันความสัมพันธ์ได้ชัดเจนกว่า\par
\begin{center}\small ตารางที่ 3.11 ระดับการทดลองในสาย Multimodal Suite\end{center}
ชุดการทดลอง จำนวน คุณลักษณะ\par
คุณลักษณะที่เพิ่ม จากเส้นฐาน จุดประสงค์การประเมิน\par
text\_attr\_hybrid 23 ใช้ข้อความและแอ ตทริบิวต์เป็นหลัก\par
ใช้เป็นเส้นฐาน (Baseline) สำหรับเปรียบเทียบ\par
image\_stats 37 เพิ่ม Metadata และ Hash ภาพ\par
ทดสอบประโยชน์ของภาพใน ระดับสถิติ\par
ทดสอบประโยชน์ของการ วิเคราะห์บริบทข้ามสื่อ (Cross- modal) ช่วยแยกคู่กำกวมได้ ไหม\par
เพิ่มการเชื่อมโยง ภาพกับข้อความ โดยตรง\par
image\_context 41\par
\subsection*{3.4.4 การฝึกโครงข่ายประสาทเทียมเพื่อเป็นเส้นฐานอ้างอิง (Neural Reference Model)}
แม้งานวิจัยนี้จะเลือกใช้แบบจำลองตระกูล Tree Ensemble เป็นหลัก แต่ได้มีการพัฒนา\par
สายโครงข่ายประสาทเทียมไว้เป็นแบบจำลองอ้างอิง เพื่อพิสูจน์ว่าความซับซ้อนที่เพิ่มขึ้นให้\par
ประโยชน์คุ้มค่าหรือไม่ กระบวนการพัฒนาในสายนี้ประกอบด้วย\par
\noindent • การจัดการข้อมูล ใช้ข้อมูลชุด Train, Validation และ Test เดียวกับสายอื่น แต่มีการทำ\par
สุ่มลดจำนวนข้อมูลกลุ่มลบ (Undersampling Negatives) ในชุด Train เพื่อป้องกัน\par
ไม่ให้ข้อมูลกลุ่มลบครอบงำการฝึกสอน (Class Imbalance) จากนั้นจึงทำการปรับสเกล\par
ข้อมูล\par
% ---- source PDF page 60 ----
\clearpage
\noindent • สถาปัตยกรรม ใช้แบบจำลอง IdentityMLP ขนาด Hidden Layers [256, 128, 64]\par
โดยแต่ละชั้นประกอบด้วย Linear, BatchNorm1d, ReLU และ Dropout\par
\noindent • การฝึกสอน ใช้ฟังก์ชันการสูญเสียแบบ Focal Loss เพื่อรับมือกับความไม่สมดุลของ\par
ข้อมูล ใช้ตัวปรับพารามิเตอร์แบบ Adam และใช้เทคนิค Early Stopping เพื่อหยุดการ\par
ฝึกอัตโนมัติเมื่อค่าความแม่นยำ (Validation AP) เริ่มคงที่\par
\noindent • การปรับตั้งค่า หลังจากการฝึกสอนเสร็จสิ้น ระบบจะทำการปรับเทียบความน่าจะเป็น\par
และค้นหาค่าจุด Threshold ที่ให้ค่า F1 สูงที่สุดเช่นเดียวกับโมเดลกลุ่มอื่น\par
เพื่อความชัดเจนในการอ้างอิง ค่าตั้งต้น (Hyperparameters) ที่สำคัญของแบบจำลอง\par
ทุกกลุ่มในขั้นตอน Modeling ได้ถูกสรุปไว้ในตารางที่ 3.10\par
\begin{center}\small ตารางที่ 3.12 สรุปค่าตั้งต้นไฮเปอร์พารามิเตอร์ของแบบจำลองหลักในขั้นตอน Modeling\end{center}
กลุ่มแบบจำลอง ไฮเปอร์พารามิเตอร์\par
Linear\par
max\_iter = 2000, class\_weight = balanced\par
baseline\par
Tree-based\par
n\_estimators = 300, learning\_rate = 0.05, max\_depth = 6\par
boosting\par
Tree-based\par
n\_estimators = 400, max\_depth = 18, min\_samples\_leaf = 2,\par
bagging\par
class\_weight = balanced\_subsample\par
Neural\par
hidden\_layers = [256, 128, 64], dropout = 0.3, focal\_alpha =\par
reference\par
\noindent 0.25, focal\_gamma = 2.0, lr = 1e-3, epochs = 30, patience = 5,\par
train\_neg\_ratio = 3.0\par
Multimodal\par
random\_neg\_ratio = 0.75, hard\_neg\_ratio = 2.0, seed = 42\par
suite\par
% ---- source PDF page 61 ----
\clearpage
\subsection*{3.4.5 การปรับเทียบ (Calibration) และการกำหนดจุดตัด (Threshold)}
ไม่ว่าจะเป็นแบบจำลองจากสายการทดลองใด ทุกแบบจำลองจะต้องผ่านกระบวนการ\par
จัดการผลลัพธ์ 2 ขั้นตอนเสมอ โดยดำเนินการบน Validation Set เท่านั้น ได้แก่\par
\noindent 1) Probability Calibration นำคะแนนความน่าจะเป็นที่ได้จากแบบจำลองไปผ่าน\par
อัลกอริทึม Isotonic Regression บนชุด Validation เพื่อปรับการกระจายตัวของ\par
คะแนนให้สะท้อนความน่าจะเป็นที่แท้จริง\par
\noindent 2) Threshold Selection ค้นหาจุดตัดที่สร้างสมดุลระหว่างการจับคู่ผิด (False Positive)\par
และการพลาดคู่จริง (False Negative) ให้สอดคล้องกับวัตถุประสงค์การใช้งาน\par
ผลลัพธ์จากแบบจำลองในงานวิจัยนี้จะต้องถูกส่งต่อไปยังระบบตัดสินใจของ CRM หากไม่\par
มีการปรับเทียบคะแนน การตั้งค่าจุดตัดเพื่อใช้งานจริงจะขาดความเสถียร นอกจากนี้ การบังคับใช้\par
ให้ทั้งสองขั้นตอนดำเนินการบนชุด Validation เท่านั้น เป็นหลักการสำคัญที่ช่วยรักษาความ\par
น่าเชื่อถือทางวิชาการและป้องกันการอธิบายผลเกินจริง (Overfitting)\par
\subsection*{3.4.6 หลักเกณฑ์และผลการคัดเลือกแบบจำลองสุดท้าย (Final Model Selection)}
การคัดเลือกแบบจำลองเพื่อใช้งานจริง พิจารณาจาก 4 มิติหลัก ได้แก่\par
\noindent • ประสิทธิภาพบนชุดข้อมูลที่ควบคุม Leakage\par
\noindent • ความเหมาะสมกับลักษณะของตารางคุณลักษณะ\par
\noindent • ความสามารถในการอธิบายข้อผิดพลาด\par
\noindent • ความพร้อมในการเชื่อมต่อกับระบบ CRM\par
แบบจำลองที่ถูกเลือกเป็น Final Model คือ Gradient Boosting จากชุดการทดลอง\par
image\_context สาเหตุหลักมาจากข้อมูลคุณลักษณะส่วนใหญ่เป็นค่าความคล้ายคลึงและข้อมูล\par
เชิงสถิติ ซึ่งแบบจำลองกลุ่ม Tree Ensemble สามารถจับความสัมพันธ์แบบไม่เป็นเชิงเส้นได้ดี\par
เยี่ยม โดยไม่จำเป็นต้องพึ่งพาโครงข่ายประสาทเทียมที่ใช้ทรัพยากรสูงกว่า ดังสรุปผลการประเมิน\par
ในตารางที่ 3.11\par
% ---- source PDF page 62 ----
\clearpage
\begin{center}
\small
\begin{tabularx}{\linewidth}{p{2.7cm}p{3.1cm}p{1.5cm}XX}
\toprule
สายการทดลอง & แบบจำลอง/ชุดคุณลักษณะ & จำนวนคุณลักษณะ & ผล Validation หรือเกณฑ์คัดเลือก & ผลการทดสอบ (Test) \\
\midrule
1. Classical Leakage-safe & Gradient Boosting & 22 & Average Precision = 0.9733\newline F1 = 0.9368 & Average Precision = 0.9703\newline F1 = 0.9337 \\
2. Multimodal Suite (Final Model) & Gradient Boosting (\texttt{image\_context}) & 41 & Composite Score = 0.9613\newline Average Precision = 0.9776\newline F1 = 0.9339 & Average Precision = 0.9789\newline F1 = 0.9340 \\
3. Neural Reference & IdentityMLP & 22 & Average Precision $\approx$ 1.0000\footnotemark & Average Precision = 0.7344\newline F1 = 0.8296 \\
\bottomrule
\end{tabularx}
\par ตารางที่ 3.13 สรุปผลเชิงตัวเลขและการคัดเลือกแบบจำลองในแต่ละสายการทดลอง
\end{center}
\footnotetext{ค่า Average Precision $\approx$ 1.0000 เป็นค่าที่ปรากฏในรายงานเดิมของสาย Neural Reference แต่ไม่ได้ระบุชุดข้อมูลไว้อย่างชัดเจน จึงไม่ควรนำมาเปรียบเทียบกับผล Test โดยตรง ส่วนผล Test ที่ตรวจสอบได้คือ Average Precision = 0.7344 และ F1 = 0.8296}
% ---- source PDF page 63 ----
\clearpage
\begin{center}\small รูปที่ 3.7 ลำดับขั้น Model Development Pipeline\end{center}
\begin{center}\small รูปที่ 3.8 ความสัมพันธ์ของสายการทดลองและจุดบรรจบ\end{center}
% ---- source PDF page 64 ----
\clearpage
\section*{3.5 การประเมินผล (Evaluation)}
การประเมินผลในงานวิจัยนี้ ไม่ได้วัดเฉพาะความแม่นยำของแบบจำลองเพียงอย่างเดียว\par
แต่ถูกออกแบบให้ประเมินประสิทธิภาพการทำงานของระบบตลอดทั้งสาย (End-to-End\par
Pipeline) เนื่องจากผลลัพธ์ของงาน Identity Resolution ในระบบ CRM ต้องอาศัยการทำงาน\par
ร่วมกันหลายขั้นตอน\par
\subsection*{3.5.1 กรอบการประเมินผลและตัวชี้วัดที่ใช้}
\noindent • ระดับการประเมินออกเป็น 4 ระดับ ในการพัฒนาระบบแต่ละช่วง ดังนี้\par
\noindent 1. การประเมินขั้นคัดกรองข้อมูล (Retrieval Evaluation) ตรวจสอบว่าระบบลด\par
ขนาดพื้นที่ค้นหาได้จริงหรือไม่ และทำคู่ข้อมูลจริงหลุดหายไปหรือไม่\par
\noindent 2. การประเมินแบบจำลอง (Model Evaluation) ตรวจสอบความสามารถในการ\par
แยกแยะคู่ที่เป็นบุคคลเดียวกัน (MATCH) ออกจากคู่ที่ไม่ใช่ (NO\_MATCH)\par
\noindent 3. การประเมินการปรับเทียบ (Calibration Evaluation) ตรวจสอบว่าคะแนนความ\par
น่าจะเป็น (Probability Score) ที่ได้ มีความน่าเชื่อถือพอที่จะใช้ตั้งจุดตัด\par
(Threshold) หรือไม่\par
\noindent 4. การประเมินระดับใช้งานจริง (Production Evaluation) ตรวจสอบว่าเมื่อนำ\par
ระบบไปใช้จริง ภาระงานและผลลัพธ์ปลายทางมีความเหมาะสมกับการทำงานใน\par
CRM หรือไม่\par
\noindent • ตัวชี้วัดที่ใช้ (Evaluation Metrics) เนื่องจากชุดข้อมูลมีปัญหาความไม่สมดุลของคลาส\par
(Class Imbalance) สูงมาก การใช้ความแม่นยำ (Accuracy) เพียงอย่างเดียวจะไม่\par
สะท้อนคุณภาพที่แท้จริง งานวิจัยนี้จึงเลือกใช้ตัวชี้วัดดังนี้\par
o Precision และ Recall ใช้ดูความสมดุล (Trade-off) ระหว่างการจับคู่ผิดกับการ\par
พลาดคู่จริง\par
o F1-score ค่าเฉลี่ยฮาร์มอนิกของ Precision และ Recall\par
o Average Precision (AP) เหมาะกับการประเมินคุณภาพการจัดอันดับในชุดข้อมูล\par
ที่ไม่สมดุล\par
% ---- source PDF page 65 ----
\clearpage
o ROC-AUC ประเมินความสามารถการแยกแยะคลาสในภาพรวม โดยไม่อิงจุดตัดใด\par
จุดตัดหนึ่ง\par
o Expected Calibration Error (ECE) ตรวจสอบว่าคะแนนที่โมเดลให้ สื่อถึงระดับ\par
ความมั่นใจได้แม่นยำแค่ไหน\par
o Confusion Matrix และ Classification Report ใช้สำหรับดูรายละเอียดเชิงลึก\par
ว่าโมเดลทำนายคลาสใดผิดพลาดในสัดส่วนเท่าใด\par
\begin{center}\small ตารางที่ 3.14 สรุประดับการประเมินผลและตัวชี้วัดหลักที่ใช้ในงานวิจัย\end{center}
ระดับการ ประเมิน\par
คำถามวิจัยที่ต้องการตอบ ตัวชี้วัดหลัก (Key Metrics)\par
\noindent 1. Retrieval ระบบลด Search Space ได้ ดีแค่ไหน? คู่จริงหลุดไปหรือไม่?\par
Exact matches, Candidate pairs, Search-space reduction, Ground- truth coverage 2. Model แบบจำลองแยก MATCH กับ NO\_MATCH ได้ดี เพียงใด?\par
Precision, Recall, F1-score, ROC-AUC, Average Precision, Confusion Matrix\par
\noindent 3. Calibration คะแนนที่ปรับเทียบแล้ว สื่อ\par
ECE ก่อนและหลัง Calibration, Threshold stability 4. Production จุดตัดที่เลือก ส่งผลดีต่อ\par
ถึงความมั่นใจได้ดีขึ้นหรือไม่?\par
MATCH, REVIEW, NO\_MATCH, Review queue size, Unified profiles\par
ภาระงานและการใช้งานจริง หรือไม่?\par
\subsection*{3.5.2 การประเมินขั้น Retrieval และ Candidate Generation}
จุดประสงค์ของขั้นตอนนี้คือ การตรวจสอบว่าวิธีการ Exact-first และการจัดกลุ่ม\par
Blocking สามารถลดปริมาณคู่ข้อมูลมหาศาล ให้เหลือเฉพาะกลุ่มตัวเลือก (Candidate Set) ที่\par
ระบบประมวลผลไหว โดยประกอบด้วยการประเมิน 3 มิติหลัก\par
\noindent 1. คุณภาพของ Exact-first\par
\noindent 2. ประสิทธิภาพของการทำ Blocking\par
% ---- source PDF page 66 ----
\clearpage
\noindent 3. ความครอบคลุมคู่จริง (Ground-truth Coverage)\par
\begin{center}\small ตารางที่ 3.15 กรอบการประเมินขั้น Retrieval และ Candidate Generation\end{center}
มิติที่ประเมิน ตัวชี้วัดหลัก จุดประสงค์\par
Exact-first\par
จำนวน Exact matches,\par
ตรวจสอบว่ากฎที่เข้มงวด ให้คู่ที่\par
quality\par
Precision\par
เชื่อถือได้มากเพียงใด\par
Blocking\par
Candidate pairs, Search-\par
ตรวจสอบว่าระบบลดภาระการ\par
efficiency\par
space reduction\par
คำนวณได้มากเพียงใด\par
Coverage of\par
Ground-truth coverage ตรวจสอบว่าคู่จริงไม่หลุดหายไป และ\par
true pairs\par
ถูกส่งต่อให้แบบจำลอง\par
\subsection*{3.5.3 การประเมินแบบจำลองเชิงการจำแนก (Classification Model Evaluation)}
จุดประสงค์ของขั้นตอนนี้คือ ตรวจสอบว่า เมื่อกลุ่ม Candidate Set ผ่านขั้นตอนการคัด\par
กรองมาแล้ว แบบจำลองจะสามารถจำแนกคู่ที่เป็นบุคคลเดียวกัน (MATCH) ออกจากคู่ที่ไม่ใช่\par
บุคคลเดียวกัน (NO\_MATCH) ได้ดีเพียงใด ในการประเมินผลได้มีการออกแบบการประเมิน\par
ออกเป็น สายการทดลองแบบคู่ขนาน ตอบคำถามเชิงโครงสร้างของระบบ 3 ประการ ได้แก่\par
\noindent 1. เชิงตระกูลแบบจำลอง (Model Family) แบบจำลองที่มีกลไกต่างกัน ชนิดใดสามารถ\par
จัดการกับโครงสร้างข้อมูลแบบตารางของงานนี้ได้ดีกว่ากัน\par
\noindent 2. เชิงชนิดของคุณลักษณะ (Feature Modality) การลงทุนดึงข้อมูลภาพและบริบทข้าม\par
สื่อมาใช้งาน เพิ่มความแม่นยำได้คุ้มค่ากับทรัพยากรการประมวลผลที่ต้องเสียไปหรือไม่\par
\noindent 3. เชิงความซับซ้อน แบบจำลองเชิงลึก (Neural Network) ให้ผลลัพธ์ที่ก้าวกระโดดจาก\par
แบบจำลองการเรียนรู้ของเครื่อง (Machine Learning) พื้นฐาน มากพอที่จะคุ้มค่ากับ\par
ต้นทุนการดูแลรักษาระบบในระยะยาวหรือไม่\par
เพื่อให้การประเมินมีความน่าเชื่อถือตามหลักวิชาการ ชุดข้อมูลจะถูกแบ่งหน้าที่อย่าง\par
เด็ดขาด โดยใช้ชุดฝึกสอน (Train Set) เพื่อให้โมเดลเรียนรู้ ใช้ชุดตรวจสอบ (Validation Set)\par
% ---- source PDF page 67 ----
\clearpage
เพื่อคัดเลือกโมเดลและกำหนดจุดตัด (Operating Point) และใช้ชุดทดสอบ (Test Set) สำหรับ\par
การประเมินผลขั้นสุดท้ายเพียงครั้งเดียวเท่านั้น ห้ามนำผลจากชุดทดสอบมาใช้ช่วยตัดสินใจเลือก\par
ระหว่างทางโดยเด็ดขาด\par
\begin{center}\small ตารางที่ 3.16 กรอบการประเมินแบบจำลองในสายการทดลองหลัก\end{center}
สายการทดลอง จุดประสงค์ของการประเมิน ตัวชี้วัดหลัก\par
Classical\par
ใช้เปรียบเทียบ model family บน\par
Precision, Recall, F1-score,\par
ข้อมูลเชิงตารางภายใต้เงื่อนไขที่ควบคุม\par
ROC-AUC, Average\par
leakage-safe\par
leakage\par
Precision\par
Multimodal\par
ใช้ตรวจผลของการเพิ่มข้อมูลภาพและ\par
Precision, Recall, F1-score,\par
ROC-AUC, Average\par
suite\par
บริบทข้ามสื่อ\par
Precision\par
Neural\par
ใช้เป็นตัวอ้างอิงเพื่อเปรียบเทียบกับ\par
Precision, Recall, F1-score,\par
ROC-AUC, Average\par
reference\par
แนวทางที่ซับซ้อนกว่า\par
Precision\par
\subsection*{3.5.4 การประเมิน Calibration และผลในระดับ Production Workflow}
ขั้นตอนนี้ทำหน้าที่เป็นเชื่อมระหว่างตัวเลขประสิทธิภาพของแบบจำลอง กับการนำ\par
ผลลัพธ์ไปประยุกต์ใช้งานจริงในระบบธุรกิจ โดยไม่ได้มองแค่ว่าโมเดลทำนายถูกหรือผิด แต่\par
ประเมินว่าคะแนนจากโมเดลสามารถนำไปควบคุมผลลัพธ์เชิงปฏิบัติการได้หรือไม่ แบ่งการ\par
พิจารณาออกเป็น 3 ระดับ ดังนี้\par
\noindent 1) การประเมินการปรับเทียบ (Calibration Evaluation)\par
ตรวจสอบว่าคะแนนความน่าจะเป็นดิบสะท้อนความมั่นใจที่แท้จริงของ\par
แบบจำลองหรือไม่ โดยวัดความคลาดเคลื่อนจากค่า Expected Calibration Error\par
(ECE) ทั้งก่อนและหลังปรับเทียบ รวมถึงดูความเสถียรของการกำหนดจุดตัด บน\par
% ---- source PDF page 68 ----
\clearpage
Validation Set ขั้นตอนนี้สำคัญมาก เพราะหากคะแนนขาดความน่าเชื่อถือ จะทำให้ชั้น\par
การตัดสินใจ (Decision Layer) ทำงานผิดพลาด\par
\noindent 2) การประเมินชั้นการตัดสินใจ (Decision Layer Evaluation)\par
ตรวจสอบว่าเมื่อนำคะแนนที่ปรับเทียบแล้วมากำหนดจุดตัดระบบสามารถแบ่ง\par
ผลลัพธ์ออกเป็น 3 กลุ่มปฏิบัติการได้อย่างสมเหตุสมผลหรือไม่ ได้แก่\par
\noindent • MATCH ข้อมูลตรงกัน (รวมโปรไฟล์อัตโนมัติ)\par
\noindent • REVIEW ข้อมูลมีความกำกวม (ส่งให้มนุษย์ตรวจสอบ)\par
\noindent • NO\_MATCH ข้อมูลไม่ตรงกัน (ปฏิเสธการรวมโปรไฟล์)\par
\noindent 3) การประเมินระดับการใช้งานจริง (Production Workflow Evaluation):\par
ประเมินผลกระทบที่มีต่อระบบ CRM และผู้ใช้งาน โดยพิจารณาจากภาระงานปลายทาง\par
(Downstream Workload) เช่น ขนาดคิวที่ต้องตรวจสอบและปริมาณโปรไฟล์ที่ถูกรวมสำเร็จ\par
เพื่อยืนยันว่าจุดตัดที่เลือกสอดคล้องกับข้อจำกัดด้านทรัพยากรบุคคลและนำไปใช้งานได้จริง\par
\begin{center}\small ตารางที่ 3.17 กรอบการประเมินผลระดับ Calibration และ Production Workflow\end{center}
ระดับการ\par
คำถามวิจัยที่ต้องการตอบ ตัวชี้วัดหลัก (Key Metrics)\par
ประเมิน\par
\noindent 1. Calibration คะแนนที่ปรับเทียบแล้ว สื่อถึง\par
ECE ก่อนและหลัง Calibration,\par
ระดับความมั่นใจได้ดีและเสถียรขึ้น\par
Threshold stability\par
หรือไม่?\par
\noindent 2. Decision\par
จุดตัดที่กำหนด สามารถแบ่งกลุ่ม\par
Match threshold, Review\par
MATCH, REVIEW และ\par
threshold, Tier precision, Final\par
Layer\par
NO\_MATCH ได้เหมาะสมหรือไม่?\par
precision/recall\par
\noindent 3. Production\par
ระบบทั้งเส้นทาง (End-to-End)\par
Review queue size, Unified\par
สามารถนำผลไปใช้งานในระบบ\par
profiles, Lead scores,\par
Workflow\par
CRM ได้จริงหรือไม่?\par
Downstream workload\par
% ---- source PDF page 69 ----
\clearpage
\begin{center}\small รูปที่ 3.9 Evaluation Framework: กรอบการประเมินผลของระบบในหลายระดับ\end{center}
% ---- source PDF page 70 ----
\clearpage
% ส่วนเพิ่มเติมในบทที่ 3 ต่อจากหัวข้อ 3.5 และก่อนบทที่ 4
\section*{3.6 การพัฒนานโยบายการตัดสินใจด้วยขั้นตอนวิธีเชิงพันธุกรรม}
งานในส่วนนี้เป็นการพัฒนาต่อจากกระบวนการเชื่อมโยงอัตลักษณ์ที่อธิบายไว้ในหัวข้อ 3.1--3.5 โดยยังใช้การเตรียมข้อมูล การสร้างคู่ผู้สมัคร และคะแนนจากแบบจำลองเป็นพื้นฐานเดิม สิ่งที่เพิ่มเติมคือวิธีนำคะแนนและหลักฐานของแต่ละคู่มาใช้ตัดสินใจว่า ควรรวมข้อมูลอัตโนมัติ ส่งให้เจ้าหน้าที่ตรวจสอบ หรือไม่รวมข้อมูล ผู้วิจัยจึงเพิ่มขั้นตอนการค้นหานโยบายการตัดสินใจด้วยขั้นตอนวิธีเชิงพันธุกรรม (Genetic Algorithm: GA) เพื่อให้การเลือกเกณฑ์มีหลักฐานจากชุดข้อมูลตรวจสอบรองรับ และสะท้อนภาระงานที่เกิดขึ้นในระบบ CRM ได้ชัดเจนขึ้น

การพัฒนาส่วนนี้อยู่ภายใต้กรอบ CRISP-DM เดิม โดยขยายรายละเอียดในขั้นการสร้างแบบจำลองและการประเมินผล รวมถึงเตรียมผลการตัดสินใจสำหรับกระบวนการปลายทาง ขอบเขตของงานที่เพิ่มจึงครอบคลุมการออกแบบนโยบาย การค้นหาพารามิเตอร์ การทดสอบ และการส่งออกผลรายคู่ ส่วนผลเชิงตัวเลขจากการเปรียบเทียบนโยบายควรนำเสนอในบทผลการทดลองแยกจากวิธีดำเนินงานนี้

\subsection*{3.6.1 เหตุผลและขอบเขตของการพัฒนาเพิ่มเติม}
ในระบบเดิม แบบจำลองให้คะแนนว่าคู่โปรไฟล์มีแนวโน้มเป็นบุคคลเดียวกันมากน้อยเพียงใด จากนั้นจึงใช้เกณฑ์สองระดับแบ่งคำตัดสินเป็น MATCH, REVIEW และ NO\_MATCH วิธีนี้ทำให้ระบบมีทางเลือกสำหรับคู่ที่ยังไม่แน่ชัด อย่างไรก็ตาม การกำหนดเกณฑ์ด้วยค่าคงที่อาจยังไม่สอดคล้องกับสัดส่วนข้อมูล ความผิดพลาด และจำนวนคู่ที่เจ้าหน้าที่ต้องตรวจสอบในแต่ละชุดข้อมูล

ตัวอย่างเช่น การลดเกณฑ์ MATCH อาจทำให้ระบบเชื่อมโยงคู่ที่ถูกต้องได้มากขึ้น แต่ก็อาจรวมคนละคนเข้าด้วยกันมากขึ้นด้วย ในทางกลับกัน การเพิ่มเกณฑ์อาจลดการรวมผิด แต่ทำให้คู่จำนวนมากค้างอยู่ใน REVIEW การเลือกนโยบายจึงต้องพิจารณาทั้งความถูกต้องของการรวมข้อมูลและภาระงานที่ตามมา ผู้วิจัยใช้ GA ค้นหาชุดพารามิเตอร์ที่ทำให้ต้นทุนการตัดสินใจบนชุด validation ต่ำลง แล้วจึงตรวจสอบผลด้วยตัวชี้วัดด้านความแม่นยำและปริมาณงานร่วมกัน

ในการเปรียบเทียบนโยบาย ผู้วิจัยตรึงคะแนนจากแบบจำลองไว้ เพื่อให้ความแตกต่างของผลลัพธ์อธิบายได้จากการเปลี่ยนชั้นการตัดสินใจ ส่วนการปรับคุณลักษณะหรือฝึกแบบจำลองในสายการทดลองอื่นถือเป็นงานอีกส่วนหนึ่ง น้ำหนักที่ GA ค้นหาในขั้นตอนนี้จึงเป็นน้ำหนักของนโยบายการตัดสินใจ ไม่ใช่น้ำหนักภายในแบบจำลอง Gradient Boosting และไม่ใช่ค่าความสำคัญของคุณลักษณะจากแบบจำลองโดยตรง

\subsection*{3.6.2 ลำดับการพัฒนาจากกฎเดิมสู่นโยบายแบบถ่วงน้ำหนัก}
ผู้วิจัยพัฒนานโยบายเป็นลำดับ โดยเริ่มจากกฎที่ใช้คะแนนจากแบบจำลองและเกณฑ์สองระดับเป็นฐานเปรียบเทียบ จากนั้นเพิ่มการใช้ความคล้ายของชื่อช่วยตัดสินคู่ที่อยู่ในช่วง REVIEW และจึงขยายไปสู่นโยบายที่รวมหลักฐานหลายด้านเป็นคะแนนเดียว ตารางที่ 3.18 สรุปความแตกต่างของแนวทางที่ใช้ในการศึกษา

\begin{center}\small
\begin{tabularx}{\linewidth}{p{3.4cm}XX}
\toprule
แนวทาง & สิ่งที่ใช้ตัดสินใจ & บทบาทในการทดลอง \\
\midrule
เกณฑ์สองระดับจากคะแนนโมเดล & คะแนนโมเดลและเกณฑ์ MATCH กับ REVIEW & ใช้ตรวจว่าการปรับเกณฑ์เพียงอย่างเดียวเพียงพอหรือไม่ \\
GA ที่ใช้ชื่อช่วยตัดสิน & คะแนนโมเดล เกณฑ์สองระดับ และเกณฑ์ความคล้ายของชื่ออีกสองค่า & เป็นแนวทางต่อยอดระยะแรกและฐานเปรียบเทียบของนโยบายใหม่ \\
GA แบบถ่วงน้ำหนักขั้นเดียว & คะแนนโมเดลร่วมกับหลักฐานที่มีข้อมูลของแต่ละคู่ และเกณฑ์สองระดับ & ค้นหาน้ำหนักและเกณฑ์ร่วมกัน แล้วตัดสินจากคะแนนรวมครั้งเดียว \\
\bottomrule
\end{tabularx}
\par ตารางที่ 3.18 ลำดับการพัฒนานโยบายการตัดสินใจ
\end{center}

นโยบายที่ใช้ชื่อช่วยตัดสินมีโครโมโซม 4 ค่า ได้แก่ เกณฑ์ MATCH เกณฑ์ REVIEW เกณฑ์ความคล้ายของชื่อสำหรับเลื่อนคู่ขึ้นเป็น MATCH และเกณฑ์สำหรับลดคู่ลงเป็น NO\_MATCH โดยเกณฑ์ของชื่อทำงานเฉพาะคู่ที่คะแนนโมเดลอยู่ในช่วง REVIEW เท่านั้น ข้อจำกัดของแนวทางนี้คือ ชื่อถูกกำหนดให้เป็นหลักฐานหลักตั้งแต่ก่อนเริ่มค้นหา GA จึงเลือกค่าจุดตัดของชื่อได้ แต่ยังไม่ได้เปรียบเทียบว่าหลักฐานชนิดอื่นเหมาะสมกว่าหรือไม่

เพื่อศึกษาข้อจำกัดดังกล่าว ผู้วิจัยเพิ่มการทดลองใช้หลักฐานแต่ละกลุ่มเป็นตัวช่วยตัดสิน และต่อยอดเป็นนโยบายแบบถ่วงน้ำหนักขั้นเดียว (One-Stage Weighted Policy) ซึ่งรวมคะแนนโมเดลและหลักฐานที่มีอยู่ก่อนเปรียบเทียบกับเกณฑ์ MATCH และ REVIEW แนวทางนี้ไม่มีเกณฑ์เลื่อนขึ้นและลดลงของชื่อซ้อนอยู่ในช่วง REVIEW อีกชั้นหนึ่ง แต่ยังเปิดให้ GA เลือกน้ำหนักของชื่อได้ตามผลบน validation

\subsection*{3.6.3 การเตรียมข้อมูลสำหรับชั้นการตัดสินใจ}
ข้อมูลสำหรับนโยบายต้องอ้างอิงคู่เดียวกับที่ระบบให้คะแนนไว้จริง ผู้วิจัยจึงจัดเตรียมตารางคุณลักษณะสำหรับการตัดสินใจ โดยเก็บรหัสโปรไฟล์ทั้งสองฝั่ง คะแนนโมเดล กลุ่มข้อมูลที่ใช้ประเมิน ค่าความจริงสำหรับคำนวณผล และหลักฐานแต่ละด้านไว้ด้วยกัน พร้อมตัวบ่งชี้ว่าหลักฐานด้านนั้นใช้งานได้หรือไม่ การมีรหัสคู่ช่วยตรวจสอบว่าคะแนนและคุณลักษณะไม่ได้ถูกนำไปจับกับคู่ผิดระหว่างการประมวลผลแบบแบ่งส่วน

โครงสร้างนโยบายกำหนดช่องข้อมูลไว้ 8 ช่อง ได้แก่ คะแนนโมเดล ชื่อ ประวัติย่อ ลิงก์ ตำแหน่งที่อยู่ สัญญาณทางสังคม ภาพ และรูปแบบการเขียน ดังตารางที่ 3.19 ค่าหลักฐานถูกจัดให้อยู่ในช่วง 0 ถึง 1 และมีทิศทางเดียวกัน คือค่าสูงแสดงความสอดคล้องมากกว่า ส่วนคุณลักษณะที่เดิมเป็นค่าความแตกต่างต้องแปลงทิศทางก่อนนำมาใช้

\begin{center}\small
\begin{tabularx}{\linewidth}{p{2.5cm}X}
\toprule
กลุ่มข้อมูล & ความหมายในชั้นการตัดสินใจ \\
\midrule
คะแนนโมเดล & คะแนนความเป็นบุคคลเดียวกันจากแบบจำลองที่ตรึงไว้ \\
ชื่อ & ความคล้ายของชื่อผู้ใช้หรือชื่อที่แสดงในโปรไฟล์ \\
ประวัติย่อ & ความคล้ายของข้อความหรือความหมายในประวัติย่อ \\
ลิงก์ & ความสอดคล้องของ URL หรือโดเมนที่อ้างอิง \\
ตำแหน่งที่อยู่ & ความคล้ายของข้อมูลสถานที่เมื่อมีข้อมูลที่ใช้งานได้ \\
สัญญาณทางสังคม & ความสอดคล้องของการกล่าวถึงบัญชีหรือแฮชแท็กในคุณลักษณะที่เตรียมไว้ \\
ภาพ & หลักฐานจากภาพหรือข้อความประกอบภาพตามที่มีในตารางคุณลักษณะ \\
รูปแบบการเขียน & ความสอดคล้องของลักษณะการเขียน เช่น ความยาวและการใช้เครื่องหมาย \\
\bottomrule
\end{tabularx}
\par ตารางที่ 3.19 กลุ่มข้อมูลที่โครงสร้างนโยบายรองรับ
\end{center}

การรองรับ 8 ช่องไม่ได้หมายความว่าทุกรอบทดลองใช้ข้อมูลได้ครบทั้ง 8 ด้าน ในรอบข้อมูลเต็มวันที่ 14 กันยายน 2569 ตารางคุณลักษณะไม่มีช่อง social และไม่มีรายการ location ที่พร้อมใช้งาน จึงปิดการใช้งานสองด้านนี้และกำหนดน้ำหนักเป็นศูนย์ตั้งแต่เริ่มค้นหา ส่วนกลุ่มที่มีข้อมูลให้ใช้เปิดให้ GA ค้นหาน้ำหนักได้ตามการตั้งค่าของรอบนั้น การเปิดใช้งานในรอบนี้พิจารณาว่ามีหลักฐานด้านนั้นปรากฏใน validation หรือไม่ โดยไม่ได้กำหนดเกณฑ์ขั้นต่ำด้านสัดส่วนความครอบคลุมและจำนวนคู่บวกเพิ่มเติม

ผู้วิจัยตรวจสอบสัดส่วนคู่ที่มีหลักฐานพร้อมใช้งานควบคู่กับน้ำหนักที่ได้ เพื่อแยกกรณี ``ไม่มีข้อมูลให้ใช้'' ออกจากกรณี ``มีข้อมูล แต่ GA เลือกน้ำหนักเป็นศูนย์'' นอกจากนี้ หลักฐานกลุ่มภาพอาจรวมสัญญาณจากคำบรรยายภาพด้วย สัดส่วนความพร้อมของกลุ่มภาพจึงไม่ควรตีความเป็นสัดส่วนคู่ที่มีภาพจริงครบทั้งสองฝั่งโดยอัตโนมัติ

\subsection*{3.6.4 การคำนวณคะแนนเมื่อข้อมูลแต่ละคู่มีไม่ครบ}
โปรไฟล์แต่ละคู่มีข้อมูลไม่เท่ากัน หากแทนข้อมูลที่ขาดด้วยศูนย์แล้วเฉลี่ยทุกด้านเท่ากัน คู่ที่ไม่มีประวัติย่อหรือภาพอาจได้คะแนนต่ำลงเพียงเพราะไม่มีข้อมูล ผู้วิจัยจึงแยกค่าหลักฐานออกจากตัวบ่งชี้ความพร้อมใช้งาน และปรับสัดส่วนน้ำหนักตามข้อมูลที่มีจริงในแต่ละคู่

กำหนดให้ $s_{ik}$ เป็นค่าหลักฐานช่องที่ $k$ ของคู่ที่ $i$, $m_{ik}$ เป็นตัวบ่งชี้ความพร้อมใช้งาน โดยมีค่า 1 เมื่อใช้ได้และ 0 เมื่อใช้ไม่ได้ และ $w_k$ เป็นน้ำหนักที่ GA ค้นหา คะแนนรวมของแต่ละคู่คำนวณได้ดังนี้
\begin{equation*}
S_i=\frac{\sum_{k=1}^{8} w_k m_{ik}s_{ik}}{\sum_{k=1}^{8} w_k m_{ik}},\qquad w_k\geq 0,\quad \sum_{k=1}^{8}w_k=1
\end{equation*}

สูตรนี้ใช้เมื่อผลรวมน้ำหนักของช่องที่พร้อมใช้งานมากกว่าศูนย์ ตัวอย่างเช่น หากนโยบายให้น้ำหนักคะแนนโมเดล 0.2 ชื่อ 0.5 และประวัติย่อ 0.3 แต่คู่หนึ่งไม่มีประวัติย่อ ระบบจะหารด้วย 0.7 ซึ่งเป็นน้ำหนักรวมของสองด้านที่ยังใช้ได้ แทนการหารด้วย 1 จึงไม่ลดคะแนนเพราะการขาดประวัติย่อเพียงอย่างเดียว หากไม่มีหลักฐานที่มีน้ำหนักให้ใช้เลย ระบบกำหนดคู่นั้นเป็น REVIEW

จากนั้นระบบเปรียบเทียบคะแนนรวมกับเกณฑ์สองระดับเพียงครั้งเดียว ดังนี้
\begin{equation*}
D_i=\begin{cases}
\mathrm{MATCH}, & S_i\geq t_{\mathrm{match}},\\
\mathrm{NO\_MATCH}, & S_i<t_{\mathrm{review}},\\
\mathrm{REVIEW}, & t_{\mathrm{review}}\leq S_i<t_{\mathrm{match}}.
\end{cases}
\end{equation*}

MATCH หมายถึงคู่ที่ผ่านเกณฑ์สำหรับการเชื่อมโยงอัตโนมัติ REVIEW หมายถึงคู่ที่ยังต้องตรวจสอบเพิ่มเติม และ NO\_MATCH หมายถึงคู่ที่ไม่ผ่านเกณฑ์การเชื่อมโยงในรอบนั้น คะแนน $S_i$ เป็นคะแนนรวมของนโยบาย จึงไม่ควรเรียกว่าเป็นความน่าจะเป็นที่ผ่านการปรับเทียบแล้ว เว้นแต่มีการประเมินการปรับเทียบของคะแนนรวมนี้เพิ่มเติม

\subsection*{3.6.5 การกำหนดโครโมโซมและเงื่อนไขการค้นหา}
โครโมโซมของนโยบายแบบถ่วงน้ำหนักประกอบด้วย 10 ค่า คือเกณฑ์สองระดับและน้ำหนักอีก 8 ค่า เรียงตามลำดับดังนี้
\begin{equation*}
\boldsymbol{g}=[t_m,t_r,w_p,w_n,w_b,w_u,w_l,w_s,w_i,w_y]
\end{equation*}
โดย $t_m$ และ $t_r$ แทนเกณฑ์ MATCH และ REVIEW ส่วนน้ำหนัก $w_p$ ถึง $w_y$ แทนคะแนนโมเดล ชื่อ ประวัติย่อ ลิงก์ ตำแหน่งที่อยู่ สัญญาณทางสังคม ภาพ และรูปแบบการเขียนตามลำดับ

การค้นหาในโค้ดกำหนดให้ $0.50\leq t_m\leq0.999$ และ $0.05\leq t_r\leq t_m$ น้ำหนักทุกค่าต้องไม่ติดลบ และปรับผลรวมให้เท่ากับ 1 หลังสร้างหรือเปลี่ยนแปลงโครโมโซม กลุ่มข้อมูลที่ไม่ได้เปิดใช้งานมีน้ำหนักเป็นศูนย์ตลอดการค้นหา เงื่อนไขเหล่านี้ช่วยให้โครโมโซมทุกชุดนำไปตัดสินใจได้จริงและคงลำดับของเกณฑ์สองระดับไว้เสมอ

\subsection*{3.6.6 การกำหนดต้นทุนสำหรับประเมินนโยบาย}
ผู้วิจัยประเมินนโยบายด้วยต้นทุนสามประเภท ได้แก่ การรวมโปรไฟล์ของคนละคนเข้าด้วยกัน การไม่รวมคู่ที่เป็นคนเดียวกัน และการส่งคู่ให้เจ้าหน้าที่ตรวจสอบ โดยใช้สมมติฐานค่าแรง 100 บาทต่อชั่วโมง เวลาตรวจสอบคู่ละ 1 นาที เวลาแก้ไขการรวมผิด 20 นาที และเวลาจัดการคู่ที่ตกหล่น 10 นาที สมมติฐานเหล่านี้เป็นค่าที่กำหนดสำหรับการทดลอง ยังไม่ใช่ต้นทุนที่วัดจากการใช้งาน CRM จริง

เมื่อคิดเฉพาะค่าแรง จะได้ต้นทุนประมาณ 33.33 บาทต่อการรวมผิดหนึ่งคู่ 16.67 บาทต่อคู่ที่ตกหล่นหนึ่งคู่ และ 1.67 บาทต่อการตรวจสอบหนึ่งคู่ เมื่อนำต้นทุนคู่ที่ตกหล่นมาเป็นหน่วยอ้างอิง จะได้อัตราส่วน 2 ต่อ 1 ต่อ 0.1 และกำหนดฟังก์ชันต้นทุนดังนี้
\begin{equation*}
C(\boldsymbol{g})=2\,FP+FN+0.1\,R
\end{equation*}

ในสมการนี้ $FP$ คือจำนวนคู่คนละคนที่ถูกตัดสินเป็น MATCH, $FN$ คือจำนวนคู่คนเดียวกันที่ถูกตัดสินเป็น NO\_MATCH และ $R$ คือจำนวนคู่ทั้งหมดที่ถูกส่งเป็น REVIEW การนับ $FN$ ในต้นทุนจึงหมายถึงการปฏิเสธคู่ที่ควรเชื่อมโยงโดยตรง ส่วนคู่บวกที่อยู่ใน REVIEW ถูกคิดเป็นต้นทุนการตรวจสอบ และไม่ได้ถูกคิดซ้ำเป็น $FN$ ในสมการเดียวกัน

GA ค้นหานโยบายที่ทำให้ต้นทุนดังกล่าวต่ำที่สุดบน validation การลดต้นทุนอาจมาพร้อมกับการเปลี่ยนแปลงของ precision, recall และจำนวน REVIEW จึงต้องรายงานตัวชี้วัดเหล่านี้ร่วมกัน ไม่สามารถสรุปจากต้นทุนเพียงค่าเดียวว่านโยบายดีกว่าในทุกด้าน และไม่ควรแปลงผลต้นทุนเป็นเงินที่องค์กรประหยัดได้จริงโดยไม่มีข้อมูลการปฏิบัติงานรองรับ

\subsection*{3.6.7 ขั้นตอนการค้นหานโยบายด้วย GA}
การค้นหาเริ่มจากสร้างประชากรของชุดพารามิเตอร์หลายชุด แต่ละชุดแทนนโยบายที่เป็นไปได้หนึ่งนโยบาย ผู้วิจัยสุ่มค่าเกณฑ์และน้ำหนักเริ่มต้น แล้วปรับให้อยู่ภายใต้เงื่อนไขที่กำหนด จากนั้นนำทุกนโยบายไปตัดสินคู่ใน validation และคำนวณต้นทุนของแต่ละชุด

ในแต่ละรุ่น ระบบเรียงประชากรตามต้นทุนจากน้อยไปมาก สำหรับตัวค้นหานโยบายแบบถ่วงน้ำหนัก หากต้นทุนเท่ากันจะพิจารณา precision ที่สูงกว่า จำนวน FP ที่น้อยกว่า recall ที่สูงกว่า และจำนวน REVIEW ที่น้อยกว่าตามลำดับ แล้วเก็บนโยบายอันดับต้นไว้เป็นกลุ่มชั้นนำเพื่อใช้สร้างรุ่นถัดไป

การสร้างลูกใช้การเลือกยีนแต่ละตำแหน่งจากพ่อหรือแม่ด้วยโอกาสเท่ากัน จากนั้นกลายพันธุ์แต่ละยีนด้วยความน่าจะเป็น 0.30 โดยเพิ่มค่ารบกวนจากการแจกแจงปกติ ส่วนเบี่ยงเบนมาตรฐานสำหรับเกณฑ์เท่ากับ 0.025 และสำหรับน้ำหนักเท่ากับ 0.05 หลังกลายพันธุ์ ระบบตรวจขอบเขตของเกณฑ์ ปิดน้ำหนักของกลุ่มที่ไม่ใช้งาน และปรับผลรวมน้ำหนักอีกครั้ง กระบวนการนี้ทำซ้ำตามจำนวนรุ่นที่กำหนด

การทดลองบนข้อมูลเต็มใช้ประชากร 64 ชุด จำนวน 60 รุ่น เก็บประชากรชั้นนำ 8 ชุด และทำซ้ำด้วยค่า seed 42, 43 และ 44 ทั้งนโยบายคะแนนอย่างเดียว นโยบายใช้ชื่อช่วยตัดสิน และนโยบายแบบถ่วงน้ำหนักใช้ขนาดประชากรและจำนวนรุ่นเท่ากันในรอบนี้ หลังจบการค้นหา ผู้วิจัยเลือกผลของแต่ละนโยบายข้าม seed โดยเรียงตามต้นทุน FP, FN, REVIEW และ seed ตามลำดับ พร้อมเก็บประวัติการค้นหาไว้ตรวจสอบความเปลี่ยนแปลงของคำตอบ

การใช้หลาย seed ช่วยตรวจสอบผลจากการสุ่มเริ่มต้น แต่ไม่ได้รับประกันว่าพบคำตอบที่ดีที่สุดของพื้นที่ค้นหาทั้งหมด นอกจากนี้ ค่าตั้งต้นของการทดลองย่อยอาจต่างจากรอบข้อมูลเต็ม จึงต้องอ้างอิงการตั้งค่าจากไฟล์ของรอบที่นำมารายงานเสมอ

\section*{3.7 การออกแบบการประเมินนโยบายที่พัฒนาเพิ่มเติม}
\subsection*{3.7.1 การแบ่งบทบาทของข้อมูลและการตรึงนโยบาย}
ผู้วิจัยแยกบทบาทของการฝึกแบบจำลอง การเลือกนโยบาย และการทดสอบออกจากกัน โดยใช้คะแนนและการแบ่งกลุ่มที่บันทึกไว้ในชุดข้อมูลของแต่ละการทดลอง การค้นหาพารามิเตอร์ของ GA ใช้เฉพาะคู่ validation ส่วน test ใช้ประเมินนโยบายหลังเลือกเสร็จแล้ว การเปรียบเทียบแต่ละนโยบายจึงใช้คู่ข้อมูล คะแนนโมเดล และนิยามต้นทุนเดียวกัน

สำหรับรอบข้อมูลเต็ม มีคู่ทั้งหมด 2,086,245 คู่ แบ่งเป็นคู่ที่ใช้คะแนนโมเดล 2,073,842 คู่ และคู่ในชั้น Exact-First 12,403 คู่ โดย GA ใช้คู่ validation ที่ผ่านการให้คะแนนโมเดลจำนวน 1,010,926 คู่ในการค้นหา หลังเลือกนโยบาย ผู้วิจัยบันทึกค่าเกณฑ์ น้ำหนัก seed รายชื่อกลุ่มข้อมูลที่เปิดใช้งาน และรหัสตรวจสอบไฟล์ต้นทางไว้ก่อนเรียกขั้นตอนประเมิน test

การตรึงนโยบายก่อนทดสอบช่วยป้องกันการนำผล test ในรอบนั้นกลับไปปรับพารามิเตอร์ อย่างไรก็ตาม การมีไฟล์ตรึงนโยบายไม่ได้ยืนยันด้วยตัวเองว่าคะแนนต้นทางทุกชุดถูกสร้างโดยปราศจากข้อมูลรั่วไหล หรือว่าชุดทดสอบไม่เคยถูกใช้ในงานก่อนหน้า การกล่าวถึงการแยกข้อมูลระดับบุคคลหรือการทดสอบกับชุดใหม่จึงต้องอาศัยหลักฐานของกระบวนการต้นทางเพิ่มเติมด้วย

\subsection*{3.7.2 ฐานเปรียบเทียบและขอบเขตการเปลี่ยนแปลง}
การประเมินใช้สี่แนวทาง ได้แก่ ผลการตัดสินใจจากระบบเดิม นโยบายคะแนนโมเดลอย่างเดียวที่ค้นหาเกณฑ์สองระดับด้วย GA นโยบายที่ใช้ชื่อช่วยตัดสินด้วยพารามิเตอร์สี่ค่า และนโยบายแบบถ่วงน้ำหนักสิบค่า การมีฐานคะแนนอย่างเดียวช่วยแยกผลของการปรับเกณฑ์ออกจากผลของการเพิ่มหลักฐานหลายด้าน ส่วนนโยบายใช้ชื่อช่วยตัดสินทำให้เห็นผลของการเปลี่ยนรูปแบบกฎจากการตัดสินในช่วง REVIEW มาเป็นคะแนนรวมขั้นเดียว

คู่ที่ผ่านชั้น Exact-First ยังคงคำตัดสิน MATCH ตามระบบเดิมทุกแนวทาง จึงไม่ใช่ส่วนที่ GA ปรับปรุง ความผิดพลาดของชั้นดังกล่าวยังต้องนับเมื่อรายงานผลรวม ขณะเดียวกัน ผลบนคู่ที่ใช้คะแนนโมเดลเพียงอย่างเดียวต้องแยกจากผลที่รวม Exact-First แล้ว เพื่อไม่ให้ความแม่นยำของกฎเดิมถูกตีความเป็นผลของ GA ทั้งหมด

\subsection*{3.7.3 ตัวชี้วัดและการอ่านผลลัพธ์สามทาง}
ผู้วิจัยรายงานต้นทุน จำนวน FP จำนวน FN จำนวน REVIEW, precision, recall และ F1 ร่วมกัน โดย precision ของการรวมอัตโนมัติคำนวณจากคู่บวกที่เป็น MATCH หารด้วยคู่ทั้งหมดที่เป็น MATCH ส่วน recall คำนวณจากคู่บวกที่เป็น MATCH หารด้วยคู่บวกทั้งหมดในขอบเขตที่กำลังประเมิน

\begin{equation*}
\mathrm{Precision}=\frac{TP}{TP+FP},\qquad
\mathrm{Recall}=\frac{TP}{N_+},\qquad
F1=\frac{2\,\mathrm{Precision}\,\mathrm{Recall}}{\mathrm{Precision}+\mathrm{Recall}}
\end{equation*}

กำหนดให้ $TP$ เป็นคู่บวกที่ตัดสินเป็น MATCH และ $N_+$ เป็นคู่บวกทั้งหมดในชุดประเมิน คู่บวกที่อยู่ใน REVIEW ยังไม่ถือว่ารวมสำเร็จ จึงยังอยู่ในตัวหารของ recall ด้วย ดังนั้น ค่า FN ที่แสดงในฟังก์ชันต้นทุนไม่สามารถนำไปแทนตัวหารของ recall เป็น $TP+FN$ ได้ทันที เพราะจะขาดคู่บวกที่รอตรวจสอบ หากตัวหารเป็นศูนย์ต้องระบุว่าไม่มีตัวอย่างรองรับการตีความ แทนการมองค่าที่โปรแกรมเติมให้เป็นหลักฐานของประสิทธิภาพ

การรายงานแยกผลบน validation ผลบน test และผลการนำไปใช้กับคู่ทั้งหมด โดยผลรวมทุกคู่ใช้แสดงปริมาณงานที่นโยบายก่อให้เกิดขึ้นในข้อมูลชุดนั้น และไม่ใช้แทนผลทดสอบบนข้อมูลที่กันไว้ กลุ่มที่บันทึกเป็น drop ต้องแสดงแยกตามสถานะของข้อมูลต้นทาง และไม่เรียกว่าเป็นชุดทดสอบเพิ่มเติม

นอกจากนี้ การประเมินบนคู่ผู้สมัครยังไม่ครอบคลุมคู่จริงที่หลุดไปตั้งแต่ขั้นสร้างคู่ ผู้วิจัยจึงอ่านผลนโยบายร่วมกับความครอบคลุมของ candidate generation ที่อธิบายไว้ในหัวข้อ 3.5 การปรับ GA ไม่สามารถกู้คืนคู่ที่ไม่เคยถูกส่งเข้ามาให้ตัดสินได้ หากไม่ได้ปรับขั้นสร้างคู่ด้วย

\subsection*{3.7.4 การตรวจสอบน้ำหนักและข้อมูลที่ขาดหาย}
ผู้วิจัยรายงานน้ำหนักที่ตรึงไว้พร้อมสัดส่วนความพร้อมใช้งานของหลักฐาน เพื่อช่วยอธิบายว่าระบบอาศัยข้อมูลใดในการตัดสินใจ น้ำหนักรวมของนโยบายยังต้องแยกจากน้ำหนักที่มีผลจริงต่อแต่ละคู่ เพราะสูตรคำนวณจะปรับสัดส่วนใหม่เมื่อบางด้านไม่มีข้อมูล ตัวอย่างเช่น น้ำหนักภาพในนโยบายอาจมีค่าน้อย แต่ในคู่ที่ไม่มีหลักฐานชนิดอื่น สัดส่วนที่มีผลจริงอาจแตกต่างจากน้ำหนักรวมมาก

การวิเคราะห์จึงพิจารณาทั้งน้ำหนัก ความครอบคลุม ความเปลี่ยนแปลงระหว่าง seed และข้อผิดพลาดของคู่ที่มีหรือไม่มีหลักฐานแต่ละด้าน การอ้างว่ากลุ่มใดสำคัญที่สุดต่อแบบจำลองยังต้องใช้การวิเคราะห์แบบจำลองโดยตรงเพิ่มเติม ไม่สามารถสรุปจากน้ำหนักนโยบายเพียงอย่างเดียวได้

\section*{3.8 ขั้นตอนการดำเนินงานส่วนที่พัฒนาเพิ่มเติม}
เพื่อให้เห็นความเชื่อมโยงกับงานเดิม ผู้วิจัยจัดลำดับงานเพิ่มเติมตั้งแต่ตรวจสอบข้อมูลต้นทางจนถึงจัดเตรียมผลสำหรับระบบ CRM เป็น 9 ขั้นตอนดังนี้

\subsection*{ขั้นตอนที่ 1 ทบทวนผลและข้อจำกัดของระบบเดิม}
รวบรวมคะแนนรายคู่ คำตัดสินเดิม และผลประเมินของกระบวนการ Exact-First ร่วมกับเกณฑ์สองระดับ ตรวจสอบจำนวนคู่ที่รวมผิด คู่ที่ปฏิเสธผิด และคู่ที่รอตรวจสอบ เพื่อกำหนดว่าการพัฒนาครั้งนี้จะปรับวิธีตัดสินใจในส่วนใด ผลที่ต้องได้คือฐานเปรียบเทียบที่มีขอบเขตและนิยามตัวชี้วัดตรงกัน

\subsection*{ขั้นตอนที่ 2 จัดเตรียมคู่ข้อมูลและตรวจความถูกต้อง}
แยกคู่ Exact-First ออกจากคู่ที่ใช้คะแนนโมเดล ตรวจรหัสคู่ซ้ำ จำนวนแถว และรหัสตรวจสอบไฟล์ต้นทาง จากนั้นจัดคู่ validation และ test ตามข้อมูลการแบ่งกลุ่มที่บันทึกไว้ ขั้นตอนนี้ช่วยให้เปรียบเทียบทุกนโยบายกับประชากรเดียวกัน และติดตามได้ว่าผลที่รายงานมาจากข้อมูลรุ่นใด

\subsection*{ขั้นตอนที่ 3 สร้างหลักฐานสำหรับการตัดสินใจ}
เชื่อมข้อมูลคู่กับคุณลักษณะชื่อ ประวัติย่อ ลิงก์ ตำแหน่งที่อยู่ สัญญาณทางสังคม ภาพ และรูปแบบการเขียนตามที่มีจริง จัดทิศทางคะแนนให้สอดคล้องกัน และสร้างตัวบ่งชี้ความพร้อมใช้งานแยกจากค่าคะแนน สำหรับข้อมูลจำนวนมากให้ประมวลผลเป็นส่วนย่อย และตรวจรหัสคู่ทุกส่วนก่อนรวมผล

\subsection*{ขั้นตอนที่ 4 ตรวจความครอบคลุมและกำหนดกลุ่มที่ใช้งาน}
นับจำนวนและสัดส่วนคู่ที่ใช้หลักฐานแต่ละด้านได้ ตรวจกรณีข้อมูลหายหรือไม่มีช่องข้อมูล และกำหนดรายชื่อกลุ่มที่เปิดให้ GA ค้นหาน้ำหนักจาก validation สำหรับรอบข้อมูลเต็ม ปิด location และ social ตามสภาพข้อมูลจริง พร้อมบันทึกเหตุผล เพื่อไม่ตีความการไม่มีข้อมูลเป็นผลเชิงลบของหลักฐานด้านนั้น

\subsection*{ขั้นตอนที่ 5 กำหนดนโยบายเปรียบเทียบและต้นทุน}
เตรียมนโยบายคะแนนอย่างเดียว นโยบายใช้ชื่อช่วยตัดสิน และนโยบายแบบถ่วงน้ำหนัก กำหนดต้นทุนเป็น $2FP+FN+0.1R$ เหมือนกันทุกแนวทาง และระบุให้ชัดว่าคู่บวกใน REVIEW คิดเป็นงานตรวจสอบ การกำหนดวิธีนับก่อนเริ่มค้นหาช่วยลดความคลาดเคลื่อนในการเปรียบเทียบผล

\subsection*{ขั้นตอนที่ 6 ค้นหาพารามิเตอร์บน validation}
ใช้ GA ค้นหาค่าเกณฑ์และน้ำหนักตามโครงสร้างของแต่ละนโยบาย ทำซ้ำด้วย seed ที่กำหนด และบันทึกต้นทุนกับพารามิเตอร์ระหว่างการค้นหา เปรียบเทียบผลจากการตั้งค่าของรอบเดียวกัน แล้วเลือกชุดที่ดีที่สุดตามลำดับเกณฑ์ที่กำหนดไว้ โดยยังไม่นำผล test มาใช้เลือก

\subsection*{ขั้นตอนที่ 7 บันทึกและตรึงนโยบายก่อนทดสอบ}
จัดเก็บเกณฑ์ น้ำหนัก seed การตั้งค่า GA กลุ่มข้อมูลที่เปิดใช้งาน และรหัสตรวจสอบไฟล์ต้นทางลงในไฟล์นโยบายที่ตรึงไว้ ขั้นตอนนี้ทำให้ย้อนตรวจได้ว่าชุดพารามิเตอร์ที่ใช้ประเมิน test เป็นชุดเดียวกับที่เลือกจาก validation และสามารถนำกลับมาใช้ตัดสินข้อมูลได้ซ้ำ

\subsection*{ขั้นตอนที่ 8 ประเมิน test และนำไปตัดสินคู่ทั้งหมด}
นำพารามิเตอร์ที่ตรึงไว้ไปประเมิน test โดยไม่ปรับค่าเพิ่มเติม จากนั้นคำนวณคำตัดสินสำหรับคู่ทั้งหมดโดยคง Exact-First ตามเดิม และส่งออกผลแยกตามนโยบาย รายงานผล validation, test, drop และภาพรวมแยกกัน เพื่อให้เห็นทั้งคุณภาพในขอบเขตทดสอบและปริมาณงานที่เกิดเมื่อใช้กับข้อมูลทั้งหมด

\subsection*{ขั้นตอนที่ 9 วิเคราะห์ข้อผิดพลาดและเตรียมผลส่งต่อ CRM}
ตรวจคู่ที่ตัดสินต่างจากระบบเดิม โดยเฉพาะคู่ที่เปลี่ยนจาก REVIEW เป็น MATCH หรือ NO\_MATCH เพื่ออธิบายว่าภาระตรวจสอบที่ลดลงสัมพันธ์กับความผิดพลาดอย่างไร จัดเก็บรหัสคู่ คะแนน คำตัดสิน และรุ่นนโยบายสำหรับส่งต่อกระบวนการรวมโปรไฟล์หรือคิวตรวจสอบ ผลส่งออกนี้เป็นหลักฐานของการประมวลผลในโครงงาน ยังไม่ใช่หลักฐานว่าเจ้าหน้าที่ได้ตรวจครบทุกคู่ หรือระบบ CRM ใช้งานจริงและเกิดผลประหยัดตามสมมติฐานแล้ว

\begin{center}\small
\begin{tabularx}{\linewidth}{p{3.4cm}X}
\toprule
ช่วงดำเนินงาน & ผลงานที่จัดเก็บเพื่อตรวจสอบย้อนกลับ \\
\midrule
เตรียมข้อมูล & ตารางหลักฐานรายคู่ ตัวบ่งชี้ความพร้อมใช้งาน และข้อมูลตรวจสอบแหล่งที่มา \\
ค้นหานโยบาย & ผล validation ของแต่ละ seed และประวัติการค้นหา GA \\
ตรึงนโยบาย & เกณฑ์ น้ำหนัก การตั้งค่า และไฟล์นโยบายก่อนประเมิน test \\
ประเมินและส่งออก & ตารางตัวชี้วัดแยกกลุ่มข้อมูล และคะแนนกับคำตัดสินรายคู่ของแต่ละนโยบาย \\
\bottomrule
\end{tabularx}
\par ตารางที่ 3.20 ผลงานจากขั้นตอนการพัฒนานโยบายเพิ่มเติม
\end{center}

การดำเนินงานส่วนนี้ทำให้กระบวนการเดิมมีขั้นตอนเลือกนโยบายที่ตรวจสอบย้อนกลับได้ ตั้งแต่หลักฐานที่ใช้ ค่าพารามิเตอร์ที่เลือก ไปจนถึงผลกระทบต่อการรวมข้อมูลและคิวตรวจสอบ โดยยังรักษาความเชื่อมโยงกับการเตรียมข้อมูล การสร้างคู่ และการให้คะแนนที่ได้พัฒนาไว้ก่อนหน้า

\clearpage
\chapter*{บทที่ 4}
\begin{center}\Large\bfseries ผลการทดลอง\end{center}
\section*{4.1 การเตรียมข้อมูล}
ชุดข้อมูลที่ใช้ในการพัฒนาและทดสอบระบบครั้งนี้เป็นชุดข้อมูลจำลอง (Synthetic Dataset) ที่สร้างขึ้นโดยอ้างอิงจากเทมเพลตโครงสร้างข้อมูลที่บริษัทจัดให้ ครอบคลุมโปรไฟล์ผู้ใช้ จาก 3 แพลตฟอร์ม ได้แก่ Twitter, Instagram และ Google+ Archive รวมทั้งสิ้น 36,807 โปร ไฟล์ โดยแต่ละโปรไฟล์ประกอบด้วยข้อมูลชื่อผู้ใช้ (Username), ชื่อจริง (Full Name), ประวัติย่อ (Bio), ลิงก์อ้างอิง (URL) และตำแหน่งที่อยู่ (Location) ซึ่งมีระดับความสมบูรณ์ของข้อมูล ที่ แตกต่างกันไปตามแต่ละแพลตฟอร์ม\par
แพลตฟอร์ม จ ำนวนโปรไฟล์ สัดส่วน (\%) Twitter 15,234 41.4\% Instagram 12,871 35.0\% Google+ 8,702 23.6\% รวม 36,807 100\%\par
ก่อนที่โปรไฟล์ทั้งหมดจะเข้าสู่กระบวนการ Classification จำเป็นต้องผ่าน ขั้นตอน Blocking ก่อน เพื่อลดพื้นที่การค้นหาจากจำนวน Candidate Pairs ทั้งหมด C(36,807, 2) \ensuremath{\approx} 677 ล้านคู่ ให้เหลือเฉพาะคู่ที่มีโอกาสเป็นบุคคลเดียวกัน โดยใช้ Blocking Strategy แบบผสมผสานระหว่าง Username Prefix Blocking และ Full Name Jaro-Winkler Similarity ซึ่งสามารถลด Candidate Space ลงเหลือ 873,956 คู่ คิดเป็น Reduction Ratio 99.9\% ในขณะที่ยังคงรักษา Blocking Recall ไว้ใน ระดับสูง จากนั้นได้ทำการติด Label ด้วยกระบวนการ Human-in-the-Loop ได้ชุด ข้อมูลที่มี Label ทั้งสิ้น 162,442 pairs ก่อนแบ่งออกเป็น Train/Val/Test ตามแนวทาง Entity-Aware Split เพื่อป้องกัน Data Leakage ดังแสดงในตารางที่ 4.2\par
\begin{center}\small ตารางที่ 4.2 Entity-Aware Train/Val/Test Split (70/15/15)\end{center}
Split จ ำนวน Entities\par
จ ำนวน Pairs\par
Positive Negative\par
Train 8,472 114,007 16,258 97,749 Val 1,815 24,113 3,510 20,603 Test 1,816 24,322 3,438 20,884\par
% ---- source PDF page 71 ----
\clearpage
รวม 12,103 162,442 23,206 139,236\par
รายการ ค่าที่ได้ All cross-platform pairs 449,149,239 Exact-first auto-accepted pairs 12,403 Candidate pairs for model 2,073,842 Search-space reduction 99.54\% Ground-truth coverage 88.67\%\par
ข้อมูล Google+ มีลักษณะพิเศษที่กว่าอีกสองแพลตฟอร์ม เนื่องจาก Google+ ปิดตัวลง ในปี 2019 ข้อมูลที่เหลืออยู่จึงอยู่ในรูปแบบ Archive ที่ไม่สมบูรณ์ ทั้งในแง่ Metadata และข้อมูล Engagement บางส่วนที่หายไปพร้อมกับการปิดระบบ สถานการณ์นี้เป็นตัวแทนของปัญหาข้อมูล ไม่ครบในสภาพแวดล้อมจริง ไม่ใช่เพียงข้อจำกัด ของงานวิจัยนี้เพียงอย่างเดียว จำนวนโปรไฟล์ แยกตามแพลตฟอร์มแสดงไว้ในตารางที่ 4.1\par
\begin{center}\small ตารางที่ 4.1 จำนวน valid profiles แยกตามแพลตฟอร์มที่ใช้จริง\end{center}
แพลตฟอร์ม จำนวนโปรไฟล์ สัดส่วน (\%) Twitter 13,959 37.9 Instagram 10,956 29.8 Google+ 11,889 32.3 รวม 36,804 100.0\par
ก่อนที่โปรไฟล์ทั้งหมดจะเข้าสู่ขั้น candidate scoring ระบบต้องผ่านกระบวนการลด พื้นที่ค้นหาก่อน เพราะหากเปรียบเทียบทุกคู่แบบ Cross-platform โดยตรงจะมีจำนวนคู่สูงถึง 449,149,239 คู่ การลด Search space จึงเป็นเงื่อนไขจำเป็นของ ไม่ใช่เพียงขั้นตอนเสริม โดย Pipeline ใช้ Exact-First ร่วมกับ Blocking เพื่อดึงคู่ที่หลักฐานชัดมากออกมาจับคู่กันอัตโนมัติ ทันที และทำการส่งคู่ที่ยังมีความกำกวมเข้าสู่ขั้น Candidate scoring\par
\begin{center}\small ตารางที่ 4.2 ผลของขั้นเตรียม candidate set ก่อนเข้าสู่ candidate scoring\end{center}
สำหรับชุดข้อมูลที่ใช้พัฒนาแบบจำลองการให้คะแนน (Candidate-scoring Model) ระบบได้ใช้เทคนิคการแบ่งข้อมูลแบบป้องกันการรั่วไหล (Leak-safe Split) โดยมีเงื่อนไขว่าต้องไม่\par
% ---- source PDF page 72 ----
\clearpage
มีการทับซ้อนของข้อมูลระหว่างชุด Train, Validation และ Test รวมถึงต้องไม่มีคู่โปรไฟล์ที่จับคู่ กับตัวเองหลงเหลืออยู่ การแบ่งข้อมูลด้วยวิธีนี้ช่วยป้องกันปัญหาผลการประเมินสูงเกินจริง ซึ่งเป็น ปัญหาที่พบได้บ่อยในงานระบุตัวตน เมื่อข้อมูลของบุคคลเดียวกันรั่วไหลไปปรากฏอยู่ในหลายชุด ข้อมูลพร้อมกัน\par
\begin{center}\small ตารางที่ 4.3 โครงสร้างของ supervised split สำหรับ main run\end{center}
Split Profiles Positive\par
Random Negatives\par
Hard Negatives Rows รวม\par
Pairs\par
Train 22,478 20,470 15,346 5,987 41,803 Validation 7,925 4,386 3,289 905 8,580 Test 6,401 4,387 3,290 636 8,313\par
แสดงให้เห็นว่าการใคก Positive pairs เพียงอย่างเดียว แต่ใช้ทั้ง Random negatives` และ `hard negatives` เพื่อทำให้ model เห็นกรณีที่แยกได้ง่ายและกรณี ที่สับสนคล้ายกันมากในเวลาเดียวกัน การออกแบบนี้ส่งผลโดยตรงต่อคุณภาพของ calibrated candidate scoring ในขั้นถัดไป\par
\begin{center}\small รูปที่ 4.1 สรุปผลการเตรียมข้อมูล ได้แก่ จำนวนโปรไฟล์ที่ใช้งานได้แยกตามแพลตฟอร์ม (รวม 36,604 โปรไฟล์), การแบ่ง Train/Val/Test แบบ Leak-Safe และการเติบโตของ Feature Space ในแต่ละชุด\end{center}
Feature ที่ออกแบบไว้\par
การแบ่งข้อมูลในลักษณะ Entity-Aware คือการแยก Entities ออกก่อน แล้วให้ Pairs ที่ เกี่ยวข้องติดตามไปยัง Split เดียวกัน ซึ่งแตกต่างจาก Random Split ตรงที่ Entity เดียวกันจะ ไม่กระจายอยู่ทั้งใน Train และ Test พร้อมกัน การออกแบบนี้ช่วยป้องกันปัญหา Optimistic Bias ที่มักเกิดขึ้นในงาน Entity Matching เมื่อใช้การแบ่งแบบสุ่มทั่วไป\par
% ---- source PDF page 73 ----
\clearpage
\begin{center}\small รูปที่ 4.2 ตัวอย่างโปรไฟล์จาก 3 แพลตฟอร์มหลังผ่าน Standardization\end{center}
\section*{4.2 วิธีวัดประสิทธิภาพ}
แนวทางการวัดผลในงานวิจัยครั้งนี้แบ่งออกเป็น 3 ระดับตามลำดับของ pipeline จริง ได้แก่ ระดับการลดพื้นที่ค้นหา, ระดับ candidate scoring และระดับ decision thresholding สำหรับการใช้งานจริงใน CRM การแยกตัวชี้วัด ออกตามหน้าที่ของแต่ละขั้นมีความสำคัญ เพราะค่า metric ที่ดีในระดับหนึ่ง ไม่ได้แปลว่าระบบพร้อมใช้งานจริงโดยอัตโนมัติเสมอไป\par
\subsection*{4.2.1 ตัวชี้วัดสำหรับ Search Space Reduction และ Candidate Retrieval}
ในขั้นก่อนเข้าสู่ candidate scoring ผู้วิจัยใช้ตัวชี้วัดหลัก 2 ค่า • Reduction Ratio สะท้อนว่าระบบลดจำนวนคู่ที่ต้องตรวจต่อได้มากเพียงใด ยิ่งสูงยิ่งดี\par
เพราะหมายถึง pipeline ประหยัดภาระการคำนวณในขั้น scoring ได้มาก ชึ้น\par
Reduction Ratio = 1 −( Pairs after\par
)\par
Pairs before\par
\noindent • Ground-truth Coverage สะท้อนถึงสัดส่วนการคงอยู่ของคู่ข้อมูลจริงที่ต้องนำไป\par
ประมวลผลต่อในขั้นถัดไป หากค่าความครอบคลุม มีระดับที่ต่ำจนเกินไป ตัวแบบจะสูญเสีย\par
โอกาสในการกู้คืนคู่ข้อมูลจริงที่ถูกคัดทิ้งไปในขั้นตอนก่อนหน้า ซึ่งวัดสัดส่วนของ True\par
Positive Pairs ที่สามารถผ่าน Blocking เข้ามาได้เทียบกับ True Positive Pairs ทั้งหมด\par
ค่านี้ยิ่งสูงยิ่งดี เนื่องจากหาก Blocking ตัด Positive Pairs ออกไป\par
Ground −truth Coverage = TP passed\par
TP total\par
% ---- source PDF page 74 ----
\clearpage
\subsection*{4.2.2 ตัวชี้วัดสำหรับ Candidate Scoring}
ในระดับ candidate scoring ใช้ตัวชี้วัดมาตรฐานของงานจัดอันดับและจำแนก ได้แก่\par
\noindent • Precision วัดว่าจากทั้งหมดที่แบบจำลองทำนายว่า MATCH มีกี่เปอร์เซ็นต์ที่เป็น\par
True Positive จริง\par
Precision = TP (TP + FP)\par
\noindent • Recall ซึ่งวัดว่าจาก Positive Pairs ทั้งหมด แบบจำลองสามารถจับได้กี่เปอร์เซ็นต์\par
Recall = TP (TP + FN)\par
\noindent • F1-score ค่าเฉลี่ยฮาร์โมนิกของ Precision และ Recall เหมาะสำหรับข้อมูลที่มี\par
Class Imbalance เพราะให้น้ำหนักทั้งสองค่าเท่ากัน\par
F1 −Score = 2 × (Precision × Recall)\par
(Precision + Recall)\par
\noindent • Average Precision (AP) เป็นตัวชี้วัดที่ใช้ประเมินคุณภาพการจัดอันดับของ\par
แบบจำลองบนข้อมูลที่มีความไม่สมดุลของคลาส โดยคำนวณจากพื้นที่ใต้เส้นโค้ง\par
Precision-Recall ซึ่งสะท้อนว่าแบบจำลองสามารถจัดคู่ข้อมูลที่เป็นบวกจริงให้อยู่ใน\par
ลำดับต้นได้ดีเพียงใด ทั้งนี้สามารถเขียนได้ในรูป\par
AP = \ensuremath{\sum}(R n − R n−1 )P n\par
n\par
โดยที่\par
P n คือค่า Precision ที่จุดลำดับที่ n R n คือค่า Recall ที่จุดลำดับที่ n\par
\noindent • ROC-AUC ใช้วัดความสามารถของแบบจำลองในการแยกคู่ MATCH ออกจาก\par
NO\_MATCH โดยไม่ยึดกับ threshold ใด threshold หนึ่ง เส้น ROC สร้างจาก\par
ความสัมพันธ์ระหว่าง True Positive Rate (TPR) และ False Positive Rate\par
(FPR) และค่า AUC คือพื้นที่ใต้เส้นโค้งดังกล่าว เข้าใกล้ 1.0 หมายถึงแบบจำลองมี\par
ความสามารถในการแยกแยะสูงมาก\par
% ---- source PDF page 75 ----
\clearpage
1\par
AUC = \ensuremath{\int}TPR(FPR)d(FPR)\par
0\par
นอกจากตัวชี้วัดหลักดังกล่าว ยังมีการใช้ Confusion Matrix เพื่อวิเคราะห์ลักษณะความ\par
ผิดพลาดของแบบจำลอง และใช้ precision@k เพื่อประเมินคุณภาพของการจัดอันดับคู่ข้อมูลที่มี\par
คะแนนสูงสุดบน full candidate pairs ซึ่งมีความสำคัญต่อการนำผลลัพธ์ไปใช้จัดลำดับงานตรวจ\par
สอบในขั้น review\par
\subsection*{4.2.3 ตัวชี้วัดสำหรับ Probability Calibration และ Decision Thresholding}
เมื่อนำคะแนนจากแบบจำลองไปใช้งานจริงในระดับ Production เป้าหมายของการประเมิน\par
จะเปลี่ยนจากการวัดความแม่นยำเชิงสถิติทั่วไป ไปสู่ ตัวชี้วัดเชิงการตัดสินใจ (Decision-centric\par
Metrics) เพื่อตอบโจทย์กระบวนการทำงานของระบบ CRM โดยให้ความสำคัญกับตัวชี้วัดหลัก 3\par
ส่วน ดังนี้\par
Expected Calibration Error (ECE) ใช้ตรวจสอบว่าคะแนนความน่าจะเป็นที่ผ่าน\par
การปรับเทียบ (Calibration) แล้ว มีความสอดคล้องกับสัดส่วนความเป็นจริงมากน้อยเพียงใด\par
ค่า ECE ที่ต่ำจะเป็นเครื่องยืนยันว่า การตั้งค่าจุดตัด (Threshold) ในระดับสูง (เช่น 0.95\par
หรือ 0.98) มีความน่าเชื่อถือ และสามารถนำไปใช้อ้างอิงเป็นเกณฑ์ในการปฏิบัติงานได้จริง\par
\noindent •\par
ECE = m = 1\ensuremath{\sum}m n \ensuremath{\mid}B m \ensuremath{\mid}acc(B m ) −conf(B m )\par
โดยที่\par
B m คือกลุ่มข้อมูลใน bin ที่ m \ensuremath{\mid}B m \ensuremath{\mid} คือจำนวนตัวอย่างใน bin นั้น n คือจำนวนตัวอย่างทั้งหมด acc(B m ) คือ accuracy จริงของ bin ที่ m conf(B m ) คือค่า confidence เฉลี่ยของ bin ที่ m\par
% ---- source PDF page 76 ----
\clearpage
\noindent • Final Match-only Precision และ Final Match-only Recall เป็นตัวชี้วัดที่ใช้ตอบ\par
โจทย์ปลายทางโดยตรงว่า จุดปฏิบัติการ (Operating Point) หรือค่าจุดตัดที่ระบบ\par
เลือกใช้นั้น มีความแม่นยำและปลอดภัยเพียงพอสำหรับการทำ "การรวมโปรไฟล์อัตโนมัติ\par
(Auto-merge)" ในระบบ CRM หรือไม่ โดยไม่ก่อให้เกิดความผิดพลาดที่ยอมรับไม่ได้\par
Final Match - only Recall = TP match\par
TP all\par
Final Match −only Precision = TP match + FP match TP match โดยที่\par
\noindent • TP match คือจำนวนคู่จริงที่ระบบจัดอยู่ในกลุ่ม MATCH\par
\noindent • FP match คือจำนวนคู่ผิดที่ระบบจัดอยู่ในกลุ่ม MATCH\par
\noindent • TP all คือจำนวนคู่จริงทั้งหมดในข้อมูลอ้างอิง\par
\noindent • Operational Outcomes และ Review Queue Size พิจารณาจากสัดส่วนของผลลัพธ์\par
ที่ถูกแบ่งออกเป็น 3 กลุ่มปฏิบัติการ (MATCH, REVIEW, NO\_MATCH) รวมถึง ขนาด\par
ของคิวที่พนักงานต้องตรวจสอบด้วยตาเปล่า (Review Queue Size) เพื่อประเมินว่า\par
ระบบสามารถช่วยลดภาระงานของมนุษย์ได้จริง และไม่สร้างคอขวด (Bottleneck) ใน\par
กระบวนการทำงาน\par
Review Queue Size = N(T review \ensuremath{\leq}score < T match )\par
Match Count = N(score > T match ) + N (exact match)\par
NO\_\_MATCH Count = N(score < T review ) โดยที่\par
\noindent • T match คือ threshold สำหรับ MATCH\par
\noindent • T review คือ threshold ต่ำสุดของช่วง REVIEW\par
\section*{4.3 Candidate-Scoring Model ที่ใช้ก่อน Decision Thresholding}
main run ที่ถูกใช้เป็นฐานของ Project 1 คือ image\_context\_r075\_h20\_s42 จากชุดทดลอง multimodal suite ซึ่งสืบทอด pair construction, leakage-safe split และ calibration logic มาจาก rebuilt hybrid pipeline เดิม ดังนั้นผลในบทนี้จึงไม่ใช่ผลจาก pair source ที่รั่วหรือ split ที่มี entity overlap ระหว่าง train, validation และ test\par
% ---- source PDF page 77 ----
\clearpage
แบบจำลองในงานนี้ไม่ได้รับ raw profiles เข้าไปโดยตรง แต่แปลงโปรไฟล์สองฝั่งให้เป็น pair-level feature vector ก่อนเสมอ โดยโครงสร้างของ feature space เติบโตเป็นลำดับจาก 22 baseline text+attribute features ไปเป็น 23 features เมื่อเพิ่ม SBERT, ขยายเป็น 37 features เมื่อเพิ่ม image statistics และสิ้นสุดที่ 41 features เมื่อเพิ่ม caption-to-text cross signals ใน run แบบ image\_context วิธีออกแบบเช่นนี้ทำให้ระบบทำงานในฐานะ tabular candidate scorer ที่เรียนรู้จากหลักฐานเชิงความคล้ายของคู่โปรไฟล์หลายประเภทพร้อมกัน\par
\begin{center}\small รูปที่ 4.3 โครงสร้างของ main-run model pipeline สำหรับ Project 1 เริ่มจากการสร้าง pair\end{center}
features จากข้อมูลชื่อ ข้อความ attribute ภาพ และ caption จากนั้นให้ candidate-scoring\par
model ประเมินคะแนน ก่อนผ่าน probability calibration และ decision thresholding\par
จากรูปที่ 4.2 จะเห็นว่าตัวแบบในงานนี้ไม่ใช่ end-to-end deep model ที่เรียนรู้จากข้อมูลดิบ ทั้งก้อน แต่เป็น pipeline เชิงวิศวกรรมข้อมูลที่แยกขั้นชัดเจน กล่าวคือผู้วิจัยสร้าง feature table ระดับคู่ก่อน แล้วจึงให้ model family ที่เหมาะกับข้อมูลเชิงตารางเรียนรู้รูปแบบการแยก MATCH / NO\_MATCH วิธีนี้เหมาะกับขอบเขตของ Project 1 เพราะอธิบายได้ชัด ตรวจสอบได้ง่าย และรองรับการนำคะแนนไป calibrate และ threshold ต่อในเชิงปฏิบัติการ\par
\subsection*{4.3.1 การเลือกโมเดลและการเปรียบเทียบ candidate-scoring configurations}
ภายใน main run ระบบฝึกและเปรียบเทียบอย่างน้อย Gradient Boosting, Random Forest และ Logistic Regression แล้วเลือกโมเดลที่ให้ผลดีที่สุดบน validation set ซึ่ง experiment\_report.json ระบุชัดว่า best\_model = gb หรือ Gradient Boosting ตารางที่ 4.4 ผล validation leaderboard ภายใน main run\par
Model Val AP Val ROC-AUC ข้อสรุปใน run Gradient Boosting 0.9795 0.9720 ดีที่สุด Random Forest 0.9778 0.9696 รองลงมา Logistic Regression 0.9431 0.9245 ต่ำกว่าสองตัวแรกชัดเจน\par
% ---- source PDF page 78 ----
\clearpage
เพื่อไม่ให้การเลือก Gradient Boosting อิงเพียง leaderboard ภายใน run เดียว ผู้วิจัยยัง benchmark model families หลายชนิดบน final 41-feature split เดียวกันเพิ่มเติม ซึ่งรวมถึง MLP, Extra Trees, AdaBoost, Logistic Regression และ Linear SVM ด้วย ผลในชุด benchmark นี้ชี้ว่า Gradient Boosting ให้สมดุลของ AP, ROC-AUC, F1 และ Precision ดีที่สุดในภาพรวม แม้บางโมเดลจะมีค่าในบาง metric ใกล้เคียงกัน\par
\begin{center}\small รูปที่ 4.4 การเปรียบเทียบ model families บน final 41-feature split โดยใช้สีแสดงคุณภาพสัมพัทธ์ภายในแต่ละ metric และใช้ตัวเลขในเซลล์แสดงค่าจริงของ AP ROC-\end{center}
AUC F1 Precision และ Recall พร้อมกันนั้น ผู้วิจัยยังเปรียบเทียบ candidate-scoring configurations ที่ใช้ pair set และ split เดียวกัน เพื่อพิสูจน์ว่า image\_context เหมาะสมกว่าการหยุดที่ text\_attr\_hybrid หรือ image\_stats ตารางที่ 4.5 ผลเปรียบเทียบ candidate-scoring configurations ที่ใช้เลือก main run\par
Configuration Feature\par
Count AP ROC-\par
AUC F1 Precision Recall\par
Text/Attribute Hybrid 23 0.9782 0.9715 0.9330 0.9538 0.9132\par
Image Statistics 37 0.9789 0.9731 0.9340 0.9537 0.9152\par
Image Context 41 0.9789 0.9734 0.9340 0.9507 0.9179\par
% ---- source PDF page 79 ----
\clearpage
\begin{center}\small รูปที่ 4.5 การเปรียบเทียบ metric ของ candidate-scoring configurations โดยใช้สีแสดงคุณภาพสัมพัทธ์ภายในแต่ละ metric และใช้ตัวเลขในเซลล์แสดงค่าจริงของ AP ROC-\end{center}
AUC F1 Precision และ Recall\par
ผลในตารางที่ 4.5 และรูปที่ 4.4 แสดงว่า image\_context เป็น configuration ที่ เหมาะสมที่สุดสำหรับ Project 1 เพราะให้ AP, ROC-AUC และ Recall สูงที่สุดในภาพรวม แม้ Image Statistics จะมี F1 สูงกว่าในหลักทศนิยมลึกบางจุด แต่ความต่างนั้นน้อยมากเมื่อเทียบกับ ข้อได้เปรียบของ Image Context ในด้าน ranking quality และความครอบคลุมของคู่จริง\par
เพื่อให้การเปรียบเทียบครบในระดับโครงสร้างความผิดพลาด ผู้วิจัยจึงพิจารณา confusion matrix ของทั้งสาม configuration ควบคู่ไปด้วย เพราะค่า AP หรือ F1 เพียงอย่าง เดียวอาจยังไม่อธิบายได้ชัดว่าความต่างของแต่ละ run เกิดจากการลด false positives หรือการ ลด false negatives เป็นหลัก\par
\begin{center}\small ตารางที่ 4.6 Confusion matrix เปรียบเทียบ candidate-scoring configurations บน test split\end{center}
Configuration Threshold TN FP FN TP Text/Attribute Hybrid 0.44 3,732 194 381 4,006\par
Image Statistics 0.43 3,731 195 372 4,015 Image Context 0.35 3,717 209 360 4,027\par
% ---- source PDF page 80 ----
\clearpage
\begin{center}\small รูปที่ 4.6 Confusion matrices ของ candidate-scoring configurations ทั้งสามแบบบน test split\end{center}
แสดง trade-off ระหว่าง false positives และ false negatives ที่ใช้ประกอบการเลือก main run\par
ผลในตารางที่ 4.6 และรูปที่ 4.5 ชี้ให้เห็นลักษณะ trade-off ที่สำคัญ 2 ประเด็น ประเด็นแรก Image Statistics ลด FN ลงจาก 381 เหลือ 372 เมื่อเทียบกับ Text/Attribute Hybrid โดยแทบไม่เพิ่ม FP เลย\par
ซึ่งยืนยันว่าการเพิ่มสัญญาณภาพระดับสถิติมีผลเชิงบวกจริงต่อการกู้คู่จริง ประเด็นที่สอง Image Context ลด FN ต่อเนื่องลงเหลือ 360 แต่ยอมให้ FP เพิ่มเป็น 209 ภายใต้ threshold ที่ต่ำกว่า (0.35) ผลเช่นนี้สอดคล้องกับค่า Recall ที่สูงที่สุดของ run นี้ และเป็นเหตุผลสำคัญที่ผู้วิจัยเลือก image\_context เป็น main run แล้วนำไป calibrate และตั้ง t h r e s h o l d เชิง p r o d u c t i o n ต่อ แทนที่จะเลือก c o n fi g u r a t ion ที่เข้มงวดกว่าแต่กู้คู่จริงได้น้อยกว่า\par
\subsection*{4.3.2 ผลเชิงจำแนกของ Main Run บน Test Split}
เมื่อใช้ threshold เชิงจำแนก 0.35 บน test split main run ให้ Precision = 0.9507, Recall = 0.9179 และ F1 = 0.9340 พร้อม confusion matrix TN = 3,717, FP = 209, FN = 360 และ TP = 4,027 ตารางที่ 4.7 ผลของ main run บน test split ก่อนเข้าสู่ production thresholding\par
รายการ ค่าที่ได้ Best model Gradient Boosting Feature count 41 Threshold สำหรับ test classification 0.35 Test AP 0.9789 Test ROC-AUC 0.9734 Test F1 0.9340\par
% ---- source PDF page 81 ----
\clearpage
Test Precision 0.9507 Test Recall 0.9179 TN / FP / FN / TP 3,717 / 209 / 360 / 4,027\par
\begin{center}\small รูปที่ 4.7 Confusion matrix ของ main run บน test split ที่ threshold = 0.35 แสดงสมดุลระหว่าง true positives false positives true negatives และ false negatives\end{center}
ของแบบจำลองก่อนเข้าสู่ขั้น decision thresholding ระดับ production ผลนี้ชี้ว่า main run สามารถรักษาสมดุลระหว่างความเข้มงวดและความครอบคลุมได้ดี กล่าวคือจำนวน false positives อยู่ในระดับต่ำเมื่อเทียบกับ true positives และจำนวน false negatives ก็ไม่ได้สูงเกินไปจนทำให้ recall ตกลงอย่างมีนัยสำคัญ หากมองในมุมของ pipeline ทั้งระบบ ผลในระดับนี้เหมาะกับการส่งต่อเข้าสู่ขั้น calibration และ thresholding มากกว่าการใช้ตัดสินใจ merge โดยตรง\par
\subsection*{4.3.3 การกระจายคะแนนและ Threshold Trade-off}
เมื่อพิจารณาการกระจายของ calibrated probabilities บน test split จะเห็นว่าคู่ NO\_MATCH ส่วนใหญ่กระจุกตัวอยู่ใกล้ศูนย์ ขณะที่คู่ MATCH กระจุกตัวอยู่ใกล้หนึ่งอย่างชัดเจน โดยค่าเฉลี่ยของกลุ่ม NO\_MATCH อยู่ที่ประมาณ 0.1095 และค่ามัธยฐานอยู่ที่ 0.0684 ส่วนกลุ่ม MATCH มีค่าเฉลี่ย 0.9012 และค่ามัธยฐานสูงถึง 0.9966 การแยกตัวของสอง distribution นี้เป็นสัญญาณว่าตัวแบบมีความสามารถในการจัดลำดับคู่ที่น่าใช่และไม่น่าใช่ออกจากกันได้ดี\par
% ---- source PDF page 82 ----
\clearpage
\begin{center}\small รูปที่ 4.8 การกระจายของ calibrated scores บน test split แยกตาม Actual MATCH และ Actual\end{center}
NO\_MATCH พร้อมเส้นอ้างอิงที่ threshold = 0.35 และเส้น production thresholds ที่ 0.95 และ\par
\noindent 0.98\par
\begin{center}\small รูปที่ 4.7 ช่วยให้เห็นชัดว่า threshold สำหรับการประเมินเชิงจำแนก (0.35) และ threshold สำหรับการใช้งานจริง (0.95 / 0.98) อยู่คนละบริบทกัน โดย threshold 0.35 อยู่ ในช่วงที่ optimize ความสมดุลของ F1 บน test split ขณะที่ threshold ปลายทางถูกเลื่อนไป ยังปลายขวาของ distribution เพื่อคุม precision สำหรับการ merge จริงให้เข้มงวดขึ้นมาก\end{center}
หากพิจารณาความสัมพันธ์ระหว่าง Precision, Recall และ F1 เมื่อขยับ threshold จะ เห็น trade-off ชัดเจน กล่าวคือ threshold ต่ำช่วยให้ recall สูง แต่ precision ลดลงเร็ว ขณะที่ threshold สูงช่วยคุม precision ได้ดีขึ้นแลกกับ recall ที่ลดลง ซึ่งเป็นเหตุผลสำคัญที่ทำให้ production operating point ต้องแยก MATCH ออกจาก REVIEW แทนการใช้ threshold เดียวบังคับทั้งระบบ\par
\begin{center}\small รูปที่ 4.9 พฤติกรรมของ Precision Recall และ F1 ของ main run เมื่อเปลี่ยน threshold บน test split แสดงความแตกต่างระหว่าง threshold เชิงจำแนกที่ 0.35 กับ production review threshold\end{center}
ที่ 0.95\par
% ---- source PDF page 83 ----
\clearpage
ที่ threshold 0.35 ระบบได้ Precision = 0.9507, Recall = 0.9179 และ F1 = 0.9340 ซึ่งเหมาะสำหรับการประเมินเชิงแบบจำลอง แต่เมื่อยก threshold ไปถึง 0.95 ค่า precision บน test split จะสูงขึ้นเป็นประมาณ 0.9936 ขณะที่ recall ลดลงเหลือประมาณ 0.8110 และ F1 ลดลงมาอยู่ที่ประมาณ 0.8931 ผลนี้อธิบายได้ชัดเจนว่าทำไม production จึงต้องเพิ่มชั้น REVIEW เข้ามาช่วยรับคู่ที่ยังมีศักยภาพแต่ไม่ควรถูก auto-merge ทันที\par
\subsection*{4.3.4 การสอบเทียบความน่าจะเป็นด้วย Isotonic Regression}
หลังจากเลือกโมเดลที่ดีที่สุดแล้ว pipeline ไม่ได้นำค่าความน่าจะเป็นดิบจาก predict\_proba() ไปใช้ตัดสินใจใน p r o d u c t i o n ทันที แต่ใช้ IsotonicRegression(out\_of\_bounds="clip") เป็นตัว calibrate คะแนนก่อนเสมอ โดยขั้นตอนจริงคือ fit isotonic calibrator บน validation raw scores กับ validation labels จากนั้นจึงนำ calibrator เดียวกันไปแปลงคะแนนของ train, validation และ test ให้กลายเป็น calibrated probabilities ก่อนเลือก operating point ระดับ production\par
ผลจาก score files ของ main run แสดงว่า isotonic calibration ทำงานได้ดีมาก โดย validation set ซึ่งเป็นชุดที่ใช้ fit calibrator มี ECE = 0.0000 และ test set มี ECE = 0.0044 เท่านั้น ค่าที่ต่ำมากนี้สะท้อนว่า calibrated probabilities บน test ยังคงเกาะใกล้ observed positive rate ในแต่ละช่วงคะแนนอยู่มาก\par
\begin{center}\small รูปที่ 4.10 Reliability ของ calibrated probabilities หลังผ่าน Isotonic Regression เปรียบเทียบ\end{center}
validation set และ test set โดยใช้ขนาด marker แทนจำนวนคู่ในแต่ละ probability bin\par
และเส้นตั้งแสดง production thresholds ที่ 0.95 และ 0.98\par
% ---- source PDF page 84 ----
\clearpage
\begin{center}\small รูปที่ 4.9 ให้ข้อสรุปสำคัญว่า calibrated score ของ main run มีความน่าเชื่อถือ เพียงพอสำหรับการกำหนด operating point เชิง production กล่าวคือ validation reliability curve เกือบซ้อนทับเส้นทแยงตามธรรมชาติของ isotonic fit ขณะที่ test curve ยังเกาะใกล้ แนวอุดมคติอยู่มาก แม้บาง bin ขนาดเล็กจะมีความผันผวนตามจำนวนตัวอย่างที่น้อย\end{center}
\subsection*{4.3.5 คุณภาพการจัดอันดับบน Full Candidate Pairs}
นอกเหนือจากผลบน test split ผู้วิจัยยังประเมินคุณภาพการจัดอันดับของ calibrated candidate scores บน candidate pairs จริงทั้งหมด 2,073,842 คู่ที่ผ่าน blocking เข้ามา ผลลัพธ์แสดงว่า candidate\_avg\_precision = 0.6017 และ candidate\_roc\_auc = 0.9600 พร้อมค่า precision@100 = 0.9800, precision@500 = 0.9800, precision@1000 = 0.9820 และ precision@5000 = 0.9472\par
\begin{center}\small รูปที่ 4.11 คุณภาพการจัดอันดับของ candidate scoring บน full candidate pairs โดยแสดง precision@100 precision@500 precision@1000 และ precision@5000 ร่วมกับค่า Candidate\end{center}
AP และ ROC-AUC ค่าเหล่านี้มีความหมายเชิงระบบมาก เพราะชี้ว่าแม้ candidate pool จริงจะมี class i m b a l a n c e สูงและมี c a n d i d a t e c o u n t ระดับล้านคู่ แต่คะแนนของแบบจำลองยังสามารถดันคู่จริงจำนวนมากขึ้นไปอยู่ในอันดับต้น ๆ ได้อย่างมีประสิทธิภาพ ซึ่งเป็นคุณสมบัติที่จำเป็นมากสำหรับ thresholding และ review workflow\par
% ---- source PDF page 85 ----
\clearpage
\section*{4.4 Exact-First ร่วมกับ Calibrated Candidate Scoring สำหรับ Decision}
Thresholding\par
เมื่อได้ candidate-scoring model ที่ผ่านการ calibration แล้ว ขั้นถัดไปคือการนำ คะแนนดังกล่าวไปเชื่อมกับ exact-first เพื่อสร้างกฎตัดสินใจที่ใช้ได้จริงใน CRM จุดตั้งต้น ของ ปัญหานี้คือ valid profiles 36,804 โปรไฟล์ก่อให้เกิดคู่แบบ cross-platform ทั้งหมด 449,149,239 คู่ ซึ่งไม่สามารถเปรียบเทียบทุกคู่โดยตรงได้ในทางปฏิบัติ ดังนั้น pipeline จึงต้อง ลด search space ก่อน\par
ผลรันจริงแสดงว่า exact-first และ blocking สามารถคัดคู่ exact ที่ merge ได้ทันที 12,403 คู่ และลดคู่ที่ต้องให้แบบจำลองประเมินต่อเหลือ 2,073,842 คู่ คิดเป็นการลดพื้นที่ค้นหา ลง 99.54\% พร้อมรักษา coverage ของ ground-truth matches ไว้ 88.67\% ตารางที่ 4.8 Retrieval funnel ของ Project 1 ก่อนเข้าสู่การตัดสินใจปลายทาง\par
ขั้นตอน จำนวนคู่ All cross-platform pairs 449,149,239 Exact-first auto-accepted pairs 12,403 Candidate pairs for model 2,073,842 Final MATCH only 20,549 REVIEW queue 86,296\par
\begin{center}\small รูปที่ 4.12 Retrieval funnel ของ Project 1 แสดงการลด search space จาก all cross-platform\end{center}
pairs ไปสู่ exact-first pairs model candidates final MATCH only และ REVIEW queue\par
% ---- source PDF page 86 ----
\clearpage
\begin{center}\small รูปที่ 4.13 โครงสร้างของ Project 1 ในขั้น decision thresholding เริ่มจาก exact-first และ blocking ก่อนส่ง candidate pairs เข้าสู่ calibrated candidate scoring แล้วแยกผลลัพธ์เป็น\end{center}
MATCH REVIEW และ NO\_MATCH\par
\begin{center}\small ตารางที่ 4.8 และรูปที่ 4.11-4.12 ชี้ให้เห็นว่า `exact-first` ไม่ได้เป็นเพียง heuristic เสริม แต่เป็นองค์ประกอบหลักของ Project 1 เพราะช่วยดึงคู่ที่ชัดมากออกจากระบบตั้งแต่ต้น และเปิดทางให้โมเดลโฟกัสเฉพาะคู่ที่กำกวมกว่า ซึ่งเป็นส่วนที่ calibration และ thresholding จะสร้างคุณค่าจริง\end{center}
\section*{4.5 ผลลัพธ์ปลายทางของ MATCH, REVIEW และ NO\_MATCH สำหรับ CRM}
ในระดับ production operating point ที่ถูกนำไปใช้จริง ระบบกำหนด match\_threshold = 0.98 และ review\_threshold = 0.95 เพื่อควบคุมความเสี่ยงของการ merge ผิดในบริบท CRM ซึ่ง false merge ถือเป็นความเสียหายที่รุนแรงกว่าการส่งคู่ไปให้มนุษย์ตรวจซ้ำ กติกาการตัดสินใจจริงจึงเป็นดังนี้\par
\noindent 1. ถ้า username หรือ URL ตรงกันแบบ exact ให้ตัดสินเป็น MATCH ทันที 2. ถ้าไม่ใช่ exact pair แต่ calibrated score >= 0.98 ให้ตัดสินเป็น MATCH 3. ถ้าคะแนนอยู่ในช่วง 0.95 <= score < 0.98 ให้ส่งเข้า REVIEW 4. ถ้าคะแนนต่ำกว่า 0.95 ให้ตัดสินเป็น NO\_MATCH ตารางที่ 4.9 Operating point ที่ใช้จริงในขั้น decision thresholding\par
ชั้นการตัดสินใจ เงื่อนไข จำนวนคู่ บทบาทใน workflow\par
Exact-First Username หรือ URL\par
ตรงกันแบบ exact 12,403 ส่งเข้า `MATCH` ทันที\par
% ---- source PDF page 87 ----
\clearpage
High-score Match Calibrated score\par
`>= 0.98` 8,146 ส่งเข้า `MATCH`\par
แบบอัตโนมัติ\par
Review Band `0.95 <= score <\par
\noindent 0.98` 86,296 ส่งเข้า `REVIEW`\par
ให้ผู้เชี่ยวชาญ Low-score Reject `score < 0.95` 1,979,400 ตัดเป็น `NO\_MATCH`\par
เมื่อใช้ threshold ปลายทาง 0.98 / 0.95 กับคะแนนที่ผ่านการ calibration แล้ว ระบบได้ผลลัพธ์ปลายทางเป็น MATCH = 20,549 คู่, REVIEW = 86,296 คู่ และ NO\_MATCH = 1,979,400 คู่ จาก final decisions ทั้งหมด 2,086,245 คู่ กล่าวอีกนัยหนึ่งคือ ระบบสามารถจัดการคู่ส่วนใหญ่ได้อัตโนมัติ และกันคู่ที่ยังไม่แน่ชัดออกมาให้มนุษย์ตรวจเพียง ประมาณ 4.14\% ของ final decisions เท่านั้น\par
คุณภาพของ MATCH ที่ระบบยอม merge อัตโนมัติเป็นจุดที่สำคัญที่สุดสำหรับ CRM หลังรวม exact matches กับ high-score matches เข้าด้วยกัน ระบบได้ final match-only precision = 0.9550 และ final match-only recall = 0.6711 ความหมายเชิงปฏิบัติของตัวเลขนี้คือ ในทุก 100 คู่ที่ระบบ auto-merge จะมี false merge ประมาณ 4-5 คู่ ขณะเดียวกันระบบกู้คู่จริงได้เองประมาณ 67.11\% ของคู่จริงทั้งหมด ส่วนที่ยังไม่มั่นใจพอจะไม่ถูก merge ทันที แต่ถูกส่งไปยัง review queue แทน ตารางที่ 4.10 คุณภาพของแต่ละ tier ที่เกี่ยวข้องกับการตัดสินใจปลายทาง\par
กลุ่มตัดสินใจ จำนวนคู่ คู่จริงในกลุ่ม Precision Recall ต่อคู่จริงทั้งหมด ความหมายเชิงปฏิบัติ\par
Exact auto 12,403 12,340 0.9949 0.4220 merge ได้ทันทีด้วยกฎ exact Score `>= 0.98` 8,146 7,284 0.8942 0.2491 auto-merge จาก calibrated score REVIEW 86,296 4,065 0.0471 0.1390 ส่งให้ผู้เชี่ยวชาญตรวจ Final MATCH only\par
20,549 19,624 0.9550 0.6711 คุณภาพรวมของ\par
auto-merge\par
% ---- source PDF page 88 ----
\clearpage
\begin{center}\small รูปที่ 4.14 ผลของ thresholding ต่อปริมาณงานและคุณภาพการตัดสินใจ แสดงจำนวนคู่ในแต่ละ tier\end{center}
precision ของกลุ่ม auto-merge และ recall contribution ของ exact auto score band สูง\par
และกลุ่มที่ถูกส่งเข้า review ผลในตารางที่ 4.10 และรูปที่ 4.13 สะท้อนบทบาทของแต่ละชั้นได้ชัดมาก โดย exact- first ช่วยกู้ recall ได้ 42.20\% ตั้งแต่ก่อนใช้แบบจำลอง และทำด้วย precision เกือบสมบูรณ์ (0.9949) จากนั้นคะแนน calibrated ช่วงสูง (>= 0.98) ช่วยเพิ่ม recall อัตโนมัติอีก 24.91\% ของคู่จริงทั้งหมด แม้ precision ของชั้นนี้จะต่ำกว่า exact rules แต่เมื่อนำมารวมกันใน Final MATCH only ระบบยังรักษา precision รวมไว้ได้ 0.9550 ในอีกด้านหนึ่ง REVIEW tier มี precision เพียง 0.0471 ซึ่งไม่ควรถูกตีความว่าเป็น จุดอ่อนของระบบ แต่ควรถูกตีความว่าเป็น buffer zone ที่ถูกออกแบบไว้โดยเจตนา กล่าวคือ ระบบกันคู่ที่ยังไม่ควร auto-merge ออกไปให้มนุษย์ตัดสินแทน เพื่อรักษาคุณภาพ ข้อมูลลูกค้าในระดับองค์กร\par
\subsection*{4.5.1 Production Decision Matrix ที่ Operating Point สุดท้าย}
หากพิจารณาผลลัพธ์ปลายทางในรูปของ d e c i s i o n m a t r i x จะเห็นโครงสร้างการกระจายของคู่จริงและคู่ไม่จริงในแต่ละ tier ชัดเจนยิ่งขึ้น กล่าวคือจากคู่จริงทั้งหมด 29,243 คู่ ระบบจัดให้ 19,624 คู่เข้าสู่ MATCH, 4,065 คู่เข้าสู่ REVIEW และยังเหลือคู่จริง 5,554 คู่ที่ถูกตัดเป็น NO\_MATCH ขณะเดียวกันในฝั่งคู่ที่ไม่จริง ระบบยังคงกันไว้ใน NO\_MATCH ได้ 1,973,846 คู่ และมีคู่ไม่จริงที่หลุดเข้า MATCH เพียง 925 คู่\par
% ---- source PDF page 89 ----
\clearpage
\begin{center}\small รูปที่ 4.15 Production decision matrix ที่ operating point สุดท้าย แสดงความสัมพันธ์ระหว่าง\end{center}
Actual MATCH และ Actual NO\_MATCH กับผลการตัดสินใจใน tier MATCH REVIEW และ\par
NO\_MATCH\par
\begin{center}\small รูปที่ 4.14 ช่วยให้ตีความ operating point ได้ชัดขึ้นในเชิงปฏิบัติ กล่าวคือ MATCH tier ถูกออกแบบให้กินคู่จริงเป็นหลัก ขณะที่ REVIEW tier ทำหน้าที่เป็นเขตกันชนสำหรับคู่ที่ยังมีศักยภาพแต่ไม่ควรถูก merge อัตโนมัติ และ NO\_MATCH tier รับภาระในการคัดคู่ไม่จริงส่วนใหญ่ของระบบไว้ ผลในลักษณะนี้สอดคล้องกับเป้าหมายของ CRM ซึ่งให้ความสำคัญกับการลด false merge มากกว่าการพยายาม merge ให้ได้ครอบคลุมที่สุดตั้งแต่รอบแรก\end{center}
\subsection*{4.5.2 การวิเคราะห์ False Positives และ False Negatives จากตัวอย่างข้อผิดพลาด}
เพื่อไม่ให้การอภิปรายผลหยุดอยู่เพียงที่ m e t r i c รวม ผู้วิจัยจึงวิเคราะห์ไฟล์ตัวอย่างข้อผิดพลาดจริง fp\_top.csv และ fn\_top.csv อย่างละ 200 คู่ ซึ่งเก็บคู่ที่แบบจำลองผิดพลาดเด่นที่สุดจาก test split ของ run หลัก ผลการสรุป pattern แสดงให้เห็นว่ากลุ่ม false negatives และ false positives มีธรรมชาติแตกต่างกันอย่างชัดเจน ตารางที่ 4.11 สรุป pattern ของข้อผิดพลาดจากตัวอย่าง `fp\_top.csv` และ `fn\_top.csv`\par
ตัวชี้วัด False Positives False Negatives ขนาดตัวอย่างที่วิเคราะห์ 200 200\par
คู่แพลตฟอร์มที่พบบ่อยที่สุด `googleplus -> twitter` 50\par
`googleplus -> twitter` 110\par
คู่ (25.0\%)\par
คู่ (55.0\%)\par
% ---- source PDF page 90 ----
\clearpage
ทั้ง `username\_jaro` และ `fullname\_jaro` ต่ำกว่า 0.5 17 คู่ (8.5\%) 82 คู่ (41.0\%)\par
อย่างน้อยหนึ่งชื่อมี similarity ต่ำกว่า 0.5 69 คู่ (34.5\%) 135 คู่ (67.5\%)\par
ทั้ง `username\_jaro` และ `fullname\_jaro` ตั้งแต่ 0.9 ขึ้นไป\par
10 คู่ (5.0\%) 3 คู่ (1.5\%)\par
มีภาพอย่างน้อยหนึ่งฝั่ง 51 คู่ (25.5\%) 71 คู่ (35.5\%) มี caption signal 20 คู่ (10.0\%) 20 คู่ (10.0\%)\par
\begin{center}\small รูปที่ 4.16 การวิเคราะห์ข้อผิดพลาดจากตัวอย่าง false positives และ false negatives\end{center}
แสดงการกระจุกตัวตามคู่แพลตฟอร์มและลักษณะของสัญญาณชื่อ ภาพ และ caption\par
ที่เกี่ยวข้องกับความผิดพลาด\par
ผลในตารางที่ 4.11 และรูปที่ 4.15 ให้ข้อสังเกตสำคัญ 2 ประเด็น ประเด็นแรก false negatives กระจุกตัวอยู่ในคู่ googleplus -> twitter อย่างเด่นชัด และส่วนใหญ่มีสัญญาณจากชื่ออ่อนมาก กล่าวคือในตัวอย่าง 67.5\% มีอย่างน้อยหนึ่งชื่อที่ similarity ต่ำกว่า 0.5 และ 41.0\% มีทั้ง username และ fullname ต่ำกว่า 0.5 พร้อมกัน ซึ่งสะท้อนว่าข้อผิดพลาดประเภทนี้มักเกิดเมื่อผู้ใช้เปลี่ยนชื่อผู้ใช้หรือชื่อแสดงผลข้ามแพลตฟอร์มจ นสัญญาณแบบ name-based แทบไม่เหลือ\par
ประเด็นที่สอง false positives จำนวนมากกลับเกิดในบริบทตรงข้าม คือเป็นคู่ลบที่ดู เหมือนคู่จริงมากเกินไป โดยเฉพาะกลุ่ม hard negatives ที่มีชื่อหรือชื่อผู้ใช้คล้ายกันมาก ตัวอย่างเช่นกรณี mitchellhall / mitchell hall และ pullbackes / pullbackes ซึ่งได้คะแนน สูงมากแม้เป็นคนละ user\_folder จริง รูปแบบนี้อธิบายได้ว่าทำไมระบบจึงยังต้องมี REVIEW tier เพื่อกันคู่ที่ดูน่าใช่มาก แต่ยังไม่ปลอดภัยพอสำหรับการ merge อัตโนมัติในทุกกรณี\par
% ---- source PDF page 91 ----
\clearpage
เมื่อพิจารณาร่วมกับผล feature importance จะเห็นว่าข้อผิดพลาดทั้งสองแบบสอดคล้องกับธรรมชาติของแบบจำลอง กล่าวคือระบบเก่งมากในกรณีที่มีสัญญาณจากชื่อชัด แต่ยังมีข้อจำกัดเมื่อชื่ออ่อนเกินไปหรือเมื่อคู่ลบมีชื่อเหมือนกันผิดปกติ ดังนั้นข้อเสนอเชิงเทคนิคสำหรับการพัฒนาต่อจึงไม่ใช่เพียงการปรับ threshold เท่านั้น แต่ควรเพิ่ม coverage ของหลักฐานเสริมที่ช่วยแยกสองกรณีนี้ได้ เช่น image signal ที่มีทั้งสองฝั่ง, caption-context ที่สมบูรณ์ขึ้น หรือ contextual features อื่นที่ช่วยยืนยันตัวตนเมื่อสัญญาณจากชื่อไม่เพียงพอ\par
\section*{4.6 ความสำคัญของคุณลักษณะทั้ง 41 ตัวที่สนับสนุนการตัดสินใจ}
เพื่อทำความเข้าใจว่าระบบอาศัยหลักฐานประเภทใดในการตัดสินใจ ผู้วิจัยวิเคราะห์ feature importance ของ main run ซึ่งใช้ feature ทั้งหมด 41 ตัว ผลที่ได้ชี้ชัดว่าระบบยังพึ่งพาสัญญาณจากชื่อเป็นหลักอย่างมาก โดยเฉพาะกลุ่ม fullname\_* และ username\_* ขณะที่สัญญาณจาก bio, style และ image-cross-signal ทำหน้าที่เป็นหลักฐานเสริมมากกว่าจะเป็นแกนหลักของการตัดสินใจ\par
\begin{center}\small ตารางที่ 4.12 คุณลักษณะที่สำคัญสูงสุดของ main run\end{center}
อันดับ Feature Importance 1 `fullname\_token\_sort` 0.4165 2 `username\_token\_sort` 0.2263 3 `fullname\_jaro` 0.1561 4 `bio\_tfidf\_cosine` 0.0588 5 `username\_jaro` 0.0464 6 `fullname\_lev` 0.0409 7 `bio\_sbert\_cosine` 0.0183 8 `username\_lev` 0.0089\par
% ---- source PDF page 92 ----
\clearpage
9 `platform\_pair\_code` 0.0076 10 `style\_biolen\_ratio` 0.0040\par
หากรวมความสำคัญตามตระกูลของ feature จะพบว่า Name similarity ครองสัดส่วนรวมถึง 89.51\% ของความสำคัญทั้งหมด ขณะที่ Bio text มี 7.71\%, Style / writing มี 1.21\%, Platform pair มี 0.76\% และ Image-caption cross-signal มี 0.65\% ที่สำคัญคือใน 41 features นี้มีถึง 18 ตัวที่ importance เป็นศูนย์ใน main run\par
\begin{center}\small รูปที่ 4.17 ความสำคัญของคุณลักษณะทั้ง 41 ตัวใน main run แสดงให้เห็นว่า feature กลุ่มชื่อและ\end{center}
username เป็นหลักฐานหลักของระบบ ขณะที่ feature จำนวนหนึ่งมีผลต่อการตัดสินใจต่ำมากหรือเป็นศูนย์ในรอบนี้\par
% ---- source PDF page 93 ----
\clearpage
\begin{center}\small รูปที่ 4.18 การรวม feature importance ตามตระกูลหลัก แสดงสัดส่วนของหลักฐานจากชื่อ bio\end{center}
style image-cross-signal และหลักฐานประเภทอื่นที่แบบจำลองใช้จริงใน main run\par
อย่างไรก็ตาม ผลในระดับ feature importance ต้องตีความร่วมกับข้อจำกัดของข้อมูลภาพด้วย โดย valid profiles ทั้งหมด 36,804 โปรไฟล์ มีเพียง 4,248 โปรไฟล์ที่มี local image, 2,173 โปรไฟล์ที่มี image metadata และ 2,173 โปรไฟล์ที่มี caption ใช้งานได้ อีกทั้งใน test pairs มีคู่ที่มีภาพอย่างน้อยหนึ่งฝั่ง 1,873 คู่ แต่ไม่มีคู่ที่มี local image พร้อมกันทั้งสองฝั่งเลย\par
\begin{center}\small รูปที่ 4.19 Coverage ของข้อมูลภาพใน main run แยกระดับ profile และ test pair แสดงให้เห็นว่าการทดลองในรอบนี้ยังเป็น partial multimodal setting มากกว่า full image-to-\end{center}
image matching\par
ผลในรูปที่ 4.16 อธิบายได้ว่าทำไม image-related features จำนวนมากยังมีน้ำหนักต่ำกว่ากลุ่มชื่อและข้อความอย่างชัดเจน กล่าวคือ image branch ช่วยเพิ่มคุณภาพของ main run จริง แต่ยังอยู่ในฐานะสัญญาณเสริมภายใต้ coverage ที่จำกัด\par
\section*{4.7 สรุปผลของ Project 1 ในขอบเขต Decision Thresholding}
ผลการวิจัยในบทนี้ยืนยันว่าแกนหลักของ Project 1 คือ Decision thresholding ที่ ทำงานต่อจาก Exact-first + calibrated candidate scoring โดยใช้ main run แบบ image\_context เป็นตัวให้คะแนนคู่โปรไฟล์ แล้วแปลงคะแนนนั้นไปเป็นกฎการตัดสินใจใน workflow จริงของ CRM\par
หากสรุปในเชิง implementation ที่พร้อมนำไปใช้ต่อได้โดยตรง สามารถสรุปได้ดังนี้ 1. ใช้ image\_context เป็น candidate scorer หลัก เพราะให้สมดุลดีที่สุดในด้าน AP,\par
ROC-AUC และ Recall บน feature space 41 ตัว\par
% ---- source PDF page 94 ----
\clearpage
\noindent 2. ใช้ IsotonicRegression เพื่อ calibrate ความน่าจะเป็นก่อนตั้ง threshold\par
ปลายทาง เพื่อให้คะแนนมีความหมายเชิงปฏิบัติการมากขึ้น 3. ใช้กฎ exact-first เพื่อดึงคู่ที่ชัดมากออกมา merge อัตโนมัติก่อน แล้วค่อยใช้\par
calibrated score กับคู่ที่ยังมีความกำกวม 4. ใช้ thresholds 0.98 สำหรับ MATCH และ 0.95 สำหรับ REVIEW\par
เพื่อควบคุมความเสี่ยงของ false merge ในบริบท CRM 5. ที่ operating point นี้ ระบบได้ final match-only precision = 0.9550 และ\par
final match-only recall = 0.6711 พร้อม review queue = 86,296 คู่ ซึ่งสะท้อนสมดุลที่ยอมรับได้ระหว่างการ merge อัตโนมัติและการส่งคู่กำกวม ให้มนุษย์ตรวจสอบต่อ ดังนั้น ข้อสรุปสุดท้ายของ Project 1 ไม่ใช่เพียงว่าแบบจำลองให้ metric ดีบน test split แต่คือระบบสามารถแปลงคะแนนของแบบจำลองให้กลายเป็นกระบวนการตัดสินใจแบบ MATCH / REVIEW / NO\_MATCH ที่ใช้งานได้จริงใน CRM โดยรักษาสมดุลระหว่างการ merge อัตโนมัติและการส่งคู่กำกวมให้มนุษย์ตรวจสอบต่ออย่างเป็นระบบ\par
\begin{center}\small รูปที่ 4.20 สรุปเส้นทางการพัฒนาระบบจาก Leakage diagnosis ไปสู่ Final multimodal\end{center}
\chapter*{บทที่ 5}
\begin{center}\Large\bfseries วิเคราะห์และสรุปผล\end{center}
\section*{5.1 สรุปผลการดำเนินงาน}
โครงงานนี้มีวัตถุประสงค์หลักในการพัฒนาระบบ Identity Resolution สำหรับข้อมูล Social Media Profiles จาก 3 แพลตฟอร์ม ได้แก่ Twitter, Instagram และ Google+ Archive รวม 36,807 profiles เพื่อสร้าง Unified Customer Profile ที่ลดความซ้ำซ้อนและเพิ่ม ความ\par
% ---- source PDF page 95 ----
\clearpage
ครบถ้วนของข้อมูล สำหรับนำไปประยุกต์ใช้ใน Lead Scoring ของระบบ Agentic CRM ใน อนาคต ผลการดำเนินงานทั้ง 2 ระยะที่ดำเนินการได้สำเร็จในงานวิจัยนี้สรุปได้ดังนี้\par
ในระยะที่ 1 การเตรียมข้อมูล ดำเนินการ Data Cleaning และ Standardization กับ Profiles ทั้งหมด 36,807 รายการ ครอบคลุมการแก้ไข Encoding, การ Normalize รูปแบบ Username และ Full Name ให้เป็นมาตรฐานเดียวกันข้ามแพลตฟอร์ม รวมถึงการสร้าง Candidate Pairs ด้วย Blocking Strategy ที่ผสม Username Prefix Blocking กับ Full Name Jaro-Winkler Similarity ได้ Candidate Pairs ทั้งสิ้น 873,956 คู่ และ Labeled Dataset จำนวน 162,442 pairs สำหรับใช้ฝึกโมเดล\par
ในระยะที่ 2 การพัฒนาโมเดล Identity Resolution ได้ออกแบบและฝึกโมเดล IdentityMLP ซึ่งเป็น Neural Network แบบ Multi-Layer Perceptron สำหรับ 3-Class Classification ใช้ FocalLoss แก้ปัญหา Class Imbalance และ Isotonic Calibration เพื่อให้ ค่า Probability Output สะท้อนความน่าจะเป็นจริง ผลการทดสอบบน Test Set ได้ F1-Score = 0.93, Precision = 0.97, Recall = 0.90 และ AUC-ROC = 0.982 เมื่อนำโมเดลไปทำ Inference บน Candidate Pairs ทั้ง 873,956 คู่ ได้ผลลัพธ์เป็น MATCH 20,480 คู่ (2.34\%), POSSIBLE\_MATCH 3,940 คู่ (0.45\%) และ NO\_MATCH 849,536 คู่ (97.21\%) ส่งผลให้ได้ Unified Profiles ทั้งสิ้น 10,723 profiles คิดเป็น Deduplication Ratio 70.9\% โดยสรุป ระบบที่พัฒนาขึ้นสามารถบรรลุวัตถุประสงค์ด้าน Identity Resolution ได้สำเร็จ ทั้งในด้านความ แม่นยำของโมเดล (F1 = 0.93) และความสามารถในการ Scale ไปยัง Candidate Pairs กว่า 8 แสนคู่ ส่วนขั้นตอน Lead Scoring ยังต้องพัฒนาต่อเมื่อ ข้อมูล Engagement และ Behavior พร้อมใช้งาน\par
งานวิจัยนี้ได้ศึกษาและพัฒนาระบบเชื่อมโยงตัวตนข้ามแพลตฟอร์ม (Cross-Platform Identity Resolution) สำหรับระบบบริหารจัดการความสัมพันธ์ลูกค้า (CRM) โดยใช้ข้อมูลโปร ไฟล์ผู้ใช้จาก Twitter, Instagram และ Google+ จำนวนรวมทั้งสิ้น 36,804 โปรไฟล์\par
ระบบถูกออกแบบให้ทำงานเป็นลำดับขั้นตอน เริ่มตั้งแต่การลดขอบเขตการค้นหา (Search Space) ด้วยเทคนิคการจับคู่แบบตรงกันทุกประการ (Exact Matching) และการจัด กลุ่มข้อมูล (Blocking) จากนั้นจึงทำการให้คะแนนความคล้ายคลึงของคู่ข้อมูลด้วยแบบจำลองการ เรียนรู้ของเครื่อง (Machine Learning) และตัดสินใจขั้นสุดท้ายสำหรับการนำไปใช้งานจริง (Production) ด้วยการใช้เกณฑ์ตัดสินใจแบบคู่ (Dual Threshold)\par
ในการวิจัยนี้ได้ทำการทดลองและเปรียบเทียบประสิทธิภาพของแบบจำลอง โดยแบ่ง ออกเป็น 3 กลุ่มการทดลองหลัก ได้แก่ 1) กลุ่มคลาสสิกแบบป้องกันข้อมูลรั่วไหล (Classical Leakage-Safe) ซึ่งใช้ชุดคุณลักษณะ (Features) จำนวน 22 ตัว 2) กลุ่มพหุรูปแบบ (Multimodal Suite) ที่มีการขยายชุดคุณลักษณะจาก 23 ตัวเป็น 41 ตัว และ 3) กลุ่มตัวแบบ อ้างอิง (Neural-Network Reference) ที่ใช้แบบจำลอง IdentityMLP เป็นตัวเปรียบเทียบ ทั้งนี้\par
% ---- source PDF page 96 ----
\clearpage
การประเมินผลการทดลองทั้งหมดอ้างอิงจากการแบ่งชุดข้อมูล (Data Split) ที่ออกแบบมาเพื่อ ป้องกันปัญหาข้อมูลตัวตนรั่วไหล (Identity Leakage) อย่างเคร่งครัด โดยควบคุมไม่ให้มีความ ซ้ำซ้อนขององค์ประกอบข้อมูลในแต่ละชุดการทดลอง (Zero Component Overlap) ตารางที่ 5.1 สรุปเปรียบเทียบประสิทธิภาพจากทุกสายการทดลอง\par
สายการทดลอง จำนวนคุณลักษณะ Test AP Test AUC Test F1 Classical Leakage- Safe (GB) 22 0.9703 0.9772 0.9337\par
Multimodal Hybrid Text 23 0.9782 0.9715 0.9330\par
Multimodal Image Stats 37 0.9789 0.9731 0.9340\par
Multimodal Image Context (main run)\par
41 0.9789 0.9734 0.9340\par
Neural-Network MLP (41 features) 41 0.9734 0.9649 0.9255\par
Neural-Network MLP (strict 22 features)\par
22 0.7344 0.9960 0.8296\par
จากผลการศึกษาและการทดลองทั้งหมด สามารถสรุปประเด็นสำคัญของงานวิจัยได้ 4 ประเด็น ดังนี้\par
\noindent 1. ปัญหาหลักของการพัฒนาระบบไม่ได้อยู่ที่การเลือกแบบจำลองแต่อยู่ที่วิธีการสร้างคู่ข้อมูล\par
และการออกแบบการแบ่งชุดข้อมูล เห็นได้จากการประมวลผล ในรอบแรกที่ค่าตัวชี้วัด\par
(Metrics) เช่น AP, AUC และ F1 มีค่าสูงถึง 1.0000 ซึ่งแท้จริงแล้วเกิดจากปัญหาข้อมูล\par
ตัวตนรั่วไหล (Identity Leakage) เมื่อทำการแก้ไขปัญหานี้แล้ว ค่าตัวชี้วัดจึงลดลงมาอยู่\par
ในระดับที่สมเหตุสมผลและสะท้อนประสิทธิภาพที่แท้จริงของแบบจำลองได้\par
% ---- source PDF page 97 ----
\clearpage
\noindent 2. แบบจำลอง Gradient Boosting เป็นทางเลือกที่เหมาะสมที่สุดสำหรับงานวิจัยนี้ โดยใน\par
การทดลองกลุ่มคลาสสิกสามารถทำคะแนน Test AP ได้ 0.9703 และ F1 ได้ 0.9337 ซึ่ง\par
สูงที่สุดในกลุ่มแบบจำลอง 3 ชนิด และเมื่อเปรียบเทียบแบบจำลองทั้ง 7 ชนิดบนชุด\par
คุณลักษณะ 41 ตัว Gradient Boosting ก็ยังคงทำคะแนน Test AP ได้สูงสุดที่ 0.9797\par
ซึ่งถือว่าสูงกว่าแบบจำลองโครงข่ายประสาทเทียม หรือ MLP ที่ทำได้ 0.9734 สาเหตุ\par
หลักเป็นเพราะ Gradient Boosting สามารถวิเคราะห์ความสัมพันธ์ที่ซับซ้อนระหว่าง\par
กลุ่มคุณลักษณะประเภทชื่อผู้ใช้ (Username) และชื่อจริง (Full Name) ได้อย่างมี\par
ประสิทธิภาพ นอกจากนี้ยังสามารถแสดงค่าความสำคัญของคุณลักษณะให้ตีความ\par
ผลลัพธ์ได้ชัดเจนกว่าแบบจำลองโครงข่ายประสาทเทียม\par
\noindent 3. การเพิ่มข้อมูลเข้าสู่ระบบทีละขั้นตอนช่วยเพิ่มความแม่นยำได้อย่างชัดเจน ตัวอย่างเช่น\par
การนำแบบจำลอง SBERT มาช่วยวิเคราะห์ความคล้ายคลึงเชิงความหมาย (Semantic\par
Similarity) ของประวัติย่อ (Bio) ทำให้ระบบสามารถจับคู่ข้อมูลที่ไม่ได้ใช้คำศัพท์\par
(Token) ตรงกันได้ ส่งผลให้ค่าความครอบคลุม (Recall) ดีขึ้น นอกจากนี้ การเพิ่ม\par
คุณลักษณะประเภทรูปภาพยังช่วยปรับปรุงค่า AP และ AUC อย่างต่อเนื่อง แม้จะเป็น\par
การพัฒนาเพียงเล็กน้อยแต่ก็แสดงให้เห็นถึงแนวโน้มที่เสถียร เนื่องจากทุกการทดลองใช้\par
ชุดคู่ข้อมูลและการแบ่งชุดข้อมูลแบบเดียวกัน\par
\noindent 4. ระบบนี้มีความพร้อมในการนำไปใช้งานจริง โดยอาศัยการกำหนดเกณฑ์แบบ Exact-First\par
ร่วมกับเกณฑ์ตัดสินใจแบบคู่ ซึ่งกำหนดเกณฑ์การจับคู่ (Match Threshold) ไว้ที่ 0.98\par
และเกณฑ์การตรวจสอบ (Review Threshold) ที่ 0.95 วิธีการนี้ช่วยลดขอบเขตการ\par
ค้นหาจาก 449 ล้านคู่ เหลือเพียง 2 ล้านคู่ (ลดลงร้อยละ 99.54) แต่ยังคงครอบคลุมคู่\par
ข้อมูลที่ถูกต้อง (True Matches) ได้ถึงร้อยละ 88.67 ผลลัพธ์สุดท้ายมีความแม่นยำใน\par
การจับคู่ (Match Precision) ที่ 0.9550 และสามารถรวบรวมข้อมูลเป็นโปรไฟล์ลูกค้า\par
% ---- source PDF page 98 ----
\clearpage
แบบรวมศูนย์ (Unified Profiles) ได้จำนวน 19,799 รายชื่อ พร้อมสำหรับจัดแบ่งกลุ่ม\par
ลูกค้าเพื่อนำไปใช้วางแผนกลยุทธ์ทางการตลาดได้ทันที\par
\begin{center}\small ตารางที่ 5.2 สรุปผลระดับ Production\end{center}
รายการ ค่า Search space ทั้งหมด 449,149,239 คู่ Candidate pairs หลัง blocking 2,073,842 คู่ Search-space reduction 99.54\% Ground-truth coverage 88.67\% Match precision (exact + model) 0.9550 Unified profiles 19,799 Review queue 86,296 คู่\par
โดยสรุป แบบจำลองที่เหมาะสมที่สุดสำหรับงานวิจัยนี้คือ Gradient Boosting ซึ่ง ดำเนินการบนกระบวนการสร้างชุดคู่ข้อมูลและการแบ่งชุดข้อมูลที่ปราศจากปัญหาข้อมูลตัวตน รั่วไหล พร้อมทั้งผนวกข้อมูลเชิงความหมายของข้อความผ่าน SBERT และข้อมูลบริบทของรูปภาพ เข้าไปเป็นลำดับขั้น โดยรอบการประมวลผลที่ให้ผลลัพธ์ดีที่สุดคือชุดการทดลอง image\_context ซึ่งใช้ตัวแปรคุณลักษณะรวมทั้งสิ้น 41 ตัว และมีจุดตัดสำหรับการนำไปใช้งานจริงที่ระดับเกณฑ์ การตัดสินใจ 0.98 และ 0.95\par
\subsection*{5.1.1 การตอบวัตถุประสงค์ของงานวิจัย}
จากผลการทดลองและการวิเคราะห์ข้อมูลทั้งหมด สามารถสรุปผลการดำเนินงานตาม วัตถุประสงค์ของงานวิจัยได้ดังนี้\par
วัตถุประสงค์ที่ 1 เพื่อศึกษาและพัฒนาระบบเชื่อมโยงตัวตนข้ามแพลตฟอร์มสำหรับ ระบบบริหารจัดการความสัมพันธ์ลูกค้า (CRM) โดยใช้แบบจำลองการเรียนรู้ของเครื่อง\par
ผลการวิจัยยืนยันว่าระบบที่พัฒนาขึ้นสามารถทำงานได้จริงอย่างมีประสิทธิภาพ โดยผล การประมวลผลหลักให้ค่า Test AP ถึง 0.9789 และ F1-Score ที่ 0.9340 ซึ่งอยู่ในระดับที่พร้อม ใช้งานจริง นอกจากนี้ ระบบยังสามารถลดขอบเขตการค้นหา 449 ล้านคู่ ลงเหลือเพียง 2 ล้านคู่ ในขณะที่ยังคงความครอบคลุมคู่ข้อมูลที่ถูกต้องได้ถึงร้อยละ 88.67 แสดงให้เห็นว่าการออกแบบ กระบวนการดึงข้อมูล ร่วมกับแบบจำลองการเรียนรู้ของเครื่อง สามารถรองรับการจัดการข้อมูล ขนาดใหญ่ในเชิงปฏิบัติได้อย่างมีประสิทธิภาพ\par
วัตถุประสงค์ที่ 2 เพื่อเปรียบเทียบประสิทธิภาพของแบบจำลองหลายประเภท\par
% ---- source PDF page 99 ----
\clearpage
ผลการเปรียบเทียบแสดงให้เห็นอย่างชัดเจนว่า แบบจำลอง Gradient Boosting ให้ ผลลัพธ์โดยรวมดีที่สุด ทั้งในด้านค่า AP, F1-Score และความแม่นยำ (Precision) ซึ่งมี ประสิทธิภาพเหนือกว่าแบบจำลอง MLP, Random Forest, Extra Trees, AdaBoost, Logistic Regression และ Linear SVM สาเหตุหลักเป็นเพราะแบบจำลองที่ใช้โครงสร้างต้นไม้เป็นฐาน (Tree-based Models) สามารถวิเคราะห์ปฏิสัมพันธ์ที่ซับซ้อนระหว่างคุณลักษณะต่าง ๆ ได้เป็น อย่างดี โดยเฉพาะกลุ่มคุณลักษณะประเภทชื่อผู้ใช้ ในขณะเดียวกัน แบบจำลองโครงข่ายประสาท เทียมแม้จะสามารถปรับตัวได้ดีขึ้นเมื่อปริภูมิคุณลักษณะ (Feature Space) ขยายตัว แต่ก็ยังคงมี ประสิทธิภาพต่ำกว่า Gradient Boosting ในทุกมิติการวัดผลหลัก อีกทั้งยังมีข้อจำกัดในด้านการ ตีความผลลัพธ์\par
วัตถุประสงค์ที่ 3 เพื่อประเมินความเป็นไปได้ในการนำระบบไปใช้งานจริงในระดับการ ทำงานจริง (Production)\par
ระบบได้ผ่านการทดสอบด้วยเกณฑ์การตัดสินใจสำหรับการใช้งานจริง (Production Thresholds) ในรูปแบบ Exact-First ร่วมกับ Dual Threshold ซึ่งให้ค่าความแม่นยำในการจับคู่ (Match Precision) สูงถึง 0.9550 และมีข้อมูลที่ถูกส่งเข้าคิวรอการตรวจสอบด้วยเจ้าหน้าที่ (Review Queue) จำนวน 86,296 คู่ ท้ายที่สุด ระบบสามารถรวบรวมโปรไฟล์ข้ามแพลตฟอร์มให้ กลายเป็นโปรไฟล์ลูกค้าแบบรวมศูนย์ (Unified Profiles) ได้จำนวน 19,799 รายการ พร้อมทั้งจัด กลุ่มเป้าหมายทางการตลาด (Lead Tiers) ออกเป็นกลุ่ม HOT 4,937 รายการ, WARM 9,039 รายการ และ COLD 5,823 รายการ ซึ่งพร้อมสำหรับการนำไปใช้บริหารลูกค้าสัมพันธ์และ วางแผนการตลาดได้ทันที\par
\section*{5.2 ปัญหาและอุปสรรคในการดำเนินงาน}
ระหว่างการวิจัยและพัฒนาระบบ ผู้วิจัยพบปัญหาและอุปสรรคสำคัญที่ต้องดำเนินการ แก้ไข ดังนี้\par
\subsection*{5.2.1 ปัญหาความไม่สมดุลของข้อมูล (Class Imbalance) และการปรับเทียบแบบจำลอง}
(Model Calibration)\par
ในระยะแรกของการฝึกสอนแบบจำลอง พบว่าสัดส่วนของข้อมูลที่จับคู่กัน (Match) และ ไม่จับคู่ (No Match) มีความแตกต่างกันสูงถึง 1:6 ส่งผลให้แบบจำลองเกิดความลำเอียงและมักจะ ทำนายว่า ไม่จับคู่ ไว้ก่อนเพื่อลดค่าความผิดพลาด ผู้วิจัยจึงแก้ไขด้วยการนำฟังก์ชันการสูญเสีย แบบโฟกัส (Focal Loss) มาใช้ ซึ่งช่วยบังคับให้แบบจำลองให้ความสำคัญกับกรณีที่ยากและ ซับซ้อนมากขึ้น นอกจากนี้ ยังพบปัญหาค่าความน่าจะเป็นที่ส่งออกมา ไม่สะท้อนความเป็นจริง จึง\par
% ---- source PDF page 100 ----
\clearpage
ได้เพิ่มกระบวนการปรับเทียบแบบไอโซโทนิก (Isotonic Calibration) ในขั้นตอนหลังการฝึกสอน ซึ่งช่วยแก้ปัญหาและทำให้การตั้งค่าเกณฑ์ตัดสินใจ (Threshold) มีความแม่นยำขึ้นอย่างชัดเจน\par
\subsection*{5.2.2 อุปสรรคในการจัดการข้อมูลที่ซับซ้อน (Hard Cases)}
แม้ระบบจะแสดงประสิทธิภาพที่น่าพึงพอใจ แต่ยังพบข้อจำกัดในกรณีที่ผู้ใช้งานตั้งชื่อ แตกต่างกันอย่างสิ้นเชิงในแต่ละแพลตฟอร์ม เนื่องจากระบบในปัจจุบันยังคงพึ่งพาคุณลักษณะ ด้านความคล้ายคลึงของตัวอักษรเป็นหลัก ทำให้การจับคู่ในกรณีข้อมูลที่ซับซ้อน (Hard Cases) ทำได้ค่อนข้างยาก\par
\section*{5.3 ข้อจำกัดของงานวิจัย}
งานวิจัยนี้มีข้อจำกัดที่ควรพิจารณาประกอบการตีความผลลัพธ์ ดังต่อไปนี้ 1. ข้อมูลที่ใช้พัฒนามาจากเพียง 3 แพลตฟอร์ม ได้แก่ Twitter, Instagram และ Google+\par
ซึ่งปัจจุบันปิดให้บริการแล้ว และยังเป็นชุดข้อมูลจำลอง จึงอาจไม่สะท้อนความซับซ้อน\par
ของพฤติกรรมลูกค้าในสภาพแวดล้อมการทำงานจริงได้ครบถ้วน อีกทั้งยังไม่สามารถ\par
ประเมินได้ว่าระบบจะทำงานได้ดีเพียงใดหากนำไปใช้กับแพลตฟอร์มยอดนิยมอื่น ๆ เช่น\par
TikTok, LinkedIn หรือ Facebook\par
\noindent 2. ระบบยังมีข้อจำกัดด้านการใช้ข้อมูลรูปภาพ เนื่องจากมีเพียงร้อยละ 11.5 ของโปรไฟล์ที่มี\par
รูปภาพ และร้อยละ 5.9 ที่มีคำบรรยายภาพ (Caption) ยิ่งไปกว่านั้น ไม่มีคู่ข้อมูลทดสอบ\par
(Test Pairs) ใดที่มีรูปภาพพร้อมกันทั้งสองฝั่ง ทำให้ไม่สามารถวัดศักยภาพของการจับคู่\par
ภาพข้ามแพลตฟอร์มโดยตรง (Full Image-to-Image Matching) ได้ในการวิจัยครั้งนี้\par
\noindent 3. งานวิจัยนี้ยังไม่ได้แบ่งข้อมูลเพื่อทดสอบตามลำดับเวลา ทำให้ไม่สามารถจำลอง\par
สถานการณ์ที่ผู้ใช้เปลี่ยนแปลงชื่อ ประวัติย่อ (Bio) หรือรูปโปรไฟล์ไปตามกาลเวลาได้ ซึ่ง\par
อาจส่งผลต่อความแม่นยำของการใช้งานระบบในระยะยาว\par
\noindent 4. แบบจำลองโครงข่ายประสาทเทียม (MLP) มีลักษณะเป็นกล่องดำ ในขณะที่ Gradient\par
Boosting แม้จะสามารถดูค่าความสำคัญของคุณลักษณะในภาพรวมได้ แต่ยังไม่สามารถ\par
อธิบายเหตุผลการตัดสินใจในระดับรายคู่ได้ ซึ่งเป็นอุปสรรคต่อการนำไปใช้งานในธุรกิจที่\par
ต้องการความโปร่งใสสูง\par
\noindent 5. จำนวนคู่ข้อมูลที่ถูกส่งเข้าคิวรอการตรวจสอบมีสูงถึง 86,296 คู่ ซึ่งยังคงเป็นภาระที่ต้อง\par
ใช้ทรัพยากรบุคคลในการดำเนินการ และอาจไม่เหมาะสมสำหรับองค์กรที่มีขนาดเล็ก\par
หรือมีทรัพยากรจำกัด\par
% ---- source PDF page 101 ----
\clearpage
\section*{5.4 ข้อเสนอแนะสำหรับการทำวิจัยในอนาคต}
จากผลการวิจัยและข้อจำกัดที่กล่าวมา ผู้วิจัยมีข้อเสนอแนะสำหรับการต่อยอดในอนาคต ดังนี้\par
\noindent 1. ควรนำระบบไปทดสอบกับข้อมูลจริงจากสภาพแวดล้อมการทำงานจริง รวมถึงขยายไปยัง\par
แพลตฟอร์มที่ใช้งานในปัจจุบัน เช่น TikTok, Facebook หรือ LinkedIn เพื่อทดสอบ\par
ความสามารถในการนำไปใช้งานทั่วไป และควรเพิ่มชุดข้อมูลภาษาไทยเข้าไปในระบบ\par
\noindent 2. เพิ่มคุณลักษณะด้านพฤติกรรม เช่น ช่วงเวลาการใช้งาน หรือหัวข้อเนื้อหา และนำเทคนิค\par
การเรียนรู้เชิงลึก เช่น Contrastive Learning หรือ Siamese Network มาสร้าง\par
เวกเตอร์ตัวแทนรูปภาพ เพื่อยกระดับความแม่นยำในการจับคู่ภาพข้ามแพลตฟอร์ม\par
โดยตรง\par
\noindent 3. การนำเทคนิคปัญญาประดิษฐ์ที่อธิบายได้ มาใช้ควรประยุกต์ใช้เทคนิค เช่น SHAP หรือ\par
LIME เพื่อให้ระบบสามารถอธิบายเหตุผลการจับคู่ข้อมูลเป็นรายคู่ได้ ซึ่งจะช่วยเพิ่มความ\par
น่าเชื่อถือให้แก่ผู้ใช้งานในระบบ CRM\par
\noindent 4. ควรเพิ่มการทดสอบการจำลองสถานการณ์ที่ข้อมูลผู้ใช้เปลี่ยนแปลงไปตามเวลา เพื่อ\par
ประเมินอายุการใช้งาน ของแบบจำลองก่อนที่จะต้องทำการฝึกสอนใหม่\par
\noindent 5. ควรเปิดใช้งานระบบให้เจ้าหน้าที่ตรวจสอบ (Human-in-the-loop) อย่างเต็มรูปแบบ\par
และนำผลการตัดสินใจเหล่านั้นกลับมาเป็นข้อมูลฝึกสอน (Training Set) เพื่อปรับปรุง\par
แบบจำลองอย่างต่อเนื่อง ซึ่งจะช่วยลดปริมาณงานใน Review Queue ลงในระยะยาว\par
\noindent 6. การพัฒนาแบบจำลองลูกผสม (Hybrid Model) และระบบให้คะแนน (Lead Scoring)\par
ควรศึกษาการใช้โครงข่ายประสาทเทียมเพื่อสกัดข้อมูลดิบ แล้วส่งต่อไปยัง Gradient\par
Boosting ในฐานะ Meta-learner นอกจากนี้ ควรต่อยอดการออกแบบโมเดลการให้\par
คะแนนลูกค้า (Lead Scoring Model) ที่พิจารณาเส้นทางลูกค้า (Customer Journey)\par
% ---- source PDF page 102 ----
\clearpage
และรูปแบบการมีส่วนร่วมข้ามแพลตฟอร์ม เพื่อยกระดับสู่ระบบบริหารจัดการลูกค้า\par
สัมพันธ์แบบอัตโนมัติ (Agentic CRM) ต่อไป\par
% ---- source PDF page 103 ----
\clearpage
\section*{บรรณานุกรม}
[1] V. Sharma and C. Dyreson, "LINKSOCIAL: Linking user profiles across multiple\par
social media platforms," in Proc. IEEE Int. Conf. Big Knowledge (ICBK),\par
Singapore, 2018, pp. 260–269, doi: 10.1109/ICBK.2018.00042.\par
[2] R. Zafarani and H. Liu, "Connecting users across social media sites: A behavioral-\par
modeling approach," in Proc. 19th ACM SIGKDD, 2013, doi:\par
\noindent 10.1145/2487575.2487648.\par
[3] S. Halimi and J. Ayday, "User identity linkage on social networks: A review of\par
modern techniques and applications," IEEE Access, vol. 12, 2024.\par
[4] A. H. Doan et al., "Human-in-the-loop challenges for entity matching: A midterm\par
report," in Proc. HILDA'17, 2017, doi: 10.1145/3077257.3077268.\par
[5] G. Li et al., "Human-in-the-loop data integration," Proc. VLDB Endow., 2017, doi:\par
\noindent 10.14778/3137765.3137833.\par
[6] A. Schafer, M. K. Hassan, and J. Muller, "Predictive lead scoring using machine\par
learning techniques," Arabian J. Sci. Eng., vol. 47, no. 9, 2022, doi:\par
\noindent 10.1007/s13369-022-06560-8.\par
[7] V. Sharma and C. Dyreson. “Dataset-LinkSocial.” [Online]. เข้าถึงได้จาก:\par
https://github.com/vishalshar/Dataset-LinkSocial. 2569\par
[8] W. Xu and B. C. M. Fung. "StyleLink: User Identity Linkage Across Social Media\par
with Stylometric Representations." Proceedings of the International AAAI\par
Conference on Web and Social Media (ICWSM), vol. 19, June 2025. pp. 2076-\par
\noindent 2088.\par
% ---- source PDF page 104 ----
\clearpage
% BEGIN GA REFERENCES
\section*{บรรณานุกรม (เพิ่มเติม)}
\noindent [9] G. Papadakis, D. Skoutas, E. Thanos, and T. Palpanas, ``A Survey of Blocking and Filtering Techniques for Entity Resolution,'' arXiv:1905.06167, 2019, revised 2020. ฉบับผู้เขียน: \url{https://arxiv.org/abs/1905.06167}. ฉบับวารสารใช้ชื่อ ``Blocking and Filtering Techniques for Entity Resolution: A Survey,'' \textit{ACM Computing Surveys}, vol. 53, no. 2, article 31, doi: 10.1145/3377455.

\noindent [10] T. Ouyang, W. Pedrycz, and N. J. Pizzi, ``Record linkage based on a three-way decision with the use of granular descriptors,'' \textit{Expert Systems with Applications}, vol. 122, pp. 16--26, 2019, doi: 10.1016/j.eswa.2018.12.038. \url{https://www.sciencedirect.com/science/article/abs/pii/S0957417418308066}.

\noindent [11] C. Sun, D. Shen, Y. Kou, T. Nie, and G. Yu, ``A genetic algorithm based entity resolution approach with active learning,'' \textit{Frontiers of Computer Science}, vol. 11, no. 1, pp. 147--159, 2017, doi: 10.1007/s11704-015-5276-6. \url{https://academic.hep.com.cn/fcs/EN/10.1007/s11704-015-5276-6}.

\noindent [12] C. Elkan, ``The Foundations of Cost-Sensitive Learning,'' in \textit{Proceedings of the Seventeenth International Joint Conference on Artificial Intelligence (IJCAI)}, pp. 973--978, 2001. ฉบับผู้เขียน: \url{https://cseweb.ucsd.edu/~elkan/rescale.pdf}.

\noindent [13] C. Guo, G. Pleiss, Y. Sun, and K. Q. Weinberger, ``On Calibration of Modern Neural Networks,'' in \textit{Proceedings of the 34th International Conference on Machine Learning}, PMLR, vol. 70, pp. 1321--1330, 2017. \url{https://proceedings.mlr.press/v70/guo17a.html}.

\noindent [14] N. Reimers and I. Gurevych, ``Sentence-BERT: Sentence Embeddings using Siamese BERT-Networks,'' in \textit{Proceedings of EMNLP-IJCNLP}, pp. 3982--3992, 2019, doi: 10.18653/v1/D19-1410. \url{https://aclanthology.org/D19-1410/}.

\noindent [15] scikit-learn developers, ``Probability calibration,'' \textit{scikit-learn documentation}. \url{https://scikit-learn.org/stable/modules/calibration.html}. เข้าถึงวันที่ 20 กันยายน 2569.

\noindent [16] scikit-learn developers, ``Common pitfalls and recommended practices,'' \textit{scikit-learn documentation}. \url{https://scikit-learn.org/stable/common_pitfalls.html}. เข้าถึงวันที่ 20 กันยายน 2569.

\noindent [17] scikit-learn developers, ``Metrics and scoring: quantifying the quality of predictions,'' \textit{scikit-learn documentation}. \url{https://scikit-learn.org/stable/modules/model_evaluation.html}. เข้าถึงวันที่ 20 กันยายน 2569.

\noindent [18] M. Mitchell, \textit{An Introduction to Genetic Algorithms}. Cambridge, MA: MIT Press, 1996. ข้อมูลสำนักพิมพ์: \url{https://mitpress.mit.edu/9780262631853/an-introduction-to-genetic-algorithms/}.

% END GA REFERENCES
\clearpage
\section*{ภาคผนวก}
% ---- source PDF page 105 ----
\clearpage
\section*{ประวัติผู้เขียน}
ชื่อ-นามสกุล นางสาววิมลสิริ ปานมา\par
รหัสนักศึกษา 66070187\par
วัน เดือน ปีเกิด 11 สิงหาคม พ.ศ. 2547\par
ประวัติการศึกษา วุฒิมัธยมศึกษาตอนปลายโรงเรียนสระบุรีวิทยาคม\par
จังหวัดสระบุรี\par
ภูมิลำเนา 77 หมู่ 1 ตำบลท่าตูม อำเภอแก่งคอย จังหวัดสระบุรี\par
Email 66070187@kmitl.ac.th\par
สาขาที่จบ วิทยาการข้อมูลและการวิเคราะห์เชิงธุรกิจ\par
รุ่นที่ 7 ปีการศึกษาที่จบ 2570\par
ชื่อ-นามสกุล นางสาวถาวรีย์ ขันสาคร\par
รหัสนักศึกษา 66070260\par
วัน เดือน ปีเกิด 2 กรกฎาคม พ.ศ. 2546\par
ประวัติการศึกษา วุฒิมัธยมศึกษาตอนปลายโรงเรียนสายปัญญารังสิต\par
ภูมิลำเนา 59/9 หมู่ 3 ตำบลบึงบอน อำเภอหนองเสือ จังหวัด\par
ปทุมธานี 12170\par
Email 66070260@kmitl.ac.th\par
สาขาที่จบ วิทยาการข้อมูลและการวิเคราะห์เชิงธุรกิจ\par
รุ่นที่ 7 ปีการศึกษาที่จบ 2570\par
\end{document}
